\chapter{Quantum Toroidal Algebras and the K3 Landscape}
\label{ch:k3-quantum-toroidal}

The K3 Yangian $Y(\mathfrak{g}_{K3})$ is the rational limit of a
two-parameter family: the K3 quantum toroidal algebra.  This chapter
develops the quantum toroidal deformation hierarchy, the AGT
correspondence for K3, the twisted M-theory compactification web, and
the tropical cluster algebra structure.

All algebras of this chapter sit at the stage-$2$ shadow level of the
two-stage factorisation $\Phi_2 = \SpCh_{\Sigma_1, C} \circ \PhiFA_2$
(Theorem~\ref{thm:phi-two-stage-factorisation};
Definition~\ref{def:phi-fa-and-sp}). The canonical stage-$1$ datum is
the $E_2$-holomorphic factorisation algebra $\PhiFA_2(D^b\Coh(K3))$ on
$K3$; the Mukai signature-$(4, 20)$ Heisenberg $H_{\mathrm{Muk}} =
\SpCh_{\Sigma_1^{\mathrm{Nakajima}}, C}(\PhiFA_2(D^b\Coh(K3)))$ is its
Nakajima-cycle specialisation. The Yangian, quantum affine, and quantum
toroidal enhancements are further stage-$2$ specialisations of the same
canonical $\PhiFA_2$, parameterised by the rational $\to$ trigonometric
$\to$ elliptic structure-function deformation and controlled by
Proposition~\ref{prop:native-en-level}: native operadic level $\Etwo$
on $K3$, specialising to $\Eone$-chiral on the reference curve $C$. The
braided $\Etwo$-enhancement manifest in the $R$-matrix lives on the
derived chiral centre of $\Rep^{\Eone}$, not on the algebra itself
(Remark~\ref{rem:e2-on-drinfeld-centre}).

%% ========================================================================
%% The K3 quantum toroidal algebra
%% ========================================================================

\section{The K3 quantum toroidal algebra}
\label{subsec:k3-quantum-toroidal}

The K3 Yangian $Y(\fg_{K3})$ is the \emph{rational} limit of a
two-parameter family: the ``K3 quantum toroidal algebra''
$U_{q,t}(\fg_{K3}^{\mathrm{tor}})$, whose second deformation
parameter~$q = e^{2\pi i \tau}$ is the nome of the elliptic curve~$E$
in $K3 \times E$.  The Yangian limit is $q \to 0$
($\tau \to i\infty$): the $E$ degenerates to a nodal rational curve,
and the second loop collapses.

\subsection{The deformation hierarchy}

The full Drinfeld--Jimbo chain for K3 has four stages:
\begin{center}
\small
\renewcommand{\arraystretch}{1.3}
\begin{tabular}{llll}
\toprule
\textbf{Stage} & \textbf{Algebra} & \textbf{Structure function} & \textbf{Status} \\
\midrule
Classical & $\fg_{K3} = H_{\mathrm{Muk}} = \SpCh_{\Sigma_1, C}(\PhiFA_2(D^b\Coh(K3)))$ & $g = 1$ (trivial) & PROVED (CY-A$_2$) \\
Yangian (rational) & $Y(\fg_{K3})$ & $g_{K3}(u) = \prod \frac{u - h_a}{u + h_a}$ & Conjectural \\
Quantum affine (trig.) & $U_q(\fg_{K3}^{\mathrm{aff}})$ & $G_a(x) = \frac{1 - q_a x}{1 - q_a^{-1} x}$ & Conjectural \\
Quantum toroidal & $U_{q,t}(\fg_{K3}^{\mathrm{tor}})$ & $G_{K3}(x) = \prod_{a=1}^{24} G_a(x)$ & Conjectural \\
\bottomrule
\end{tabular}
\end{center}
The passage from the Yangian (rational) to the quantum toroidal
(trigonometric) stage is the exponentiation
$u \mapsto x = e^u$, $h_a \mapsto q_a = e^{h_a}$,
which replaces the additive spectral parameter by a multiplicative
one and the Arnold relation by the Fay trisecant identity
(Theorem~\ref{thm:fay}).  The CY$_2$ constraint
$\sum h_a = 0$ becomes $\prod q_a = 1$.

\begin{conjecture}[K3 quantum toroidal algebra]
\label{conj:k3-quantum-toroidal}
\ClaimStatusConjectured
There exists a two-parameter algebra
$U_{q,t}(\fg_{K3}^{\mathrm{tor}})$ with the following properties.
\begin{enumerate}[label=\textup{(\roman*)}]
\item \emph{Trigonometric structure function.}
The degree-$(24,24)$ structure function is
\begin{equation}\label{eq:k3-qt-structure-fn}
  G_{K3}(x) \;=\; \prod_{a=1}^{24}
  \frac{1 - q_a x}{1 - q_a^{-1} x},
\end{equation}
where $q_a = e^{h_a}$ are the multiplicative Mukai parameters
satisfying $\prod_{a=1}^{24} q_a = 1$.

\item \emph{Inversion identity.}
$G_{K3}(x) \cdot G_{K3}(x^{-1}) = 1$ (from $\prod q_a = 1$),
the trigonometric analogue of $g_{K3}(u)\, g_{K3}(-u) = 1$.

\item \emph{Abelian factorisation.}
For $\fg = \gl_1$, the algebra decomposes as a product of $24$
rank-$1$ quantum toroidal algebras
$U_{q_a, t_a}(\widehat{\widehat{\gl}}_1)$, one per Mukai lattice
direction, coupled through the Mukai-weighted central extension
$[J_{a,m},\, J_{b,n}] = \omega^{ab}\, m\, \delta_{m+n,0}$.
The cross-family structure function is trivial:
$G_{ab}(x) = 1$ for $a \neq b$ (structure constants vanish for $\gl_1$).

\item \emph{Yangian limit.}
As $\varepsilon \to 0$ with $q_a = e^{\varepsilon h_a}$
and $x = e^{\varepsilon u}$:
\[
  G_{K3}(e^{\varepsilon u})
  \;\longrightarrow\;
  \prod_{a=1}^{24} \frac{u + h_a}{u - h_a}
  \;=\; \frac{1}{g_{K3}(u)},
\]
recovering the inverse of the rational structure function.
(The inversion reflects the convention difference between
the Ding--Iohara presentation and the Yangian RTT presentation.)

\item \emph{Bar Euler product (deformation-invariant).}
The bar Euler product is $\prod_{n \geq 1}(1-q^n)^{24}
= \eta(q)^{24}/q = \Delta(q)/q$, identical to the Yangian value.
Flatness of the trigonometric deformation preserves the PBW
filtration, so the bar Euler product depends only on the
associated graded (= $H_{\mathrm{Muk}}$, class~$G$, $24$ generators).
\emph{Note}: the counting parameter~$q$ in the bar Euler product is the
bar-complex weight grading, not the nome of~$E$.

\item \emph{Charge-$n$ representation.}
At charge $n$, $U_{q,t}(\fg_{K3}^{\mathrm{tor}})$ acts on the
K3~quantum toroidal Fock module~$\cF_n$ with character
\[
  \dim \cF_n \;=\; p_{24}(n)
  \;=\; \chi\bigl(\mathrm{Hilb}^n(K3)\bigr),
\]
the $24$-coloured partition function.
At $n = 1$: $\cF_1 = \bC^{24}$ (Mukai lattice).
At $n = 2$: $\dim \cF_2 = 324$.
The generating function is
$\sum_{n \geq 0} p_{24}(n)\, q^n
= \prod_{k \geq 1} 1/(1 - q^k)^{24}$
(the Fock space character of $H_{\mathrm{Muk}}$).
\end{enumerate}
\emph{Dependency chain}: CY-A$_3$ is proved in the infinity-categorical framework (Theorem~\ref{thm:derived-framing-obstruction}); the geometric construction of the K3 quantum toroidal algebra requires the explicit chain-level data beyond the inf-cat existence.  Conditional on Conjecture~\ref{conj:k3-yangian}
for the existence of the rational limit $Y(\fg_{K3})$.
\end{conjecture}

\subsection{The Miki automorphism: absence of \texorpdfstring{$S_3$}{S-3} for K3}

\begin{proposition}[No \texorpdfstring{$S_3$}{S-3} Miki for the K3 quantum toroidal]
\label{prop:k3-qt-no-s3-miki}
\ClaimStatusConditional
Conditional on Conjecture~\ref{conj:k3-quantum-toroidal}, the K3
quantum toroidal algebra $U_{q,t}(\fg_{K3}^{\mathrm{tor}})$ admits
no $S_3$ Miki automorphism.  It admits an $\SL_2(\Z)$ automorphism
from the mapping class group of~$E$.
\end{proposition}

\begin{proof}
The $S_3$ Miki automorphism of $U_{q,t}(\widehat{\widehat{\gl}}_1)$
arises from the Weyl group $W(T) = S_3$ of the torus
$T = (\bC^*)^3/\bC^*_{\mathrm{diag}}$ acting on~$\bC^3$
(Conjecture~\ref{conj:miki-from-e3}, conditional on the
$E_3$-chiral structure).
For $K3 \times E$, the only $1$-parameter algebraic group acting
is $\mathrm{Aut}(E)^\circ \cong E$ \emph{(}elliptic translation,
\emph{not} an algebraic torus $\mathbb{G}_m$; see
Lemma~\ref{lem:no-Gm-on-E} of \texttt{drinfeld\_center.tex} for
the precise distinction\emph{):} K3 admits no torus action, and
$E$ is an abelian variety, hence the connected automorphism
component of $K3 \times E$ contains no $\mathbb{G}_m$-axis with
which to form a $T = (\C^*)^3$ permuted by $S_3$.  The
\emph{would-be} $T_E$ in the chart-by-chart MO setup of
Proposition~\ref{prop:mo-bypass} \emph{(}\texttt{drinfeld\_center.tex},
ADE/Kummer locus only\emph{)} reduces to the discrete torsion
$\mathrm{Aut}(E)/\mathrm{Aut}(E)^\circ$, generically~$\Z/2\Z$.
The Weyl group of any one-dimensional torus is $\Z/2\Z$,
not~$S_3$.

The mapping class group $\mathrm{MCG}(E) = \SL_2(\Z)$ does act:
\begin{itemize}
  \item $S \colon \tau \mapsto -1/\tau$ inverts the nome
    $q = e^{2\pi i \tau}$.
  \item $T \colon \tau \mapsto \tau + 1$ is the Dehn twist
    ($T$ acts trivially on~$q$ since $e^{2\pi i(\tau+1)} = e^{2\pi i \tau}$;
    it acts nontrivially on the generator level as spectral flow).
\end{itemize}
The $\SL_2(\Z)$ acts on the $E$-modulus only; it does not
permute the $24$ Mukai directions of~K3.
The Miki automorphism is algebra-specific --- it requires the
$(\bC^*)^3$ torus of~$\bC^3$, and is not forced by the $E_3$
operad alone.
\end{proof}

\begin{theorem}[Miki \texorpdfstring{$S_3$}{S-3} on three structurally distinct algebras]
\label{thm:plat-Miki-S3}
\ClaimStatusProvedHere
Fix $(\epsilon_1,\epsilon_2,\epsilon_3)$ with $\epsilon_1+\epsilon_2+\epsilon_3=0$
and let $Y=Y_{\epsilon_1,\epsilon_2,\epsilon_3}(\widehat{\gl}_1)$ denote the
affine Yangian on $\bC^3$. The Miki~$2007$
$S_3$-automorphism of $\mathcal W_{1+\infty}[\lambda]$
descends to three structurally distinct levels, each equivariant under
the cyclic permutation of the epsilon parameters:
\begin{enumerate}[label=\textup{(\alph*)}]
\item \emph{Drinfeld double.} Write $Y=Y^+\otimes Y^0\otimes Y^-$ for
  the triangular decomposition into positive, Cartan, and negative
  halves. The $S_3$-permutation of $(\epsilon_1,\epsilon_2,\epsilon_3)$
  preserves the shuffle kernel and descends to a Hopf-algebra
  automorphism of $Y$, intertwining coproduct, antipode, and counit.
\item \emph{Positive half.} The identification
  $Y^+=\mathrm{CoHA}(\bC^3)$ of Schiffmann--Vasserot~$2013$
  presents $Y^+$ as an associative cohomological-Hall shuffle algebra;
  the $S_3$-action on $(\epsilon_1,\epsilon_2,\epsilon_3)$ is
  shuffle-product-equivariant, and fixes the positive-root cone.
\item \emph{Coalgebra dual.} The coalgebra dual of $Y^+$ is
  realised as a factorisation algebra on the Ran space $\mathrm{Ran}(\bC)$
  (Beilinson--Drinfeld~$2004$ Chapter~$3.4$); the $S_3$-automorphism
  commutes with the Ran fusion product, acting level-wise on the
  symmetric-power stratification $\mathrm{Ran}(\bC)=\bigsqcup_n \mathrm{Sym}^n\bC$.
\end{enumerate}
The $\mathcal W_{1+\infty}[\lambda]$-triality of Gaiotto--Rapcak~$2019$
is the image of the $Y^+$-triality of~(b) under the evaluation
homomorphism $\mathrm{ev}_\lambda\colon Y^+\to\mathcal W_{1+\infty}[\lambda]$:
the $S_3$ on $(\epsilon_1,\epsilon_2,\epsilon_3)$ descends along
$\mathrm{ev}_\lambda$ to the $S_3$ on the $\lambda$-parameter space
of $\mathcal W_{1+\infty}[\lambda]$.
\end{theorem}

\begin{proof}
\emph{(a).} The Hopf structure of $Y$ is controlled by the
Cartan-triangular decomposition; the coproduct
$\Delta\colon Y\to Y\otimes Y$ on each triangular factor is expressed
via shuffle-kernel polynomials symmetric under the $S_3$-permutation of
the epsilon parameters (Negut~$2020$ \emph{Sel.~Math.}~$26$
Theorem~$1.1$), so the cyclic permutation preserves coproduct, antipode,
and counit.

\emph{(b).} The cohomological Hall algebra
$\mathrm{CoHA}(\bC^3)$ is presented as a shuffle algebra on symmetric
rational functions in variables $(z_1,\ldots,z_n)$ with structure
function
\[
 G_{\bC^3}(z)=\prod_{i=1}^{3}(z-\epsilon_i),\qquad \sum_i\epsilon_i=0
 \;\;\Longleftrightarrow\;\; G_{\bC^3}(0)=-\epsilon_1\epsilon_2\epsilon_3,
\]
symmetric under the $S_3$-permutation of the roots of $G_{\bC^3}$
(Schiffmann--Vasserot~$2013$ \emph{Publ.~Math.~IH\'ES}~$118$
Theorem~$1.1$; Rapcak--Soibelman--Yang--Zhao~$2020$
\texttt{arXiv:2007.13365} Theorem~$1.3$). Shuffle multiplication commutes
with the $S_3$ on $(\epsilon_1,\epsilon_2,\epsilon_3)$; the identification
$\mathrm{CoHA}(\bC^3)=Y^+$ (positive half of the affine Yangian, not
$\mathcal W_{1+\infty}$) is strict and equivariant under this action.

\emph{(c).} Beilinson--Drinfeld~$2004$ Chapter~$3.4$ realises the
coalgebra dual of a chiral algebra as a factorisation algebra on the
Ran space; Lurie~$2017$ HA.5.5.4.10 identifies Ran factorisation with
$\bE_1$-chiral structure on the formal disk. The $S_3$ acts as a
relabelling of the three epsilon directions, commuting with Ran fusion
$j_*j^*(A\boxtimes A\to A)$ because fusion is determined by the
symmetric structure function $G_{\bC^3}$ alone.
\end{proof}

\begin{remark}[CoHA(\texorpdfstring{$\bC^3$}{bC-3}) \texorpdfstring{$=Y^+$}{=Y-+}, not \texorpdfstring{$\mathcal W_{1+\infty}$}{W-1+inf}]
\label{rem:coha-c3-equals-yplus-not-Winf}
Theorem~\ref{thm:plat-Miki-S3} confirms the identification
$\mathrm{CoHA}(\bC^3) = Y^+$ (the positive half of the
affine Yangian on $\bC^3$), \emph{not} $\mathcal W_{1+\infty}$
(its evaluation slice). The two structures sit in the exact sequence
\[
 Y^+ = \mathrm{CoHA}(\bC^3)
 \;\xrightarrow{\;\mathrm{ev}_\lambda\;}\;
 \mathcal W_{1+\infty}[\lambda],
\]
where $\mathrm{ev}_\lambda$ is the evaluation homomorphism at a fixed
$\lambda$ (Prochazka~$2016$ \emph{JHEP}~$2016$/$077$; Gaiotto--Rapcak~$2019$).
$Y^+$ is an associative cohomological-Hecke algebra with integrable
$\bE_2$-chiral structure; $\mathcal W_{1+\infty}$ is its
$\mathrm{ev}_\lambda$-image, a conformal vertex algebra on a curve. The
$S_3$-triality of $\mathcal W_{1+\infty}[\lambda]$ observed
physically (Gaiotto--Rapcak) is the shadow under $\mathrm{ev}_\lambda$
of the algebra-level triality of~(b). Conflating the two levels is the
CoHA-versus-vertex-algebra category error; the present theorem
places the triality where its proof actually works.
\end{remark}

\begin{remark}[Non-truncation at CY\texorpdfstring{$_3$}{-3}]
\label{rmk:plat-CY3-non-truncation}
At the Calabi--Yau slice $\sum_{i=1}^{3}\epsilon_i=0$ the qq-character
recursion of Nekrasov--Pestun--Shatashvili~$2018$
(\texttt{arXiv:1312.6689} Section~$2$) does \emph{not} truncate. The
natural guess --- that CY$_3$ forces a finite-rank collapse of the
recursion --- is refuted by the infinite-dimensional content of
$\mathcal W_{1+\infty}[\lambda]$: rather than truncating, the recursion
acquires the $S_3$-triality of Theorem~\ref{thm:plat-Miki-S3} as an
additional symmetry, consistent with $\dim\mathcal W_{1+\infty}=\infty$
and with the Feigin--Jimbo--Miwa--Mukhin~$2016$
(\texttt{arXiv:1603.02765} Theorem~$2.2$) triality theorem. The
surviving content is the $S_3$-symmetric shuffle-algebra structure, not
any polynomial cut-off in qq-character order. In combination with
Proposition~\ref{prop:k3-qt-no-s3-miki}, the K3 promotion further
breaks $S_3$ to $\Z/2\Z$: on $K3\times E$ the two remaining
Mukai-equivariant epsilon directions carry a $\Z/2\Z$ involution
descended from the elliptic-curve Fourier--Mukai duality, not an $S_3$.
\end{remark}

\begin{remark}[What the K3 quantum toroidal adds to the \texorpdfstring{$\bC^3$}{bC-3} picture]
\label{rem:k3-qt-vs-c3}
\index{quantum toroidal algebra!K3 vs C3@K3 vs $\bC^3$}
The structural differences between the $\bC^3$ quantum toroidal
$U_{q,t}(\widehat{\widehat{\gl}}_1)$ and the K3 quantum toroidal
$U_{q,t}(\fg_{K3}^{\mathrm{tor}})$ are:
\begin{center}
\small
\renewcommand{\arraystretch}{1.3}
\begin{tabular}{lll}
\toprule
\textbf{Feature} & \textbf{$\bC^3$} & \textbf{$K3 \times E$} \\
\midrule
Structure function degree & $(3,3)$ & $(24,24)$ \\
Parameters & $h_1, h_2, h_3$ & $h_1, \ldots, h_{24}$ \\
CY constraint & $h_1 + h_2 + h_3 = 0$ & $\sum_{a=1}^{24} h_a = 0$ \\
Mukai signature & $(+,+,+)$ (positive) & $(+^4, -^{20})$ (indefinite) \\
Miki automorphism & $S_3$ (full triality) & $\SL_2(\Z)$ from $E$ only \\
Fock character & $M(q) = \prod 1/(1-q^k)^k$ & $\prod 1/(1-q^k)^{24}$ \\
Bar Euler product & $\eta(q)^3/q$ & $\eta(q)^{24}/q = \Delta(q)/q$ \\
$R$-matrix type & Yang ($3 \times 3$) & Diagonal ($24 \times 24$) \\
Shadow class & $M$ (mock modular) & $G$ (free field, class $G$) \\
\bottomrule
\end{tabular}
\end{center}
The K3 version replaces the equivariant torus $(\bC^*)^3$ by the Mukai
lattice $U^3 \oplus E_8(-1)^2$, gaining rank ($24$ vs~$3$) at the cost
of the $S_3$ symmetry.  The indefinite signature $(4,20)$ poses no
obstruction to the construction (Proposition~\ref{prop:mukai-indefinite-yangian}).
\end{remark}

\noindent\textit{Verification}: $51$~tests in
\texttt{k3\_quantum\_toroidal.py}
covering trigonometric structure function ($8$~tests), inversion
identity ($4$~tests), factorisation ($4$~tests), Yangian limit
($5$~tests), Miki analysis ($5$~tests compliance),
deformation chain ($4$~tests), representation theory ($6$~tests),
bar Euler product ($3$~tests), cross-family structure ($3$~tests),
full verification ($2$~tests), and multi-path cross-checks ($7$~tests).


\noindent\textit{Verification (topological-string shadow cross-check)}:
$65$ tests in \texttt{k3e\_topological\_string\_shadow.py} covering
G\"ottsche's formula through $n = 15$, Ramanujan's $\tau$-function
through $n = 15$ (including Hecke multiplicativity), $\eta^{24}$ and
$1/\eta^{24}$ product expansions, bar-Euler versus G\"ottsche
agreement, $\kappa_\bullet$-spectrum four-distinctness, DT degree-$0$
triviality, the DT/PT identity, BPS/shadow-class identification
(rank-$1$~= G, full~= M), the Borcherds discriminant constraint, and
the wall-crossing/shadow dictionary.

\section{The AGT correspondence for K3: Nekrasov partition function vs conformal blocks}
\label{sec:k3e-agt}

The Alday--Gaiotto--Tachikawa correspondence
(Alday--Gaiotto--Tachikawa~2010, \emph{Lett.~Math.~Phys.}~$91$,
arXiv:0906.3219) on $\bC^2_{\epsilon_1, \epsilon_2}$ identifies the
$U(N)$ instanton Nekrasov partition function with a conformal block of
the extended $W_N \otimes \mathrm{Heis}$ algebra, where $W_N$ is the
Fateev--Lukyanov $\cW$-algebra at central charge
$c(W_N) = (N - 1)(1 + N(N + 1)\, b^2\, b^{-2})$ with
$b^2 = -\epsilon_1/\epsilon_2$. The $\bC^2$ case was proved in two
forms: Schiffmann--Vasserot (arXiv:1202.2756,
``Cherednik algebras, W-algebras and the equivariant cohomology of the
moduli space of instantons on $\bA^2$'', \emph{Publ.\ Math.\ IH\'ES}~118
\emph{(2013)} 213--342) establishes a $W_N \otimes \mathrm{Heis}$-action
on $\bigoplus_n H^*_T(\cM(r, n))$ intertwining the Nekrasov partition
function with the $W_N$-vacuum block; Maulik--Okounkov (arXiv:1211.1287,
``Quantum Groups and Quantum Cohomology'', \emph{Ast\'erisque}~408
\emph{(2019)}) supplies the same action via stable envelopes and
identifies the affine Yangian $Y_\hbar(\widehat{\gl}_N)$ as the
integrable lift whose evaluation slice at fixed central charge
$\lambda$ recovers $W_N$. The rank-free colimit $N \to \infty$ produces
$\cW_{1 + \infty}$ as the evaluation slice of
$\mathrm{CoHA}(\bC^3) = Y^+(\widehat{\gl}_1)$, \emph{not} as the CFT
side of AGT at a single rank $N$
\textup{(}Remark~\ref{rem:coha-c3-equals-yplus-not-Winf}\textup{).}
For $\mathrm K3$, the analogous statement (conjectural at rank $N \geq 2$)
replaces $W_N \otimes \mathrm{Heis}$ by the $24$-copy Mukai-enhanced
$W_N^{\otimes 24}$ with central charge $c_{\mathrm{total}} = 24(N - 1)$;
at rank $1$ the $W_1 = \mathrm{Heis}$ reduction recovers the Mukai
Heisenberg $H_{\mathrm{Muk}}$ at central charge~$c = 24$.

\subsection{Genus-\texorpdfstring{$1$}{1} AGT on K3}
\label{subsec:k3e-agt-genus1}

\begin{proposition}[Genus-\texorpdfstring{$1$}{1} AGT for K3]
\label{prop:k3e-agt-genus1}
\ClaimStatusProvedHere
The rank-$1$ Vafa--Witten partition function on K3 equals the genus-$1$ Heisenberg conformal block character:
\[
  Z_{\mathrm{VW}}(K3, \SU(2); q)
  = \mathrm{ch}(\mathrm{Fock}(H_{\mathrm{Muk}}))
  = \prod_{k \geq 1} \frac{1}{(1 - q^k)^{24}}
  = \frac{1}{\Delta(q)}.
\]
The left side is the G\"ottsche generating function $\sum_n \chi(\mathrm{Hilb}^n(K3)) \, q^n$ (Theorem~\ref{thm:k3e-gottsche-product}). The right side is the Fock space character of $24$ free bosons with the Mukai pairing. Both equal $1/\eta(q)^{24}$ up to the standard $q$-shift.
\end{proposition}
\begin{proof}
Both sides are the generating function of $24$-coloured partitions $p_{24}(n)$, which is the coefficient of $q^n$ in $\prod_{k \geq 1} (1-q^k)^{-24}$. On the VW side, this is Theorem~\ref{thm:k3e-gottsche-product} (G\"ottsche). On the Heisenberg side, the rank-$24$ Fock space has $24$ independent bosonic oscillators at each energy level, giving the same product.
\end{proof}

\noindent This is the \emph{classical} AGT for K3: no quantum group structure is needed. The identification is purely at the level of generating functions.

\noindent\textit{Verification}: 79 tests in \texttt{nekrasov\_agt\_k3.py}; coefficient match through $n = 6$, three independent computation paths (eta product, coloured partition recursion, cross-engine).

\subsection{The refined Vafa--Witten partition function and the \texorpdfstring{$\chi_y$}{chi-y} genus}
\label{subsec:k3e-refined-vw}

The Hirzebruch $\chi_y$ genus of K3 is the polynomial
\begin{equation}\label{eq:k3e-chi-y}
  \chi_y(K3) = \sum_{p=0}^{2} (-1)^p \chi(\Omega^p_{K3}) \, y^p
  = 2 + 20y + 2y^2,
\end{equation}
with special values:
\begin{center}
\begin{tabular}{c|ccc}
$y$ & $1$ & $0$ & $-1$ \\ \hline
$\chi_y(K3)$ & $24 = \kappa_{\mathrm{fiber}}$ & $2 = \kappa_{\mathrm{cat}}$ & $-16 = \operatorname{Tr}(\omega_{\mathrm{Muk}})$
\end{tabular}
\end{center}
The $y = 0$ specialisation recovers $\kappa_{\mathrm{cat}} = \chi(\cO_{K3}) = 2$, while $y = 1$ recovers $\kappa_{\mathrm{fiber}} = \chi(K3) = 24$.  The value $\chi_{-1}(K3) = -16$ is the Mukai trace $\operatorname{Tr}(\omega) = 4 - 20$.

The refined Vafa--Witten partition function in the $\Omega$-background has the form
\[
  Z_{\mathrm{VW}}^{\mathrm{ref}}(K3, q, y) = \prod_{n \geq 1} \frac{1}{(1 - q^n)^{\chi_y(K3)}},
\]
which at $y = 1$ recovers $1/\eta^{24}$ and at $y = 0$ gives $\prod(1-q^n)^{-2}$, the generating function of $2$-coloured partitions.

\begin{remark}[The \texorpdfstring{$\Omega$}{Omega}-background intermediary]
\label{rem:k3e-omega-background}
The refinement parameter~$y$ enters through the $\Omega$-background
equivariant $\chi_y$ genus, not through a direct identification with
Mukai lattice parameters.  The intermediary mechanism is explicit:
equivariant localisation on the $\Omega$-background produces the
Nekrasov partition function~\cite{NekrasovOkounkov2003}, whose
$(\epsilon_1, \epsilon_2)$ plane descends to the single refinement
parameter $y = -\epsilon_2/\epsilon_1$ on the unrefined $\chi_y$
locus. The equivariant blow-up sum of
Nakajima--Yoshioka~\cite{NakajimaYoshioka2003a,NakajimaYoshioka2003b}
extends this to rank~$N$ gauge theory on $\widetilde{\C}^2$, whose K3
automorphic uplift supplies the DMVV generating function of
Conjecture~\ref{conj:k3e-agt-higher-rank}.
\end{remark}

\begin{theorem}[$\Omega$-background free-field chain: K3-side realisation of the stable super-RTT kernel]
\label{rem:k3e-omega-free-field-link}
\ClaimStatusProvedHere
Fix $\mathrm K3 \times \mathbb C^2_{\epsilon_1, \epsilon_2}$ with the
$T = (\C^\times)^2$ torus acting by weights $(\epsilon_1, \epsilon_2)$
on the transverse plane to K3. Write $\hbar$ for the Nekrasov--Okounkov
genus-expansion parameter. Three identifications close on a single
chain.
\begin{enumerate}[label=\textup{(\roman*)}]
\item \emph{Free-field kernel equals super-RTT kernel.} The twisted OPE
\begin{equation}\label{eq:k3qt-omega-ope-chain}
  \phi_1(z)\, \phi_2(w) \;\sim\;
  \frac{\hbar^2}{\epsilon_1\epsilon_2}\,\log(z - w),
  \qquad \phi_i(z)\,\phi_i(w)\sim 0,
\end{equation}
of Equation~(\ref{eq:omega-background-ope}) of
Chapter~\ref{ch:k3-yangian}, Subsection
\ref{subsec:omega-background-free-field}, produces the two-point
function
$\langle t_{12}(u)\, t_{21}(v)\rangle_\Omega =
(u - v + \epsilon_1)(u - v + \epsilon_2)/((u - v)(u - v + \epsilon_1 + \epsilon_2))$
on $t_{12} = {:}e^{\phi_1 - \phi_2}{:}$,
$t_{21} = {:}e^{\phi_2 - \phi_1}{:}$.
Theorem~\ref{thm:omega-bg-two-point-match} identifies this two-point
function with the matrix element
$\langle e_1 \otimes e_0 | R_{\mathrm{KS}}(u - v) | e_0 \otimes e_1\rangle$
of the Kulish--Sklyanin super $R$-matrix on
$\C^{1 \mid 1} \otimes \C^{1 \mid 1}$ under the identification
$\hbar = \epsilon_1\epsilon_2/(\epsilon_1 + \epsilon_2)$.
\item \emph{CY$_3$ closure identifies the third weight.} Under the
CY$_3$ closure $\epsilon_1 + \epsilon_2 + \epsilon_3 = 0$ the
denominator $(u - v + \epsilon_1 + \epsilon_2) = (u - v - \epsilon_3)$
supplies the third spectral direction, and the Nekrasov--Okounkov
expansion parameter becomes
\[
  \hbar \;=\; \frac{\epsilon_1\epsilon_2}{\epsilon_1 + \epsilon_2}
      \;=\; -\frac{\epsilon_1\epsilon_2}{\epsilon_3},
\]
the genus-expansion coupling of
\cite{NekrasovOkounkov2003}.
\item \emph{Automorphic uplift to $\mathrm K3 \times E$.} The rank-$N$
equivariant instanton sum of Nakajima--Yoshioka
\cite{NakajimaYoshioka2003a,NakajimaYoshioka2003b} on
$\widetilde{\C}^2 = \mathrm{Bl}_0\,\C^2$ admits an automorphic
resummation producing the genus-$1$ partition function on
$\mathrm K3 \times E$; at rank $N = 1$ the resummation is
G\"ottsche's formula
$\sum_n \chi(\mathrm{Hilb}^n(\mathrm K3))\, q^n = 1/\eta(q)^{24}$.
The $24$-copy $(\epsilon_1, \epsilon_2)$-deformation lifts the chain
$(i)$--$(ii)$ to a Mukai-equivariant free-field system
$\{\phi_a(z)\}_{a = 1, \ldots, 24}$ diagonalising the Mukai pairing,
and the orthosymplectic involution on the $(4, 20)$-signature reduces
the $576$-entry $T$-matrix of $Y(\fgl(4 \mid 20))$ to the
$296$-entry $T$-matrix of $Y_{\osp(4 \mid 20)}$
\textup{(}Chapter~\ref{ch:k3-yangian},
Remark~\ref{rem:ks-to-osp-lift}\textup{).}
\end{enumerate}
Under $(q, t) = (e^{\hbar\epsilon_1}, e^{-\hbar\epsilon_2})$ with
$z/w = e^{\hbar(u - v)}$, the Macdonald--Shiraishi $(q, t)$-refinement
of Subsection~\ref{subsec:macdonald-shiraishi-qt} \textup{(}Chapter
\ref{ch:k3-yangian}\textup{)} lifts the rational kernel to the refined
vertex of Awata--Kanno \cite{AwataKanno2008}. In the Nekrasov--Shatashvili
limit $\epsilon_2 \to 0$, Theorem~\ref{thm:ns-limit-xxx}
\textup{(}Chapter~\ref{ch:k3-yangian}\textup{)} degenerates the kernel
to the Yang $R$-matrix $(u - v + \epsilon_1)/(u - v)$ of the
$\widehat{\gl}(1 \mid 1)$ XXX spin chain, so the
Nekrasov--Shatashvili \cite{NekrasovShatashvili2009} free energy is
the super-Whittaker vector of the Kulish--Sklyanin super-Yangian
$Y_{\mathrm{KS}}(\gl(1 \mid 1))$ at $\gl(1 \mid 1)$ rank.
\end{theorem}

\begin{proof}
(i) is Theorem~\ref{thm:omega-bg-two-point-match} of
Chapter~\ref{ch:k3-yangian}. (ii) is the algebraic identity
$\epsilon_1 + \epsilon_2 = -\epsilon_3$ substituted into the $\hbar$-identification
of (i). (iii) combines G\"ottsche~$1990$ with the
Nakajima--Yoshioka~$2003$ equivariant blow-up formula and the
Mukai-diagonalisation of Remark~\ref{rem:ks-to-osp-lift}; the
orthosymplectic reduction is the signature $(4, 20)$ applied to
the $24$-copy Kulish--Sklyanin kernel. The $(q, t)$-lift and the NS
degeneration are
Proposition~\ref{prop:qt-to-omega} and
Theorem~\ref{thm:ns-limit-xxx} of Chapter~\ref{ch:k3-yangian}.
\end{proof}

\subsection{The conjectural K3 Yangian character}
\label{subsec:k3e-yangian-character-agt}

\begin{conjecture}[K3 Yangian character = VW partition function]
\label{conj:k3e-yangian-agt}
\ClaimStatusConjectured
If the K3 Yangian $Y(\fg_{K3})$ exists (Conjecture~\ref{conj:k3-yangian}), then
\[
  \mathrm{ch}(\mathrm{Fock}(Y(\fg_{K3})))
  = Z_{\mathrm{VW}}(K3, \SU(2); q)
  = \frac{1}{\Delta(q)}.
\]
\end{conjecture}

\noindent\emph{Evidence.} For the abelian ($\gl_1$) K3 Yangian, the flat deformation hypothesis (PBW filtration preserved) implies the Fock character is deformation-invariant: $\mathrm{ch}(\mathrm{Fock}(Y(\fg_{K3}))) = \mathrm{ch}(\mathrm{Fock}(H_{\mathrm{Muk}})) = 1/\Delta(q)$. The prediction becomes nontrivial at the non-abelian level.

\subsection{The Kummer point and the orbifold AGT}
\label{subsec:k3e-kummer-agt}

At the Kummer point $K3 = T^4/\bZ_2$ (resolved), the VW partition function decomposes via the orbifold formula. The Kummer route (Steps~1--4, Proposition~\ref{prop:kummer-orbifold}, PROVED) gives:
\begin{enumerate}[label=\textbf{Step~\arabic*}:]
\item $T^4$ Heisenberg: rank-$16$ lattice from $b_\ast(T^4) = 1 + 4 + 6 + 4 + 1 = 16$.
\item $\bZ_2$-invariant sublattice: rank~$8$.
\item $16$ twisted sectors from the $2^4 = 16$ fixed points, each at conformal weight $h = 1/2$.
\item Orbifold total: $\kappa_{\mathrm{ch}} = 2 = \chi(\cO_{K3})$; exponent $24 = 8 + 16$.
\end{enumerate}
The Kummer Yangian parameters specialise through the
$\mathbb{Z}_2$-equivariant reduction of the $\Omega$-background
free-field system of
Theorem~\ref{rem:k3e-omega-free-field-link}. On the resolved $A_1$
singularity at each of the $16$ fixed points, the $T = (\C^\times)^2$
action descends to the $\mathbb{Z}_2$-quotient with equivariant weights
$(\epsilon_1, -\epsilon_1)$ on the exceptional $\bP^1$; the resulting
Kummer Yangian parameters are
$\{h_a\}_{a = 1, \ldots, 24} =
\{\epsilon_1\}^{\times 4}
 \cup \{-\epsilon_1/5\}^{\times 20}$
on the $(4, 20)$-Mukai decomposition, with
$\sum_a h_a = 4\epsilon_1 - 4\epsilon_1 = 0$ (CY$_2$~closure). The
rank-$4$ positive Mukai sublattice $U^2 \subset \mathrm{II}_{4, 20}$
carries the $T^4$-Heisenberg of Step~1 (four untwisted creation
operators) and the rank-$20$ negative Mukai sublattice
$E_8(-1)^2 \oplus \langle -2\rangle^4$ carries the twisted-sector
states of Step~3; the ratio $1 : -1/5 = 4 : 1$ tracks the $(4, 20)$
signature. The Fock character
$\prod_{k \geq 1}(1 - q^k)^{-24}$ is independent of parameter choice
by flat deformation invariance, so the orbifold Kummer character
agrees with the generic K3 character.

\begin{remark}[Kummer Yangian identification]
\label{rem:k3e-kummer-yangian-agt}
The identification of the Kummer Yangian character with $Z_{\mathrm{VW}}$ at the Kummer point (Step~5 of the Kummer route, \texttt{conj:kummer-route}) remains conjectural. The match of generating functions is guaranteed by flat deformation, but the match of the \emph{algebraic structure} (R-matrix, coproduct, fusion rules) is a deeper statement.
\end{remark}

\subsection{Modularity: VW weight \texorpdfstring{$-12$}{-12} and the K3 Yangian character}
\label{subsec:k3e-agt-modularity}

The function $\eta(\tau)^{-24}$ transforms under $\mathrm{SL}(2, \bZ)$ with weight $-24/2 = -12$. This is the standard formula: $Z_{\mathrm{VW}}(S; \tau)$ has modular weight $-\chi(S)/2$ (Vafa--Witten 1994). For K3, $\chi = 24$ gives weight~$-12$.

S-duality ($\tau \mapsto -1/\tau$) exchanges gauge groups:
\[
  Z_{\mathrm{VW}}(K3, \SU(2); -1/\tau) = \tau^{-12} \, Z_{\mathrm{VW}}(K3, \SO(3); \tau).
\]
Under the Koszul duality interpretation (Volume~I, Theorem~A): $\SU(2)$ and $\SO(3) = \SU(2)/\bZ_2$ are Langlands dual, and the S-duality map is the Koszul involution $A \leftrightarrow A^!$.

At genus~$g$, the Heisenberg conformal block has modular weight $-12g$ on the moduli space $\cM_g$. At genus~$2$, the block is a Siegel modular form on $\cH_2$ under $\mathrm{Sp}(4, \bZ)$, related to $1/\Delta_5$ via the DMVV formula.

\subsection{Wall-crossing and K3 Yangian fusion rules}
\label{subsec:k3e-agt-wall-crossing}

For rank~$1$ (Hilbert schemes), the VW partition function is stability-independent ($\mathrm{Hilb}^n(K3)$ is smooth for all~$n$; Fogarty 1968), so there is no wall-crossing and no fusion rules to test.

For rank~$\geq 2$, the walls of marginal stability in the Bridgeland stability manifold $\mathrm{Stab}(K3)$ produce nontrivial jumps governed by the Kontsevich--Soibelman formula. The conjectural K3 Yangian interpretation:
\begin{itemize}
\item The wall-crossing factors correspond to K3 Yangian R-matrix data.
\item For the abelian ($\gl_1$) Yangian, the R-matrix poles at $z = -h_i$ correspond to walls.
\item Higher-rank wall-crossing involves the charge-$2$ R-matrix, which is $576 \times 576$ ($24^2 \times 24^2$).
\end{itemize}

\begin{remark}[Higher coherence obstruction]
\label{rem:k3e-agt-higher-coherence}
Pairwise R-matrix consistency (the Yang--Baxter equation) does not
automatically imply higher-order consistency (the Zamolodchikov
tetrahedron equation).  Charge-$3$ wall-crossing therefore involves
$E_3$ corrections beyond pairwise products
(Theorem~\ref{thm:zte-failure}).
\end{remark}

\noindent\textit{Verification}: \texttt{nekrasov\_agt\_k3.py}, 79 tests: genus-$1$ AGT identification, $\chi_y$ specialisations, modular weight $-12$, Kummer decomposition $24 = 8 + 16$, Ramanujan $\tau$ through $n = 10$ (including Hecke multiplicativity), BPS asymptotic growth, refined VW at $y = 0$ and $y = 1$, cross-consistency with four independent engines (\texttt{vafa\_witten\_k3}, \texttt{vafa\_witten\_shadow}, \texttt{k3\_double\_current\_algebra}, \texttt{k3\_yangian}), and nine multi-path verification tests.

\begin{conjecture}[Higher-rank AGT for K3: from \texorpdfstring{$\fgl_1$}{fgl-1} to \texorpdfstring{$\fgl_N$}{fgl-N}]
\label{conj:k3e-agt-higher-rank}
\ClaimStatusConjectured
For $N$ D3-branes wrapping K3, the rank-$N$ Vafa--Witten partition function
$Z_{\mathrm{VW}}^{U(N)}(K3; q)$ satisfies the following.
\begin{enumerate}[label=\textup{(\alph*)}]
  \item \emph{Rank-$1$ baseline \textup{(proved)}.}
  $Z_{\mathrm{VW}}^{U(1)}(K3; q) = \prod_{n \ge 1} 1/(1 - q^n)^{24}
  = 1/\eta(\tau)^{24}$, with
  $\chi(\operatorname{Hilb}^n(K3)) = p_{24}(n)$ (G\"ottsche~1990).

  \item \emph{Higher-rank AGT \textup{(conjectural)}.}
  $Z_{\mathrm{VW}}^{U(N)}(K3; q) = Z_{W_N^{\otimes 24}}(\Sigma_1, q)$,
  where $W_N^{\otimes 24}$ is $24$ copies of the $W_N$-algebra (one per
  Mukai direction) at the self-dual central charge $c(W_N) = N - 1$
  per copy, so $c_{\mathrm{total}} = 24(N - 1)$.  At $N = 1$ this reduces
  to the rank-$24$ Heisenberg ($W_1 = \mathrm{Heis}$); at $N = 2$, to
  $24$ copies of the Virasoro algebra.

  \item \emph{Structure function scaling.}
  The rank-$N$ structure function
  $g_N(u) = g_{K3}(u)^N = \prod_{a=1}^{24}
  ((u - h_a)/(u + h_a))^N$ has degree $(24N, 24N)$
  and satisfies unitarity $g_N(u) \cdot g_N(-u) = 1$.
  The Newton power sums scale as $p_k^{(N)} = N \cdot p_k^{(1)}$.

  \item \emph{Nonabelian $R$-matrix at $N \ge 2$.}
  The full $R$-matrix on $\C^{24N} \otimes \C^{24N}$ factorises as
  $R_{K3,N}(u) = R_{\fgl_N}(u) \otimes g_{K3}(u)$,
  where $R_{\fgl_N}(u) = I + P_N/u$ is the standard Yang $R$-matrix
  and $P_N$ is the $N \times N$ permutation operator.  At $N \ge 2$
  the Yang--Baxter equation is nontrivial (off-diagonal blocks from
  $\fgl_N$).

  \item \emph{Higher-rank bar Euler product.}
  For the free-field (class~$G$) sector at rank~$N$, the bar Euler
  product is $\prod_{n \ge 1}(1 - q^n)^{-24N}$ (the Heisenberg
  contribution).  The interacting $W_N$ sector contributes additional
  spin-$s$ factors for $s = 2, \ldots, N$, each with exponent~$-24$.

  \item \emph{Kummer route at rank $N$.}
  At the Kummer point $T^4/\bZ_2$, the rank-$N$ Yangian reduces to
  $Y(\fgl_N \otimes \frakg_{\mathrm{Kum}})$ with $\frakg_{\mathrm{Kum}}$
  the rank-$8$ enhanced Kummer algebra, giving orbifold
  VW partition functions computable from the $A_1$ McKay quiver
  at rank~$N$.
\end{enumerate}

The DMVV formula (Dijkgraaf--Moore--Verlinde--Verlinde, \texttt{hep-th/9607026})
provides the all-rank generating function
$\sum_{N \ge 0} p^N Z_{\mathrm{VW}}^{U(N)} =
\prod_{n,k,l}(1 - p^n q^k y^l)^{-c(4nk - l^2)}$, from which
the rank-$N$ coefficient is extracted.  The AGT identification~(b)
is the conjectural algebraic interpretation of this coefficient
as a $W_N^{\otimes 24}$ genus-$1$ partition function.
\end{conjecture}

\noindent\textit{Verification}: $127$ tests in
\texttt{agt\_higher\_rank\_k3.py} covering rank-$1$ VW baseline
($p_{24}(n)$ for $n = 0, \ldots, 6$), rank-$2$ VW from symmetric
product and DMVV, $W_N$ vacuum characters at $N = 1, 2, 3, 4$,
AGT matching at genus~$1$ for $N = 1, 2$, structure function degree
scaling $(24N, 24N)$, unitarity, nonabelian $R$-matrix block
structure, central charge $c = 24(N-1)$, log expansion coefficients,
bar-Euler product at ranks $1$--$4$, and the Kummer route at rank $2$.

%% =====================================================================
\section{The twisted M-theory compactification web}
\label{sec:twisted-m-theory-web}
%% =====================================================================

The string/M-theory compactification landscape on $K3$ and $K3 \times E$
organizes into a web of dimensional reductions, topological twists, and
dualities.  At each node of this web we extract the chiral algebra,
$\kappa_\bullet$-spectrum, shadow class, and $R$-matrix data, obtaining
independent consistency checks.

\subsection{The web}
\label{sec:compactification-web}

The nodes of the web, stratified by the effective noncompact dimension:

\begin{center}
\renewcommand{\arraystretch}{1.3}
\begin{tabular}{lllll}
\hline
\textbf{Level} & \textbf{Node} & \textbf{Chiral algebra} & \textbf{$E_n$} & \textbf{Status} \\
\hline
$11$d & M-theory on $K3$ & $H_{\mathrm{Muk}}$ & $E_1$ & PROVED \\
$10$d & IIA on $K3$ & $6$d $(2{,}0)$ boundary & $E_2$ & PROVED \\
$10$d & IIB on $K3$ & $6$d $(2{,}0)$ (T-dual) & $E_2$ & PROVED \\
$6$d & hCS on $\mathbb{C}^3$ & $U_{q,t}(\widehat{\widehat{\gl}}_1)$ & $E_3$ & CONJECTURAL \\
$6$d & $(2{,}0)$ on $K3 \times C$ & $(4{,}4)$ sigma model & $E_2$ & PROVED \\
$6$d & hCS on $K3 \times E \times C$ & $U_{q,t}(\mathfrak{g}_{K3}^{\mathrm{tor}})$ & $E_3$ & CONJECTURAL \\
$5$d & $K3 \times C \times \mathbb{R}$ & $Y(\mathfrak{g}_{K3})$ & $E_2$ & CONJECTURAL \\
$5$d & $(2{,}0)$ on $S^1$ & $Y(\widehat{\gl}_1)$ & $E_2$ & PROVED \\
$4$d & IIA on $K3 \times T^2$ & $4$d $\mathcal{N}=4$ SYM (VW) & $E_1$ & PROVED \\
$4$d & IIB on $K3 \times T^2$ & $4$d $\mathcal{N}=4$ (S-dual) & $E_1$ & PROVED \\
$3$d & M-theory on $K3 \times T^3$ & RW(Hilb$^N$($K3$)) & $E_\infty$ & PROVED \\
$0$d & $K3$ sigma model (bosonic sector) & $H_{\mathrm{Muk}}$ & $E_1$ & PROVED \\
\hline
\end{tabular}
\end{center}

\subsection{The \texorpdfstring{$\kappa_\bullet$}{kappa-.}-spectrum across the web}
\label{sec:kappa-spectrum-web}

\noindent\textbf{Local convention.} In this subsection $\kappa_\bullet$
abbreviates the four-element set
$\{\kappa_{\mathrm{cat}}, \kappa_{\mathrm{fiber}}, \kappa_{\mathrm{ch}},
  \kappa_{\mathrm{BKM}}\}$; every value appearing below is one of
these four subscripted invariants.

The four $\kappa_\bullet$-invariants are extracted at each node.
They fall into two types: manifold invariants (topological,
independent of algebraization) and algebraization invariants (dependent
on the constructed algebra).  For the $K3 \times E$ tower:

\begin{center}
\renewcommand{\arraystretch}{1.3}
\begin{tabular}{lcccl}
\hline
\textbf{Invariant} & \textbf{K3} & \textbf{$K3 \times E$} & \textbf{Origin} & \textbf{Type} \\
\hline
$\kappa_{\mathrm{cat}}$ & $2$ & $0$ & $\chi(\mathcal{O}_X)$ (Kunneth) & manifold \\
$\kappa_{\mathrm{fiber}}$ & $24$ & $24$ & Mukai lattice rank & manifold \\
$\kappa_{\mathrm{ch}}^{\mathrm{Heis}}$ & $2$ & $3$ & Heisenberg-additive reading & algebraization \\
$\kappa_{\mathrm{BKM}}$ & --- & $5$ & Weight of $\Delta_5$ & algebraization \\
\hline
\end{tabular}
\end{center}

\noindent
Two readings of $\kappa_{\mathrm{ch}}$ coexist: the Heisenberg-additive
reading $\kappa_{\mathrm{ch}}^{\mathrm{Heis}}(X \times Y) =
\kappa_{\mathrm{ch}}^{\mathrm{Heis}}(X) + \kappa_{\mathrm{ch}}^{\mathrm{Heis}}(Y)$
used above, and the Hodge-supertrace reading
$\kappa_{\mathrm{ch}}^{\mathrm{Hodge}}(X) = \chi(\mathcal{O}_X)$
(Kunneth-multiplicative, coinciding with $\kappa_{\mathrm{cat}}$).  They
agree on compact $\mathrm{CY}_d$ with $h^{1, 0} = 0$ and split at
$\mathrm{CY}_3$: on $K3 \times E$ one has
$\kappa_{\mathrm{ch}}^{\mathrm{Heis}} = 3$ while
$\kappa_{\mathrm{ch}}^{\mathrm{Hodge}} = \chi(\mathcal{O}_{K3 \times E}) = 0$.

\noindent
The decomposition
\begin{equation}
\label{eq:kappa-bkm-decomposition}
  \kappa_{\mathrm{BKM}}(K3 \times E)
    = \kappa_{\mathrm{ch}}^{\mathrm{Heis}}(K3 \times E) + \kappa_{\mathrm{cat}}(K3)
    = 3 + 2 = 5
\end{equation}
holds numerically for $K3 \times E$ (see~\texttt{twisted\_m\_theory\_web.py}).
\emph{Warning} (Remark~\ref{rem:bkm-decomposition-adversarial}): this decomposition is a \emph{numerical coincidence} specific to $N = 1$.
For $\bZ/N\bZ$-orbifolds with $N \geq 2$, it fails; the correct universal formula is $\kappa_{\mathrm{BKM}} = c_N(0)/2$ (Borcherds weight theorem, Proposition~\ref{prop:bkm-weight-universal}).

\subsection{\texorpdfstring{$R$}{R}-matrix degeneration chain}
\label{sec:rmatrix-degeneration}

The structure function degenerates along a three-stage reduction
passing through the $(q, t)$-trigonometric Awata--Kanno stratum,
matching the $\Omega$-background free-field realisation of
Theorem~\ref{rem:k3e-omega-free-field-link}:
\[
  \underset{\text{elliptic}}{G^{\mathrm{ell}}_{\mathrm{K3}}(x, y; \{q_a\}, \tau)}
  \;\xrightarrow{\;\tau \to i\infty\;}\;
  \underset{(q, t)\text{-trig}}{G^{(q, t)}_{\mathrm{K3}}(x; \{q_a\})}
  \;\xrightarrow{\;\hbar \to 0\;}\;
  \underset{\text{rational}}{g_{\mathrm{K3}}(u; \{h_a\})}
  \;\xrightarrow{\;h_i \to 0\;}\;
  \underset{\text{trivial}}{1},
\]
with the stages fixed as follows:
\begin{itemize}
\item \textbf{Elliptic (doubly-affine):}
  $G^{\mathrm{ell}}_{\mathrm{K3}}(x, y; \{q_a\}, \tau)
  = \prod_{a = 1}^{24}
  \vartheta_1(x q_a\mid \tau)/\vartheta_1(x q_a^{-1}\mid \tau)$,
  with $\prod q_a = 1$ (CY$_2$ closure).
  Primary: Okounkov~$2017$ \texttt{arXiv:1512.07363} (elliptic
  stable envelopes and $\vartheta$-function structure constants) and
  Smirnov~\cite{Smirnov2020} (elliptic stable envelopes and
  $3$-dimensional mirror symmetry for superspin chains) for the
  elliptic stable-envelope presentation.
\item \textbf{$(q, t)$-trigonometric (Awata--Kanno):} setting
  $\tau \to i\infty$ reduces $\vartheta_1(x; \tau)$ to $1 - x$, and
  $G^{(q, t)}_{\mathrm{K3}}(x) = \prod_a (1 - q_a x)/(1 - q_a^{-1} x)$
  with $(q, t) = (e^{\hbar\epsilon_1}, e^{-\hbar\epsilon_2})$ and
  $\prod q_a = 1$. The refined vertex of Awata--Kanno
  \cite{AwataKanno2008} realises this kernel; the Zamolodchikov
  factorisation
  $R(u, v) = R_1(u)\, R_2(v)\, R_{12}(u - v)$ encodes the $E_3$-crossing
  correction on the $\bC^3$ local model, with the $24$-copy Mukai
  enhancement removing the $S_3$ triality
  (Proposition~\ref{prop:k3-qt-no-s3-miki}).
\item \textbf{Rational (Yangian):} the $\hbar \to 0$ limit with
  $q_a = e^{\hbar h_a}$ and $x = e^{\hbar u}$ reduces
  $G^{(q, t)}_{\mathrm{K3}}(x)$ to
  $g_{\mathrm{K3}}(u) = \prod_{a = 1}^{24} (u - h_a)/(u + h_a)$
  with $\sum h_a = 0$ and degree $(24, 24)$. This is the rational
  Yangian structure function of
  Conjecture~\ref{conj:k3-quantum-toroidal}(iv).
\item \textbf{Trivial (Heisenberg):} all $h_i \to 0$ forces
  $g_{\mathrm{K3}}(u) = 1$; the Yangian degenerates to the current
  algebra $U(H_{\mathrm{Muk}})$ of the rank-$24$ Mukai Heisenberg.
\end{itemize}

\subsection{Shadow class stability}
\label{sec:shadow-class-stability}

Under proper dimensional reduction (shrinking a direction), the shadow class
can only \emph{decrease}:
\[
  M \to C \to L \to G.
\]
Topological twists can \emph{increase} the shadow class: the twist from
$10$d IIA $(G)$ to $6$d hCS $(L)$ introduces the Yangian deformation,
promoting the shadow depth from~$2$ to~$3$.

The key reduction chain in the hCS hierarchy:
\begin{center}
\renewcommand{\arraystretch}{1.3}
\begin{tabular}{lccc}
\hline
\textbf{Node} & \textbf{Shadow class} & \textbf{$R$-matrix} & \textbf{Edge type} \\
\hline
$6$d: $K3 \times E \times C$ & $M$ & two-parameter & --- \\
$\downarrow$ (shrink $E$) & & & reduction \\
$5$d: $K3 \times C \times \mathbb{R}$ & $L$ & rational & --- \\
$\downarrow$ (shrink $C$) & & & reduction \\
$0$d: $K3$ sigma model & $G$ & trivial & --- \\
\hline
\end{tabular}
\end{center}

\subsection{Dimensional reduction and the correspondence programme \texorpdfstring{$\{\Phi_d\}$}{Phi-d}}
\label{sec:dimred-commutativity}

The central consistency question: does the CY-to-chiral correspondence
programme $\{\Phi_d\}$ commute with dimensional reduction?
\[
\begin{tikzcd}
  \mathcal{C}\!\mathit{Y}_d(X) \ar[r, "\Phi"] \ar[d, "\text{reduce}"]
    & A_X \ar[d, "\text{reduce}"] \\
  \mathcal{C}\!\mathit{Y}_{d-1}(X') \ar[r, "\Phi"]
    & A_{X'}
\end{tikzcd}
\]

\begin{conjecture}[Compactification commutativity]
\label{conj:compactification-commutativity}
\ClaimStatusConjectured
The correspondence programme $\{\Phi_d\}$ commutes with the
$6$d~$\to$~$5$d reduction (shrink $E$, $q \to 1$), in the sense that
$\Phi_3(K3 \times E)|_{q \to 1} = \Phi_2(K3)$:
\[
  \Phi_3(K3 \times E)\big|_{q \to 1}
  \;=\;
  \Phi_2(K3)
\]
at the level of algebras, i.e.\
$U_{q,t}(\mathfrak{g}_{K3}^{\mathrm{tor}})\big|_{q \to 1}
= Y(\mathfrak{g}_{K3})$.  Object-level $\Phi_3$ is proved in the
infinity-categorical framework
(Theorem~\ref{thm:cy-to-chiral-d3} via
Theorem~\ref{thm:derived-framing-obstruction}); the explicit
chain-level geometric construction, and the $q \to 1$ commutativity
itself, remain conditional.
\end{conjecture}

\noindent At $d = 2$, the $5$d $\to$ $3$d commutativity (shrink $C$:
$h_i \to 0$) is \emph{proved}: $Y(\mathfrak{g}_{K3})\big|_{h_i = 0}
= H_{\mathrm{Muk}} = \SpCh_{\Sigma_1, C}(\PhiFA_2(D^b\Coh(K3)))$ via CY-A$_2$.
The K3 quantum toroidal algebra of this chapter is a stage-$2$ shadow of the canonical $E_2$-hFA $\PhiFA_2(D^b\Coh(K3))$; the Mukai signature $(4,20)$ is the Nakajima-cycle specialisation datum.

\subsection{Independent consistency checks}
\label{sec:web-consistency-checks}

The compactification web provides the following independent checks,
all verified in \texttt{twisted\_m\_theory\_web.py} ($76$~tests):

\begin{enumerate}
\item \textbf{$\kappa_{\mathrm{ch}}$ preservation} under proper reductions:
  $\kappa_{\mathrm{ch}}$ is constant along reduction edges with the same
  CY dimension.

\item \textbf{$\kappa_{\mathrm{BKM}}$ value}:
  $\kappa_{\mathrm{BKM}} = c_1(0)/2 = 5$ (Borcherds weight theorem); numerically $= 3 + 2$ (\eqref{eq:kappa-bkm-decomposition}, coincidence for $N = 1$ only).

\item \textbf{$R$-matrix degeneration}: two-parameter $\to$ rational
  $\to$ trivial along the hCS hierarchy.

\item \textbf{Shadow class monotonicity}: $M \to L \to G$ along
  proper reductions.

\item \textbf{Euler characteristic}: $\chi(K3 \times E)
  = \chi(K3)\cdot \chi(E) = 24 \cdot 0 = 0$.

\item \textbf{Duality invariance}: T-duality (IIA/IIB) and S-duality
  ($4$d $\mathcal{N}=4$) preserve the $\kappa_\bullet$-spectrum.

\item \textbf{CY condition}: $h_1 + h_2 + h_3 = 0$ is enforced at
  every level by anomaly cancellation of the holomorphic CS theory.
\end{enumerate}

\subsection{Physical K3 sigma model cross-check}
\label{sec:k3-sigma-model-cross-check}

The $\cN = (4,4)$ SCFT on $K3$ at $c = 6$ is \emph{proved}
(Aspinwall--Morrison, Nahm--Wendland).  The K3 Yangian $Y(\fg_{K3})$
is \emph{conjectural} (the geometric construction of $Y(\fg_{K3})$ requires more than the inf-cat existence of CY-A$_3$, which is proved).  The sigma model provides seven
independent data that any candidate K3 chiral algebra must match.
All seven are verified in
\texttt{physical\_k3\_sigma\_check.py} ($73$~tests).

\begin{enumerate}
\item \textbf{$\cN = 4$ SCA at $c = 6$, $k_R = 1$.}
  The left-moving chiral algebra of the sigma model contains the
  small $\cN = 4$ superconformal algebra at $c = 6 k_R$ with $k_R = 1$
  (unitarity).  The Yangian acts on Heisenberg (Fock) modules; the
  $\cN = 4$ constraint restricts which modules are physical BPS states.
  The bosonic sector ($c_{\mathrm{Heis}} = 24$, Mukai lattice) and the
  sigma model ($c = 6 = 3 \dim_{\mathbb{C}} K3$) are \emph{different}
  central charges: the former counts all lattice directions, the latter
  includes the superconformal projection.

\item \textbf{Elliptic genus $\phi_{0,1}$ vs.\ bar Euler product.}
  The K3 elliptic genus is $Z_{K3} = 2\,\phi_{0,1}$, where
  $\phi_{0,1}(\tau, z)$ is the unique weak Jacobi form of weight~$0$,
  index~$1$, and the factor $2 = \kappa_{\mathrm{ch}}(K3)$.
  At $q^0$: $\phi_{0,1} = y^{-1} + 10 + y$,
  so $\kappa_{\mathrm{ch}} \cdot (1 + 10 + 1) = 24 = \chi_{\mathrm{top}}(K3)$.
  The bar Euler product $\prod(1 - q^n)^{24} = \eta(q)^{24}/q$ computes
  the Ramanujan $\tau$-function.
  These are \emph{different} modular objects: $\phi_{0,1}$ (Jacobi form,
  input to the Borcherds lift) vs.\ $1/\eta^{24}$ (modular form, output
  of the bar complex).  The Borcherds lift connects them:
  $\phi_{0,1} \mapsto \Delta_5$; the rank-$0$ restriction of
  $1/\Delta_5$ gives $1/\eta^{24}$.

\item \textbf{BPS spectrum from discriminant coefficients.}
  The discriminant function $c(D)$ of $\phi_{0,1}$ ($D = 4n - l^2$,
  index~$1$) gives BPS multiplicities: $c(-1) = 1$, $c(0) = 10$
  (normalized $\phi_{0,1}$).
  The Borcherds weight is $c(0)/2 = 5 = \kappa_{\mathrm{BKM}}$.
  Only discriminants $D \equiv 0$ or $3 \pmod{4}$ appear.
  Negative values ($c(3) = -64$) indicate fermionic/odd BKM roots.
  The Witten index is $\kappa_{\mathrm{ch}} \cdot c(-1) = 2$.

\item \textbf{Spectral parameter origin}.
  The Yangian spectral parameter $u$ arises from the
  $\Omega$-background deformation of
  $\mathrm K3 \times \C^2_{\epsilon_1, \epsilon_2}$,
  \emph{not} from the sigma model: under the complex exponential
  $z(u) = e^{u/\hbar}$ of
  Theorem~\ref{thm:omega-bg-two-point-match}
  \textup{(}Chapter~\ref{ch:k3-yangian}\textup{),}
  $u$ is the additive holomorphic coordinate on the $\C^2$ transverse
  plane and $z$ is its multiplicative lift to the twisted free-field
  OPE \eqref{eq:omega-background-ope}.
  The $24$-copy Mukai equivariant free-field system produces the
  K3 Yangian structure function $g_{\mathrm K3}(u) = \prod_{a = 1}^{24}
  (u - h_a)/(u + h_a)$ by diagonalisation of the Mukai pairing on the
  $(4, 20)$-signature lattice
  \textup{(}Theorem~\ref{rem:k3e-omega-free-field-link}\textup{).}
  Setting $u = 0$ in the coproduct $\Delta_u$ gives the cocommutative
  limit; setting $z \to w$ in the OPE
  $T(z)\,T(w) \sim c/2\,(z - w)^{-4}$ gives the conformal poles of the
  Virasoro current. These are distinct parameters: the spectral $u$ of
  the Yangian lives on the transverse $\C^2$ of the M5 uplift, the OPE
  separation $z - w$ of the sigma model lives on the worldsheet K3
  itself, and the $(q, t)$ Ding--Iohara deformation is related to
  $(\epsilon_1, \epsilon_2)$ by
  $(q, t) = (e^{\hbar\epsilon_1}, e^{-\hbar\epsilon_2})$ of
  Proposition~\ref{prop:qt-to-omega} of Chapter~\ref{ch:k3-yangian}.

\item \textbf{Modular properties.}
  $\phi_{0,1}$: weight~$0$, index~$1$, on $\SL_2(\mathbb{Z}) \ltimes
  \mathbb{Z}^2$.  $1/\eta^{24}$: weight~$-12$, on $\SL_2(\mathbb{Z})$.
  $\Delta_5$: weight~$5$, on $\mathrm{O}^+(\Lambda^{3,2})$.
  The bar Euler product matches the Ramanujan $\tau$-function:
  $\tau(1) = 1$, $\tau(2) = -24$, $\tau(3) = 252$, $\tau(4) = -1472$.

\item \textbf{Kummer enhanced symmetry} (Proposition~\ref{prop:kummer-orbifold}).
  At $K3 = T^4/\mathbb{Z}_2$: $8$ untwisted bosons (rank-$4$ lattice,
  $\mathbb{Z}_2$-invariant), $16$ twisted sectors ($h = 1/2$), enhanced
  gauge algebra $\mathfrak{so}(4)^4$ at level~$1$ (rank~$8$, $c = 8$),
  $\kappa_{\mathrm{ch}} = 2$.  The Yangian structure function
  $g_{K3}(z) = \prod_{i=1}^{24}(z - h_i)/(z + h_i)$ at the Kummer
  point has $h_i = +1$ ($i = 1, \ldots, 4$), $h_i = -1/5$
  ($i = 5, \ldots, 24$), satisfying $\sum h_i = 0$, $g(0) = 1$,
  and unitarity $g(z)\,g(-z) = 1$.  Resolution:
  $8 + 32 - 16 = 24 = \rank(\Lambda_{\mathrm{Muk}})$.

\item \textbf{Mathieu $M_{24}$ moonshine} (Eguchi--Ooguri--Tachikawa 2010).
  The massive BPS multiplicities decompose as $M_{24}$ representations:
  $A_1 = 2 \times 45 = 90$, $A_2 = 2 \times 231 = 462$,
  $A_3 = 2 \times 770 = 1540$.  The factor~$2 = \kappa_{\mathrm{ch}}$.
  The Yangian does \emph{not} see $M_{24}$ directly: the bar complex
  provides the universal coefficient $\kappa_{\mathrm{ch}} \cdot
  \dim(\rho_n)$ but cannot detect which specific finite group acts.
  This is a structural limitation: the $M_{24}$ action is a symmetry
  of the \emph{spectrum}, not of the algebra.
\end{enumerate}

\begin{remark}[Consistency does not imply existence]
\label{rem:sigma-model-consistency-caveat}
All seven checks pass (the K3 Yangian is \emph{consistent} with the
physical sigma model), but consistency does not imply existence.
The checks verify that $Y(\fg_{K3})$, if it exists, produces data
compatible with the known sigma model.  CY-A$_3$ is proved in the
infinity-categorical framework (Theorem~\ref{thm:derived-framing-obstruction});
the explicit chain-level geometric construction of $Y(\fg_{K3})$
remains conjectural.
\end{remark}

%% ========================================================================
%% Tropical limit: shadow tower, cluster algebras, tropical BPS
%% ========================================================================

\section{Tropical limit and cluster algebra structure}
\label{sec:k3e-tropical}

The large complex structure limit (maximally unipotent monodromy, i.e.\ the MUM degeneration) of $K3$ produces a \emph{tropical K3}: the dual intersection complex $\Delta(K3)$ of a type~III Kulikov degeneration is a triangulated $S^2$ with $24$~marked points, one for each direction in the Mukai lattice $\Lambda_{\mathrm{Muk}} \simeq U^3 \oplus E_8(-1)^2$.  The shadow tower, scattering diagram, and Yangian parameters each acquire a tropical (piecewise-linear) incarnation in this limit, governed by the cluster algebra structure of Gross--Hacking--Keel--Kontsevich.

\subsection{The tropical K3}
\label{subsec:tropical-k3}

\begin{definition}[Dual intersection complex]
\label{def:tropical-k3}
Let $\pi \colon \cX \to \Delta$ be a type~III Kulikov degeneration of K3 surfaces.  The \emph{dual intersection complex} $\Delta(K3) = \Delta(\pi)$ is the simplicial complex whose vertices are the irreducible components of the central fibre $X_0$, with a $k$-simplex for each $(k+1)$-fold intersection.  For a maximally degenerate K3:
\begin{enumerate}[label=\textup{(\roman*)}]
\item $\Delta(K3)$ is homeomorphic to $S^2$ (by Friedman--Morrison).
\item $|\Delta(K3)_0| = 24$ (the number of vertices equals $\rank(\Lambda_{\mathrm{Muk}})$).
\item The $1$-skeleton carries an integral affine structure on $S^2 \setminus \{24\text{ singular points}\}$, with focus-focus singularities at each vertex.
\item The edge weights are determined by the Mukai pairing $\omega^{ij}$ restricted to adjacent components.
\end{enumerate}
\end{definition}

\begin{remark}[The 24 singular points]
\label{rem:tropical-k3-24}
The $24$ focus-focus singularities on the tropical K3 are the tropical analogue of the $24$ nodal fibres in an elliptic K3 (equivalently, the $24$ roots of the discriminant locus).  The monodromy around each singularity is $\left(\begin{smallmatrix} 1 & 1 \\ 0 & 1 \end{smallmatrix}\right)$, and the total monodromy around $S^2$ is trivial (Gauss--Bonnet).  The number $24 = \chi_{\mathrm{top}}(K3) = \kappa_{\mathrm{fiber}}$ appears as a topological invariant of the degeneration.
\end{remark}

\subsection{Tropical shadow tower}
\label{subsec:tropical-shadow}

The shadow obstruction tower $\Theta_A = \kappa_{\mathrm{ch}} + C + Q + \cdots$ tropicalizes to a piecewise-linear (PL) function on $\Delta(K3)$.

\begin{conjecture}[Tropical shadow tower]
\label{conj:tropical-shadow}
\ClaimStatusConjectured
Let $A_{K3}$ be the chiral algebra of a K3 surface via CY-A$_2$ \textup{(}proved\textup{)}, and let $\Theta_A \in \mathrm{MC}(\mathrm{Def}_{\mathrm{cyc}}^{\mathrm{mod}}(A_{K3}))$ be its shadow MC element.  In the type~III Kulikov degeneration:
\begin{enumerate}[label=\textup{(\roman*)}]
\item The shadow tower $\Theta_A$ degenerates to a PL function $\Theta^{\mathrm{trop}} \colon \Delta(K3) \to \R$, affine on each maximal cell.
\item The degree-$2$ component satisfies $\kappa_{\mathrm{ch}}^{\mathrm{trop}} = 2$ \textup{(}constant\textup{)}, reflecting the topological invariance of $\kappa_{\mathrm{ch}} = \chi(\cO_{K3}) = 2$.
\item The bending locus of $\Theta^{\mathrm{trop}}$ \textup{(}the set of edges where the slope jumps\textup{)} is a subset of the walls in the Gross--Siebert scattering diagram on $\Delta(K3)$.
\item The tropical MC equation $(\partial^{\mathrm{trop}})^2 \Theta^{\mathrm{trop}} = 0$ holds, where $\partial^{\mathrm{trop}}$ is the tropical differential \textup{(}slope comparison across edges\textup{)}.
\end{enumerate}
This is conditional on the identification of the shadow MC element with the KS MC element through the motivic Hall algebra \textup{(}Conjecture~\textup{\ref{conj:k3e-two-mc})}.
\end{conjecture}

At the generic point of K3 moduli \textup{(}shadow class~$G$, free-field Heisenberg\textup{)}: $C^{\mathrm{trop}} = Q^{\mathrm{trop}} = 0$, and the tropical shadow tower reduces to the constant $\kappa_{\mathrm{ch}}^{\mathrm{trop}} = 2$.  At enhanced symmetry points \textup{(}shadow class $\geq L$, cf.\ Section~\ref{sec:k3-sigma-model-cross-check}\textup{)}: the cubic shadow $C^{\mathrm{trop}}$ localises at the enhancement vertices.

\subsection{Cluster algebra structure on the scattering diagram}
\label{subsec:tropical-cluster}

The Gross--Hacking--Keel--Kontsevich (GHKK) programme endows the scattering diagram $\mathfrak{D}_{K3 \times E}$ with a cluster algebra structure.

\begin{conjecture}[Cluster--shadow correspondence]
\label{conj:cluster-shadow}
\ClaimStatusConjectured
The scattering diagram $\mathfrak{D}_{K3 \times E}$ \textup{(}Definition~\textup{\ref{def:k3e-scattering})} carries a cluster algebra structure with:
\begin{enumerate}[label=\textup{(\roman*)}]
\item \emph{Exchange matrix}: $B_{ij} = \chi(S_i, S_j) - \chi(S_j, S_i)$, the antisymmetric part of the Euler form on $K_0(D^b(\Coh(K3)))$.
\item \emph{Cluster mutations}: wall-crossing at a wall of marginal stability $W(\gamma_1, \gamma_2)$ corresponds to the Fomin--Zelevinsky mutation $\mu_k$ at the vertex $k$ dual to the decaying BPS state.
\item \emph{Cluster variables}: the canonical theta functions $\vartheta_\gamma$ on the mirror dual algebraic torus are the cluster variables, with the structure constants computed by broken-line counting \textup{(}GPS tropical vertex\textup{)}.
\item \emph{Shadow compatibility}: the tropical shadow $\Theta^{\mathrm{trop},\leq r}$ at degree $r$ is a PL function whose bending locus lies within the walls of the scattering diagram at heights $\leq r$.
\end{enumerate}
CY-A$_3$ is proved in the infinity-categorical framework; the explicit chain-level construction for $K3 \times E$ remains conditional.  The tropical vertex computation is unconditional \textup{(}GPS is a theorem\textup{)}.
\end{conjecture}

\begin{remark}[GPS tropical vertex and the pentagon identity]
\label{rem:gps-pentagon}
The Gross--Pandharipande--Siebert tropical vertex algorithm computes the consistent scattering diagram iteratively.  In the simplest case \textup{(}two seed walls at directions $(1,0)$ and $(0,1)$ with multiplicity~$1$\textup{)}, the consistency condition forces a wall at $(1,1)$ with multiplicity~$1$: this is the tropical incarnation of the pentagon identity \textup{(}Theorem~\textup{\ref{thm:scattering-mc}(ii)}, $A_2$ case\textup{)}.  The broken-line count at charge $(a,b)$ equals the wall multiplicity, which is the tropical Hurwitz number $N_{(a,b)}^{\mathrm{trop}}$.  Verification: $100$ tests in \texttt{tropical\_shadow\_cluster.py}, including GPS consistency through height~$8$ and broken-line/wall-multiplicity cross-checks.
\end{remark}

\subsection{Tropical BPS states and broken lines}
\label{subsec:tropical-bps}

The BPS states of $K3 \times E$ tropicalize to \emph{broken lines} on the scattering diagram: PL paths in the charge lattice that bend at walls, accumulating monomial weights at each crossing.

\begin{conjecture}[Tropical BPS correspondence]
\label{conj:tropical-bps}
\ClaimStatusConjectured
At the MUM point \textup{(}large complex structure degeneration\textup{)} of K3, the tropical BPS invariant
\[
  \Omega^{\mathrm{trop}}(\gamma) \;=\; \lim_{t \to 0} \Omega(\gamma; \sigma_t)
\]
\textup{(}where $\sigma_t$ is a family of Bridgeland stability conditions degenerating to the tropical point\textup{)} equals:
\begin{enumerate}[label=\textup{(\roman*)}]
\item The wall multiplicity at direction $\gamma$ in the consistent scattering diagram \textup{(}by GPS\textup{)}.
\item The broken-line count $b(\gamma) = \#\{$broken lines of total charge $\gamma\}$.
\item The coefficient of $x^\gamma$ in the theta function $\vartheta_\gamma$ \textup{(}cluster variable\textup{)}.
\end{enumerate}
CY-A$_3$ is proved (inf-cat); the explicit chain-level $K3 \times E$ identification remains conditional.
\end{conjecture}

\subsection{Cluster structure on the K3 Yangian}
\label{subsec:tropical-yangian-cluster}

\begin{conjecture}[Cluster Yangian]
\label{conj:cluster-yangian}
\ClaimStatusConjectured
The K3 Yangian $Y(\fg_{K3})$ \textup{(}CY-A$_3$ proved inf-cat; explicit geometric construction conjectural\textup{)} carries a cluster algebra structure from the Maulik--Okounkov construction:
\begin{enumerate}[label=\textup{(\roman*)}]
\item The Yangian parameters $h_1, \ldots, h_{24}$ transform under cluster mutations as
\[
  \mu_k(h_i) = \begin{cases} h_i & i \neq k, \\ -h_k + \displaystyle\sum_{j \colon B_{kj} > 0} B_{kj}\,h_j & i = k, \end{cases}
\]
where $B_{ij}$ is the exchange matrix.
\item The structure function $g_{K3}(z) = \prod_{i=1}^{24}(z - h_i)/(z + h_i)$ satisfies $g(z)\,g(-z) = 1$ \textup{(}algebraic unitarity\textup{)} for \emph{all} parameter choices, including all mutations.
\item Different stability conditions in $\Stab^\dagger(K3)$ correspond to different cluster seeds, related by mutation sequences.  The exchange graph is the Bridgeland--Smith stability graph.
\item In the tropical limit, the $R$-matrix becomes piecewise-constant: $R^{\mathrm{trop}}_{ii}(z) = (z - h_i)/(z + h_i)$ on each chamber, with wall-crossing given by the mutation rule~\textup{(i)}.
\end{enumerate}
This is doubly conjectural: conditional on both the explicit geometric construction of $Y(\fg_{K3})$ \textup{(}CY-A$_3$ proved inf-cat; chain-level construction open\textup{)} and its cluster structure \textup{(}open\textup{)}.
\end{conjecture}

\begin{remark}[Unitarity is unconditional]
\label{rem:tropical-unitarity}
The identity $g_{K3}(z) \cdot g_{K3}(-z) = 1$ is a purely algebraic consequence of the product structure and holds for \emph{any} parameters $h_i$ (not just those satisfying the CY$_2$ constraint $\sum h_i = 0$).  This is verified computationally for $3$-, $24$-, and $2$-parameter systems in \texttt{tropical\_shadow\_cluster.py}.
\end{remark}

\noindent\textit{Verification}: $100$~tests in \texttt{tropical\_shadow\_cluster.py} covering the tropical K3 ($24$~vertices, PL functions, tropical Laplacian), shadow tower (arities~$2$--$4$, class G/L, MC equation), GPS tropical vertex (pentagon identity, consistency through height~$8$, broken-line counts), cluster algebra ($A_2$/$A_3$/local~$\bP^2$ exchange matrices, mutation involution, antisymmetry preservation), Yangian (unitarity, parameter mutation, pole detection), and $14$~multi-path cross-verification tests (GPS vs broken-line, mutation involution via two paths, wall-count monotonicity, Laplacian conservation law).


%% ===================================================================
%% BFN Coulomb branch route to the K3 Yangian
%% ===================================================================

\section{BFN quantised Coulomb branch route to the K3 Yangian}
\label{sec:bfn-coulomb-k3}

The Braverman--Finkelberg--Nakajima (BFN) programme constructs
\emph{quantised Coulomb branch algebras} $A_\hbar(G, N)$ from $3d$
$\cN = 4$ gauge theories $(G, N)$ via convolution on the affine
Grassmannian $\mathrm{Gr}_G = G(\!(t)\!)/G[\![t]\!]$
(arXiv:1601.03586, 1604.03625).  This route is \emph{independent} of
the CY-to-chiral correspondence programme $\{\Phi_d\}$ and of CY-A$_3$, providing a third
construction of the K3 Yangian alongside the CY-to-chiral route~(a) and
the CoHA route~(b).

\begin{conjecture}[BFN identification for K3, Mukai-lattice form]
\label{conj:bfn-k3-yangian-mukai}
\ClaimStatusConjectured\quad
\index{BFN Coulomb branch!K3 Yangian}%
Let $(G, N)$ be the gauge theory data of the Vafa--Witten twist
on $K3 \times \R^3$.  The BFN quantised Coulomb branch algebra satisfies
\[
  A_\hbar(K3) \;\cong\; Y(\fg_{K3}),
\]
where $Y(\fg_{K3})$ is the K3 Yangian of rank~$24$ (Mukai lattice rank).

The evidence for this identification is:
\begin{enumerate}[label=\textup{(\roman*)}]
  \item Both algebras have rank~$24 = \mathrm{rk}(\Lambda_{\mathrm{Muk}})$.
  \item Both have the same classical limit: the Mukai Heisenberg
    $H_{\mathrm{Muk}}$ with pairing of signature~$(4, 20)$.
  \item The BFN $R$-matrix from Maulik--Okounkov stable envelopes
    matches the K3 structure function
    $g_{K3}(z) = \prod_{i=1}^{24} (z - h_i)/(z + h_i)$.
  \item The BFN Fock space $K_T(\mathrm{Hilb}^n(K3 \times E))$ has
    character $\sum_n \chi(\mathrm{Hilb}^n(K3))\, q^n
    = \prod_{k \geq 1}(1 - q^k)^{-24} = 1/\Delta(q)$.
  \item At the $A_1$ orbifold point $K3 = T^4/\Z_2$
    \textup{(}resolved\textup{)}, BFN gives
    $Y(\widehat{\fsl}_2)$ \textup{(}PROVED, BFN 2016\textup{)},
    which embeds as the enhanced-symmetry subalgebra.
\end{enumerate}

The BFN route reorganises the K3 Yangian construction into subproblems
\textup{(}choosing $(G, N)$, identifying the matter representation,
verifying the quantisation\textup{)}.  Each subproblem requires
independent verification; the reorganisation does not trivially solve
any of them.
\end{conjecture}

\begin{remark}[BFN for ADE quivers: proved cases]
\label{rem:bfn-ade-proved}
For ADE quiver gauge theories, the BFN identification is a
\emph{theorem}: this is Theorem~\ref{thm:bfn-phi-ade-identification}
(\texttt{k3\_yangian\_chapter.tex}), assembling Braverman--Finkelberg--Nakajima
(arXiv:1604.03625) with Nakajima--Takayama (arXiv:1606.02002) and Kronheimer
(J.~Diff.~Geom.\ 29, 1989), and giving
$\Phi(T^*\widetilde{S}_\fg) \simeq \cA_\hbar(Q_\fg, \mathbf{v}, \mathbf{w})
\simeq Y^\mu(\widehat{\fg})_{k=1}$ for $\fg$ of ADE type.
At the Kummer point $K3 = T^4/\Z_2$: the McKay quiver is the
$\widehat{A}_1$ affine Dynkin diagram with $2$~vertices.  The
BFN Coulomb branch at gauge ranks $\mathbf{v} = (1, 1)$ and
framing $\mathbf{w} = (1, 0)$ produces $Y(\widehat{\fsl}_2)$
(proved).  Passing to the full K3 (away from the orbifold point)
requires the non-quiver BFN construction; this is the conjectural
step captured by Conjecture~\ref{conj:bfn-k3-yangian-mukai}.

The Hilbert scheme Euler characteristics $\chi(\mathrm{Hilb}^n(K3))$
provide a strong consistency check: $\chi_0 = 1$, $\chi_1 = 24$,
$\chi_2 = 324$, $\chi_3 = 3200$, $\chi_4 = 25650$, $\chi_5 = 176256$
(G\"ottsche 1990), matching $\prod(1 - q^k)^{-24}$ exactly (verified
computationally through $n = 10$).

\noindent\textit{Verification}: $93$ tests in
\texttt{bfn\_coulomb\_k3\_yangian.py} covering ADE quiver data for
all types through $E_8$, Jordan quiver BFN (proved template),
Hilbert scheme Euler characteristics through $n = 10$,
Kummer orbifold specialisation, stable envelope $R$-matrix
matching, and cross-engine consistency with \texttt{k3\_yangian.py}.
\end{remark}

\section{Root-of-unity S-matrix: abelian rigidity and the \texorpdfstring{$N \geq 3$}{N >= 3} threshold}
\label{sec:root-unity-smatrix}

At level $N = 2$, the $324 = p_{24}(2)$-module MTC
(\texttt{root\_of\_unity\_n2\_explicit.py}, $59$~tests)
has a block-diagonal S-matrix: $24$ Hadamard blocks of size~$2 \times 2$
(Ising sectors) and $276$ trivial $1 \times 1$ blocks.
All quantum dimensions equal~$1$ (pointed category), and the braiding
is completely determined by the charge lattice $(\bZ/2\bZ)^{24}$.

\begin{conjecture}[Non-abelian S-matrix rigidity at \texorpdfstring{$N = 2$}{N = 2}]
\label{conj:n2-smatrix-rigidity}
\ClaimStatusConjectured
At an $A_1$ enhancement point in K3 moduli space, the MTC S-matrix
for the $N = 2$ truncation of $U_{q,t}(\fg_{K3}^{\mathrm{tor}})$ is
\emph{unchanged}:
\[
  S^{A_1}_{\lambda\mu} \;=\; S^{\mathrm{ab}}_{\lambda\mu}
  \qquad\text{for all } \lambda,\mu \in \{1,\ldots,324\}.
\]
The non-abelian correction vanishes because the double braiding
at $q = e^{2\pi i/2} = -1$ is trivial: the Drinfeld--Jimbo
R-check for $U_q(\fsl_2)$ has eigenvalues $q = -1$ (symmetric,
multiplicity~$3$) and $-q^{-1} = +1$ (antisymmetric, multiplicity~$1$),
so $q^2 = 1$ and $(-q^{-1})^2 = 1$.
\end{conjecture}

\begin{remark}[The \texorpdfstring{$N = 2$}{N = 2} rigidity mechanism]
\label{rem:n2-rigidity-mechanism}
The Yang R-matrix $R(u) = (u\,\mathrm{Id} + \hbar\,P)/(u + \hbar)$
introduces $48$ off-diagonal entries in the spectral R-matrix at
the $A_1$ point (Remark~\ref{rem:charge2-offdiagonal}).
These modify the $u$-dependent braiding but leave the \emph{modular}
data invariant: the MTC S-matrix is the double braiding trace, and
$(q^2 = 1,\; q^{-2} = 1)$ renders all corrections trivial.

At $u = 0$, the Yang R-matrix equals the permutation $P$, and
$P^2 = \mathrm{Id}$ gives trivial monodromy directly.
The pointed structure (all $d_\lambda = 1$) forces the MTC to be
a \emph{metric group} (Deligne 2002): its modular data is
determined uniquely by the abelian group structure of the fusion ring,
which the $A_1$ enhancement preserves.
\end{remark}

\begin{conjecture}[Non-abelian threshold at \texorpdfstring{$N = 3$}{N = 3}]
\label{conj:n3-nonabelian-threshold}
\ClaimStatusConjectured
At $N = 3$ ($q = e^{2\pi i/3}$), the $A_1$ enhancement produces
\emph{genuine} off-diagonal corrections to the MTC S-matrix.
The double braiding eigenvalue $q^2 = e^{4\pi i/3} \neq 1$ is
nontrivial, and the Verlinde formula yields non-abelian fusion rules
among the $p_{24}(3) = 3200$ simple modules.
The quantum integer $[2]_q = 2\cos(2\pi/3) = -1 \neq 0$
permits non-trivial representations; $[3]_q = 0$ provides the
truncation.
\end{conjecture}

\noindent\textit{Verification}: $59$ tests in
\texttt{root\_of\_unity\_smatrix.py} covering the Drinfeld--Jimbo
R-check eigenvalues at $q = -1$, double braiding triviality,
S-matrix identity $S^{A_1} = S^{\mathrm{ab}}$, module
classification at the $A_1$ point ($5 + 44 + 44 + 231 = 324$),
Verlinde fusion ($\bZ/2\bZ$ per Ising sector, unchanged),
the $N \geq 3$ threshold analysis, and Yang R-matrix unitarity.

\subsection{The \texorpdfstring{$N = 3$}{N = 3} MTC: first non-abelian level}
\label{subsec:n3-mtc}

At $N = 3$ the threshold is crossed.  The $p_{24}(3) = 3200$
simple modules decompose into six types:
\begin{equation}\label{eq:n3-module-types}
  \underbrace{24}_{(3_a)} \;+\;
  \underbrace{24}_{(2_a,1_a)} \;+\;
  \underbrace{24}_{(1_a^3)} \;+\;
  \underbrace{552}_{(2_a,1_b)} \;+\;
  \underbrace{552}_{(1_a^2,1_b)} \;+\;
  \underbrace{2024}_{(1_a,1_b,1_c)}
  \;=\; 3200.
\end{equation}
These organise into $2600$ sectors: $24$ \emph{Ising sectors} of
size~$3 \times 3$ (charge concentrated in one Mukai direction),
$552$ \emph{Hadamard sectors} of size~$2 \times 2$ (charge split
$(2,1)$ across two directions), and $2024$ trivial $1 \times 1$
sectors (charge spread across three directions).

\begin{conjecture}[Non-abelian quantum dimensions at \texorpdfstring{$N = 3$}{N = 3}]
\label{conj:n3-nonabelian-qdims}
\ClaimStatusConjectured
Within each Ising sector, the three modules carry quantum dimensions
\[
  d_{V_0} = 1, \qquad d_{V_1} = \sqrt{2}, \qquad d_{V_2} = 1,
\]
where the $\mathfrak{sl}_2$ level-$2$ Verlinde S-matrix provides
the intra-sector braiding:
\begin{equation}\label{eq:ising-smatrix}
  S^{\mathrm{Ising}} \;=\;
  \begin{pmatrix}
    \frac{1}{2} & \frac{1}{\sqrt{2}} & \frac{1}{2} \\[4pt]
    \frac{1}{\sqrt{2}} & 0 & -\frac{1}{\sqrt{2}} \\[4pt]
    \frac{1}{2} & -\frac{1}{\sqrt{2}} & \frac{1}{2}
  \end{pmatrix}.
\end{equation}
The $24$ Type~$B_1$ modules $(2_a, 1_a)$ have quantum dimension
$\sqrt{2}$; all other $3176$ modules have quantum dimension~$1$.
The total quantum dimension squared is
$D^2 = 24 \cdot 4 + 552 \cdot 2 + 2024 \cdot 1 = 3224 > 3200$,
confirming the MTC is \emph{not pointed}.
\end{conjecture}

The Ising fusion rules within each $3 \times 3$ sector are
non-abelian:
\begin{equation}\label{eq:ising-fusion}
  V_1 \otimes V_1 \;=\; V_0 \oplus V_2,
  \qquad
  V_1 \otimes V_2 \;=\; V_1,
  \qquad
  V_2 \otimes V_2 \;=\; V_0.
\end{equation}
The fusion $V_1 \otimes V_1 = V_0 \oplus V_2$ (two summands)
is the hallmark of non-abelian braiding; it is impossible in a
pointed MTC where all fusions are single-valued.

\begin{conjecture}[Charge-1 sector at \texorpdfstring{$A_1$}{A-1} enhancement]
\label{conj:n3-charge1-a1}
\ClaimStatusConjectured
The $24 \times 24$ sub-matrix of the abelian S-matrix restricted
to the $24$ Type~$C_1$ modules $(1_a^3)$ is diagonal with
$S_{aa} = \tfrac{1}{2}$.
At an $A_1$ enhancement merging directions $(a,b)$, the double
braiding phase $q^{-2} = e^{-4\pi i/3} \neq 1$ produces a
genuine off-diagonal entry
\[
  S^{\mathrm{enh}}_{ab}
  \;=\; \frac{q^{-2}}{\sqrt{2}}
  \;=\; \frac{e^{-4\pi i/3}}{\sqrt{2}},
\]
coupling the previously decoupled modules $(1_a^3)$ and $(1_b^3)$.
The enhanced $24 \times 24$ matrix retains full rank~$24$.
\end{conjecture}

\begin{remark}[The \texorpdfstring{$N = 2 \to N = 3$}{N = 2 -> N = 3} phase transition]
\label{rem:n2-n3-phase-transition}
The passage from $N = 2$ to $N = 3$ is a \emph{phase transition}
in the MTC structure:
\begin{itemize}
\item Pointed $\to$ non-pointed
  ($d_{\max} = 1 \to d_{\max} = \sqrt{2}$).
\item Trivial double braiding $\to$ nontrivial
  ($q^2 = 1 \to q^2 = e^{4\pi i/3}$).
\item $\bZ/2\bZ$ fusion $\to$ Ising fusion
  ($V_1 \otimes V_1 = V_0 \oplus V_2$).
\item Rigid S-matrix $\to$ enhancement-sensitive S-matrix.
\item $324$ modules in $300$ sectors $\to$
  $3200$ modules in $2600$ sectors.
\end{itemize}
The quantum integer $[2]_q = -1 \neq 0$ ensures the fundamental
representation survives the truncation, while $[3]_q = 0$
provides the level-$2$ cutoff.  The Ising MTC
($\mathrm{SU}(2)_2$) is the simplest non-abelian MTC,
and it appears here as the \emph{universal local factor} in
each of the $24$ Mukai-direction sectors.
\end{remark}

\noindent\textit{Verification}: $80$ tests in
\texttt{root\_of\_unity\_n3\_mtc.py} covering module enumeration
($3200 = 24 + 24 + 24 + 552 + 552 + 2024$, six types),
sector decomposition ($24 + 552 + 2024 = 2600$ sectors),
Ising S-matrix properties (symmetric, unitary, $S^2 = I$,
$\det = -1$, quantum dimensions $1, \sqrt{2}, 1$),
Verlinde fusion rules (non-abelian: $V_1 \otimes V_1 = V_0 \oplus V_2$,
non-negative integer coefficients, Verlinde consistency
$\sum_k N^k_{ij} d_k = d_i d_j$),
DJ R-matrix eigenvalues at $q = e^{2\pi i/3}$
(nontrivial double braiding $q^2 \neq 1$),
full $3200 \times 3200$ S-matrix (symmetric, unitary, full rank,
block-diagonal with Ising and Hadamard blocks),
charge-$1$ sector ($24 \times 24$ diagonal at abelian level,
off-diagonal at $A_1$ enhancement, full rank),
and $N = 2$ vs.\ $N = 3$ comparison
(pointed $\to$ non-pointed, $D^2 = 324 \to 3224$).


\section{Quantum-toroidal Pentagon-at-\texorpdfstring{$E_1$}{E-1}: edge-architecture closure}
\label{sec:k3qt-pentagon-edge-architecture}

The K$3$ quantum-toroidal algebra $U^+_{q, t}(\widehat{\widehat{\fgl}}_1)^{K3}$
inherits the edge-architecture closure of the K$3$ Pentagon-at-$E_1$
\textup{(}Theorem~\ref{thm:k3-pentagon-E1-edge-architecture} of
Chapter~\ref{ch:k3-yangian}\textup{).}  This section records the
status at cohomological level and the chain-level routing through the
bigraded edge-character matrix.

\subsection{Cohomological status}
\label{subsec:k3qt-cohomological-status}

\begin{theorem}[K\texorpdfstring{$3$}{3} quantum-toroidal Pentagon-at-\texorpdfstring{$E_1$}{E-1};
                cohomological, conditional on Vol II fibre cells]
\label{thm:k3qt-pentagon-cohomological}
\ClaimStatusConditional
Conditional on closure of the K$3$ Heisenberg and affine-Yangian-fibre
cells in Vol II, the quantum vertex chiral group
$U^+_{q,t}(\widehat{\widehat{\fgl}}_1)^{K3}$ satisfies Pentagon-at-$E_1$
at \emph{cohomological} level.  The cocycle $[\omega]_{K3}$ vanishes in
$H^2(\mathrm{SC}^{\mathrm{ch,top}};\, \mathfrak{aut})$ via the
edge-character matrix $M$ of
Theorem~\ref{thm:k3-pentagon-E1-edge-architecture},
with diagonal entries $(0, 5, -16, 11)$ at $K3 \times E$.
\end{theorem}

\begin{proof}[Proof sketch]
The K$3$ Mukai cell decomposes into two fragments:
(a) ADE simply-laced (covered by the Vol II affine-Yangian fibre),
(b) abelian Heisenberg $\widehat{\fu}(1)^{24}$ (covered by the
$\fg \to 0$ limit, with the closed-form Heisenberg R-matrix
verified independently).  The edge-character matrix is diagonal at
$K3 \times E$ with entries
$(\Pi_{++}, \Pi_{+-}, \Pi_{-+}, \Pi_{--}) = (0, 5, -16, 11)$.  The
trace $\sum \mathrm{tr}_{\Pi_{\epsilon\epsilon'}}(\mathfrak{K}_{\mathcal{C}})
= 0 = \chi(\mathcal{O}_{K3 \times E})$ closes by Caldararu chiral
HRR + Hattori--Stallings projection
\textup{(}Theorem~\ref{thm:k3-multiproj-bigraded-lefschetz} of
Chapter~\ref{ch:k3-yangian}\textup{).}
\end{proof}

\subsection{Chain-level status via the Stasheff bridge lemma}
\label{subsec:k3qt-chain-level-status}

\begin{theorem}[K\texorpdfstring{$3$}{3} quantum-toroidal Pentagon-at-\texorpdfstring{$E_1$}{E-1} chain-level;
                conditional on the Vol~II Pentagon-coherence bridge lemma]
\label{thm:k3qt-pentagon-chain-level}
\ClaimStatusConditional
The cohomological closure of
Theorem~\ref{thm:k3qt-pentagon-cohomological} lifts to a chain-level
identity $\delta_A = 0$ in $\mathrm{Hom}_{\mathrm{Ch}}$ conditional
on the following bridge lemma (expected from Vol~II):

\medskip
\textbf{Pentagon-coherence chain-level lift.}
For each specialisation cell
\[
C \in \{\mathrm{Yangian},\;
\mathrm{Aff.Yangian},\; \mathrm{Shifted},\; \mathrm{Trunc/BFN},\;
\mathrm{Finite}\, W,\; \mathrm{Aff}.W\,\mathrm{principal},\;
\mathrm{Aff}.W\,\mathrm{non}\text{-}\mathrm{principal},\; \mathrm{BP},\;
\mathrm{Orth.coideal}\},
\]
the cohomological Pentagon coherence on
$C$ admits a chain-level lift via Stasheff $K_5$-associahedral chain
witnesses $h_e$, conditional on a detecting-family hypothesis and
Stasheff $K_5$-cocycle vanishing.
\medskip

For the K$3$ cell (ADE simply-laced + abelian Heisenberg fragments),
the bridge lemma holds as a corollary of:
(i) explicit chain-level Pentagon vanishing for the abelian
Heisenberg fragment via Schur centrality
\textup{(}Theorem~\ref{thm:heisenberg-pentagon-E1} of
Chapter~\ref{ch:e1-chiral}\textup{);}
(ii) Etingof--Kazhdan quantization for the ADE fragment with explicit
Stasheff $K_n$ witnesses
\textup{(}Drinfeld 1989; Markl--Shnider--Stasheff 2002\textup{);}
(iii) Stasheff $K_5$-cocycle vanishing
$[\eta^{(3)}]_{K3} = 0 \in H^3$, verified for K$3$ via the joint
detecting-family of the three K$3$ Yangian routes
(Borcherds singular theta, Gritsenko--Nikulin BKM, Kawai--Yoshioka).
\end{theorem}

\begin{remark}[Per-class status table]
\label{rem:k3qt-per-class-status}
The landscape dichotomy applied to quantum-toroidal Pentagon-at-$E_1$:

\begin{center}
\small
\begin{tabular}{p{0.28\textwidth}p{0.22\textwidth}p{0.42\textwidth}}
\toprule
Class & Pentagon status (cohomological) & Pentagon status (chain-level) \\
\midrule
A (K$3$, K$3 \times E$) & UNCONDITIONAL & cond.\ bridge lemma, K$3$ cell \\
A' (K$3$ orbifolds) & UNCONDITIONAL & cond.\ bridge lemma + fibre-localisation \\
A'' (Conway, Leech) & UNCONDITIONAL (algebraic) & cond.\ bridge lemma, algebraic cell \\
B$_0$ (conifold) & UNCONDITIONAL & cond.\ conifold-specific bridge \\
Algebraic Yang--Baxter sub-class ($Y(\fsl_n)$) & UNCONDITIONAL & cond.\ resurgent Drinfeld twist, $(2, 2)$-component \\
Mock-modular sub-class (quintic, LP$^2$) & CONJECTURAL & cond.\ universal mock-modular receptacle $\Rightarrow$ resurgent Drinfeld twist, $(2, \infty)$-component \\
\bottomrule
\end{tabular}
\end{center}
\end{remark}

\subsection{Discipline summary}
\label{subsec:k3qt-discipline}

\begin{remark}[Cohomological vs chain-level discipline]
\label{rem:k3qt-cohom-vs-chain-discipline}
Every Pentagon-at-$E_1$ status claim discloses four elements:
(i) volume dependency, (ii) cohomological vs chain-level modality,
(iii) specialisation cell (among the eight fibres enumerated in
Vol~II), (iv) per-input verification.  The status promotions in this
chapter
\textup{(}Theorems~\ref{thm:k3qt-pentagon-cohomological}--\ref{thm:k3qt-pentagon-chain-level}\textup{)}
satisfy all four elements.  Cohomological closure at K$3$ is
unconditional; chain-level closure is conditional on the bridge
lemma, which has explicit reduction targets for each dichotomy class.
\end{remark}

\begin{remark}[The \texorpdfstring{$V_4$}{V-4}-equivariant Lefschetz interpretation]
\label{rem:k3qt-V4-equivariant}
In the unified form of the edge-character matrix, the K$3$
quantum-toroidal Pentagon closure is the $V_4$-equivariant Lefschetz
fixed-point identity
$\sum_{\epsilon\epsilon'} \mathrm{tr}_{\Pi_{\epsilon\epsilon'}}(\mathfrak{K})
= \chi(\mathcal{O}_X)$ on the chiral Hochschild
$\mathrm{ChirHoch}^\bullet(A^{K3}_{q,t}, A^{K3}_{q,t})$, with the
faithful $V_4$-action localised on $H^{1,1}_{\mathrm{prim}}$
(rank $20$).  The matrix $M$ is integer-valued, diagonal at K$3$,
with kernel $\Pi_{++}$ at $K3 \times E$ (entry $0$).  The vanishing
trace is the Pentagon-coherence content; the $-16$ Berezinian entry
is rigidly forced by the K$3$ Mukai signature $(4, 20)$ via the
universal $(p+q)^2 = (p-q)^2 + 4pq$ identity.
\end{remark}

\medskip

What the chapter establishes.  The rational K$3$ Yangian
$Y(\fg_{K3})$ sits at the Yangian vertex of a four-stage
Drinfeld--Jimbo tower whose quantum-toroidal apex
$U_{q,t}(\fg_{K3}^{\mathrm{tor}})$ loses the $S_3$ Miki symmetry of the
$\bC^3$ case and retains only the $\SL_2(\bZ)$ mapping class of the
elliptic fibre.  Five independent bridges --- AGT identification with
the Vafa--Witten partition function, the twisted M-theory
compactification web, the tropical-cluster incarnation at the MUM
point, the BFN Coulomb-branch realisation, and the root-of-unity
$N = 2 \to N = 3$ phase transition --- cross-check the same structure
function $g_{K3}(u) = \prod_{a=1}^{24}(u - h_a)/(u + h_a)$ and its
trigonometric lift.  The Pentagon-at-$E_1$ edge-architecture closes
cohomologically at K$3$ and K$3 \times E$; the chain-level lift reduces
to a single Stasheff bridge lemma.  Each conjectural bridge carries an
explicit dependency on CY-A$_3$ beyond its inf-cat resolution; the
classical limits ($q \to 1$, $h_i \to 0$, MUM) recover proved
structures.

\begin{remark}[24-Miki presentation with \texorpdfstring{$\widetilde{M}_{24}$}{widetildeM-24}-umbral cocycle]
\label{rem:k3qt-24miki-M24-umbral}
The K3-quantum-toroidal object admits a $24$-fold presentation on the nodes of the
$24$-node discriminant curve $E^{\mathrm{nod}}_{24}$: $24$ copies
$\{e^{(i)}_r,\,f^{(i)}_r,\,\psi^{(i)\pm}_s\}$ of Miki quantum toroidal
$U_{q_1,q_2}(\widehat{\widehat{\mathfrak{gl}}}_1)$ (Miki 2007;
Feigin--Hashizume--Hoshino--Shiraishi--Yanagida 2010), one per node, with
$\widetilde{M}_{24}$-twisted identification between copies
\[
  \sigma\cdot e^{(i)}_r
  \;=\;
  \exp\!\bigl(2\pi i\, r\, m_\sigma/24\bigr)\, e^{(\sigma(i))}_r,
  \qquad \sigma\in M_{24},\;m_\sigma\in\{0,\dots,23\},
\]
where $m_\sigma$ is the umbral shift of Cheng--Duncan--Harvey 2014. The fusion
kernel across neighbouring nodes near a Humbert wall is
\[
  F_{ij}(z,w)
  \;=\;
  \frac{(1-q_1 z/w)(1-q_2 z/w)}{1-q_3 z/w},\qquad q_1 q_2 q_3 = 1,
\]
via the Feigin--Jing--Miki rank-embedding into
$U_{q_1,q_2}(\widehat{\widehat{\mathfrak{gl}}}_{24})$. See
Chapter~\ref{ch:k3-chiral-bialgebra-platonic},
Theorem~\ref{thm:k3-24miki-presentation}.
\end{remark}

\section{K3 quantum toroidal algebra as BKM shadow}
\label{sec:k3qt-bkm-shadow}

\begin{remark}[K3 quantum toroidal algebra as BKM shadow]
\label{rem:k3qt-bkm-shadow}
The K3 quantum toroidal algebra $U^{\mathrm{tor}}_{q_1,q_2,q_3}(\widehat{\widehat{\mathfrak{sl}}_N})$
developed here, with Miki triality generators
$\{e^{(i)}_r,\, h^{(i)}_r,\, f^{(i)}_r\}$ and Drinfeld coproduct, becomes
the non-abelian chiral bialgebra $\mathbf H_{\Delta_5}$ after
the following specialisation:
\begin{itemize}
\item Parameters $(q_1, q_2, q_3)$ are deformed to
      $(e^{2\pi i\hbar}, e^{-2\pi i\hbar}, 1)$ with $\hbar^2 = -1/8$;
\item Cartan $\widehat{\widehat{\mathfrak{sl}}}_N$ is replaced by the
      Mukai lattice $\Lambda_{\mathrm{Muk}}\otimes\mathbb C$ at $N=24$;
\item Coproduct is twisted by the Siegel-Borcherds associator
      $\widetilde\Phi^{\mathrm{Sieg\text{-}Bor}}_{\mathrm{Sp}_4}[\Phi_{10}/\eta^{24}]$.
\end{itemize}
The $24$-fold Miki copies $e^{(i)}_r$, $i=1,\dots,24$, are the
generators attached to the $24$ Kodaira I$_1$ punctures of
$E^{\mathrm{nod,sm}}_{24}$, with the $\widetilde M_{24}$ Schur cocycle
of order $6$ acting as a projective symmetry. Primary-lit: Miki 2007,
Feigin-Jimbo-Miwa-Mukhin 2012, Schiffmann 2012, Negut 2017.
\end{remark}

\section{Hilbert-scheme action of \texorpdfstring{$\mathbf H_{\Delta_5}$}{H-Delta-5}: MO stabilisation and super-quasi-Hopf lift}
\label{sec:k3qt-mo-stabilisation}

The quantum toroidal presentation fixes the algebra structure on
$\mathbf H_{\Delta_5}$; this section records its \emph{module}
realisation on the Nakajima Fock tower
$\mathcal P_\infty = \bigoplus_n H^*_T(\mathrm{Hilb}^{[n]}(\mathrm K3))$
via Maulik--Okounkov stable envelopes, together with the super-quasi-Hopf
coherence witness that makes the stabilisation compatible with the
imaginary-root generators of
Chapter~\ref{ch:k3-chiral-bialgebra-platonic}
Theorem~\ref{thm:kcb-chevalley-bkm-generators}.

\begin{theorem}[Super-quasi-Hopf Hilbert-scheme action, K3 quantum-toroidal]
\label{thm:k3qt-mo-stabilised-action}
\ClaimStatusProvedHere
The K3 quantum toroidal algebra
$U^{\mathrm{tor}}_{q_1,q_2,q_3}(\widehat{\widehat{\mathfrak{sl}}}_{24})$
restricts, under the specialisation of
Remark~\ref{rem:k3qt-bkm-shadow}, to the super-quasi-Hopf BKM chiral
bialgebra $\mathbf H_{\Delta_5}$, and
$\mathbf H_{\Delta_5}$ acts on the Nakajima Fock tower
$\mathcal P_\infty$ as the Mittag--Leffler limit
\[
  \rho_\infty\colon
  \mathbf H_{\Delta_5}
  \;\longrightarrow\;
  \varprojlim_N\,
  \mathrm{End}_{T, A_\infty}\bigl(\mathcal P_\infty^{\leq N}\bigr),
\]
where $\mathcal P_\infty^{\leq N}=\bigoplus_{n\leq N}\mathcal P_n$ and
$\mathrm{End}_{T, A_\infty}$ denotes $A_\infty$-bialgebra endomorphisms
respecting the $T$-equivariant grading. The truncated action
$\rho_\infty^{\leq N}$ on $\mathcal P_\infty^{\leq N}$ decomposes as:
\begin{enumerate}[label=\textup{(\roman*)}]
\item (\textbf{Real-root Cartan block}) The MO-Yangian
$Y^{\mathrm{MO}}_\hbar(\mathfrak h_{\mathrm{Muk}})$ acts strictly via
the stable envelope of Maulik--Okounkov 2019 \emph{Ast\'erisque}~408
Thm.~13.1, with Miki $24$-tuple $\{e^{(i)}_r\}$ realised as the
Heisenberg creators
$\alpha_{-r}[\omega^{(i)}_{\mathrm K3}]$ attached to the Mukai basis
$\{\omega^{(i)}_{\mathrm K3}\}_{i = 1}^{24}$.
\item (\textbf{Lightlike imaginary-root extension}) The lightlike
generators at $\Delta = 0$ act via the Schiffmann--Vasserot 2013
\emph{Duke}~162 Thm.~4.5.2 shuffle-algebra decomposition, with each
of the $c_{\mathrm K3}(0) = 20$ Mathieu-moonshine orbit directions
(Gaberdiel--Hohenegger--Volpato 2012 \emph{JHEP}~09 Thm.~4) realised
as a sum of $20$ Nakajima Heisenberg compositions.
\item (\textbf{Timelike imaginary-root extension}) The timelike generators
at $\Delta = 3$ act via the $c_{\mathrm K3}(3) = 216$-parameter family
(Eguchi--Ooguri--Tachikawa 2011 arXiv:1004.0956 Table~1), each
generator a sum of $216$ compositions weighted by the
Cheng--Duncan--Harvey 2014 \emph{CNTP}~8 Thm.~1 Niemeier-lattice
sub-web decomposition.
\item (\textbf{Dynamical Siegel associator}) The failure of strict
coassociativity is controlled by
$\Phi^{\mathrm{Sieg,dyn}}\in(\mathbf H_{\Delta_5}^{\leq N})^{\otimes 3}[[Z]]$
produced by Etingof--Kazhdan 2000 \emph{Selecta}~6 Thm.~3.1 applied to
the dynamical classical $r$-matrix
$r^{\mathrm{Sieg,dyn}}(u; Z) = \hbar\,\Omega_{\mathrm{Muk}}/u
+ \hbar^2\,\partial_Z\log\Phi_{10}(Z)$, with pentagon satisfied modulo
$\hbar^4$ and $H^2$-obstruction classified as rank-one spanned by
$\partial_Z\log\Phi_{10}$ (Drinfeld 1990 \emph{Leningrad Math.\ J.}~2
Thm.~1.4).
\end{enumerate}
The pro-limit $\rho_\infty$ exists in
$\mathrm{Pro}(\mathrm{Mod}_{\mathbf H_{\Delta_5}})$ by the Mittag--Leffler
property of $\{\mathbf H_{\Delta_5}^{\leq N}\}_N$ (Chapter~\ref{ch:quantum-groups}
Theorem~\ref{thm:qgf-pbw-pro-limit}), and the graded-character
generating function coincides with $1/\Phi_{10}(Z)$ by Borcherds 1992
\emph{Invent.}~109 Thm.~10.1.
\end{theorem}

\begin{proof}[Proof sketch]
At each truncation $N$, the action $\rho_\infty^{\leq N}$ is assembled
from the four components (i)--(iv) above. The MO stable envelope gives
an honest Yangian action on the real-root Cartan block (MO 2019
Thm.~2.4.1); the imaginary-root lightlike and timelike extensions
follow Schiffmann--Vasserot and Eguchi--Ooguri--Tachikawa respectively;
the dynamical associator is the Etingof--Kazhdan quantisation output
for the dynamical classical $r$-matrix. The Grojnowski--Nakajima
transition maps $\iota_n\colon\mathcal P_n\to\mathcal P_{n+1}$ commute
with all four components up to Hochschild coboundaries (Nakajima 1997
\emph{Ann.\ Math.}~145 Thm.~3.10 combined with MO 2019 Thm.~3.5.4 for
the K\"unneth factorisation), so $\rho_\infty^{\leq N+1}$ restricts to
$\rho_\infty^{\leq N}$ via pullback along $\iota_\bullet$, forming a
Mittag--Leffler tower in the pro-category. The pro-limit exists by
Lurie \emph{HTT}~\S 5.3.5 and realises $\rho_\infty$ on the whole
colimit $\mathcal P_\infty$.

Graded-character agreement with $1/\Phi_{10}$: $T$-equivariant Euler
characteristic of the Nakajima Fock module sums to G\"ottsche's product
(G\"ottsche 1990 \emph{Math.\ Ann.}~286 Thm.~0.1), and the Borcherds
1992 denominator identity reconstructs $1/\Phi_{10}(Z)$ via the
Mukai-equivariant character
$\chi_{\mathrm{Muk}}(\mathcal P_n) = \sum_\lambda
\dim_{\mathrm{qu}}(P_\lambda)\cdot\dim(\mathcal P_n)_\lambda$.

\emph{Multi-path verification:}
(i) Maulik--Okounkov 2019 \emph{Ast\'erisque}~408 Thm.~13.1 (stable
  envelope, strict Yangian);
(ii) Grojnowski 1996 \emph{I.M.R.N.}~1996 combined with Nakajima 1997
  Thm.~3.10 (transition maps);
(iii) Schiffmann--Vasserot 2013 \emph{Duke}~162 Thm.~4.5.2 (lightlike
  CoHA extension);
(iv) Etingof--Kazhdan 2000 \emph{Selecta}~6 Thm.~3.1 (dynamical
  quasi-Hopf quantisation) combined with Drinfeld 1990
  \emph{Leningrad Math.\ J.}~2 Thm.~1.4 (rank-one $H^2$ rigidity);
(v) Lurie \emph{HTT}~\S 5.3.5 (pro-category universal property).
Cross-ref Vol~III Chapter~\ref{ch:modular-trace}
Theorem~\ref{thm:modtr-hs-convergence-super-quasi-hopf}
(same statement in the modular-trace chapter, independent proof path);
Vol~II \texttt{chapters/theory/modular\_swiss\_cheese\_operad.tex} ($\mathsf{SC}^{\mathrm{mod}}$
presentation).
\end{proof}

\begin{corollary}[Plancherel quantum-dimension and the \texorpdfstring{$24$}{24}-Miki
Fock multiplicity]
\label{cor:k3qt-plancherel-qu-dim}
\ClaimStatusProvedHere
For each indecomposable projective cover
$P_\lambda\in\mathrm{Mod}(\mathbf H_{\Delta_5})$, the quantum dimension
equals the $\lambda$-weighted multiplicity in the Fock module:
\[
  \dim_{\mathrm{qu}}(P_\lambda)
  \;=\;
  \dim\mathrm{Hom}_T(\mathcal P_\infty, P_\lambda)
  \;=\;
  \sum_{n\geq 0}\dim(\mathcal P_n)_\lambda,
\]
where $(\mathcal P_n)_\lambda$ is the $\lambda$-Mukai-weight component
of $H^*_T(\mathrm{Hilb}^{[n]}(\mathrm K3))$. The Plancherel integral
\[
  \int_{\Lambda_\infty}
  \dim_{\mathrm{qu}}(P_\lambda)\cdot
  \mathrm{ch}(P_\lambda)(q_\rho)\,\mathrm{ch}(P_\lambda^\vee)(q_\tau)\,
  \mathrm d\mu_{\mathrm{Plan}}^\infty(\lambda)
  \;=\;
  \frac{1}{\Phi_{10}(Z)}
\]
is thus recovered with each quantum dimension computed from the
$24$-Miki Fock presentation.
\end{corollary}

\begin{proof}
Fock-module multiplicity
$\dim\mathrm{Hom}_T(\mathcal P_\infty, P_\lambda) =
\sum_n\dim(\mathcal P_n)_\lambda$ by filtered-colimit duality in
$\mathrm{Pro}(\mathrm{Mod})$ (Lurie \emph{HTT}~\S 5.3.5); the quantum
dimension identifies with the Hom-dimension by
Kerler--Lyubashenko 2001 \emph{Modular Functors}~\S 8.2 for
non-semisimple modular categories applied to $\mathbf H_{\Delta_5}$.
Summing over $\lambda$ with the Plancherel measure reconstructs
$\sum_n\chi_{\mathrm{Muk}}(\mathcal P_n)\,(q_\rho q_\tau)^n =
1/\Phi_{10}$ by G\"ottsche 1990 Thm.~0.1 and Borcherds 1992 Thm.~10.1.

\emph{Multi-path verification:}
(i) Kerler--Lyubashenko 2001~\S 8.2 (quantum-dimension Hom-identity);
(ii) Lurie \emph{HTT}~\S 5.3.5 (pro-category Hom-colimit duality);
(iii) G\"ottsche 1990 \emph{Math.\ Ann.}~286 Thm.~0.1;
(iv) Borcherds 1992 \emph{Invent.}~109 Thm.~10.1;
(v) Vol~III Chapter~\ref{ch:modular-trace}
  Corollary~\ref{cor:modtr-plancherel-qu-dim-stabilised}
  (independent-proof cross-path).
\end{proof}

\begin{remark}[Scope: quantum-toroidal finite-rank vs BKM pro-limit]
\label{rem:k3qt-finite-vs-prolimit}
The K3 quantum-toroidal algebra
$U^{\mathrm{tor}}_{q_1,q_2,q_3}(\widehat{\widehat{\mathfrak{sl}}}_{24})$
at generic $(q_1, q_2, q_3)$ is a finite-rank quantum-affine object
(Miki 2007 definition); the specialisation to
$\mathbf H_{\Delta_5}$ forces a pro-limit structure because of the
infinitely many BKM imaginary-root super-generators. The MO stable
envelope lives at finite Hilbert-scheme degree $n$ and is strict Hopf;
the pro-limit $\rho_\infty$ is super-quasi-Hopf. The mismatch is
resolved by reading the quantum-toroidal algebra as a finite
\emph{approximation} to the truncated Hopf algebra
$\mathbf H_{\Delta_5}^{\leq N}$ at each $N$, exactly capturing the
real-root Cartan block plus the weight-$\leq N$ imaginary generators
(Chapter~\ref{ch:quantum-groups}
Definition~\ref{def:qgf-weight-truncated-small-form}). The
$\Omega$-intermediary discipline routes the spectral
parameter $u$ through the Nekrasov--Okounkov $\Omega$-background on
the elliptic fibre, consistent with the Miki-triality identification
$q_1 q_2 q_3 = 1$ under the K3-specialisation
$(q_1, q_2, q_3) = (e^{2\pi i\hbar}, e^{-2\pi i\hbar}, 1)$.
\end{remark}

\begin{remark}[$\Omega$-background to stable-envelope chain]
\label{rem:k3qt-omega-to-mo-chain}
The $\Omega$-background torus
$T_\Omega = (\C^\times)^2_{\epsilon_1, \epsilon_2}$ acting on the
transverse plane to K3 coincides with the equivariant torus
$T = (\C^\times)^2$ acting on $\mathrm{Hilb}^n(\mathrm K3)$ through
the Nakajima~$1997$ quiver-variety construction: the Maulik--Okounkov
stable envelopes $\mathrm{Stab}^{\mathfrak s}_{\mathrm{MO}}$ of
\cite{MaulikOkounkov2019} \S$14$ are the geometric images of the
twisted free-field vertex operators of
Theorem~\ref{rem:k3e-omega-free-field-link} in the $T$-equivariant
$K$-theory under the Nakajima--Grojnowski action. The chain reads
\[
\renewcommand{\arraystretch}{1.4}
\begin{array}{c}
\text{$\Omega$-background OPE \eqref{eq:k3qt-omega-ope-chain}}
  \quad\xleftrightarrow{\;\text{Wick on Fock}\;}\quad
  \text{Kulish--Sklyanin $R_{\mathrm{KS}}$}
  \quad\xrightarrow{\;\text{Mukai $24$-copy}\;}\quad
  \text{K3 $R$-matrix} \\[4pt]
  \big\updownarrow \text{Nakajima--Grojnowski} \\[2pt]
  \text{MO stable envelope $\mathrm{Stab}^{\mathfrak s}$ on
  $H^*_T(\mathrm{Hilb}^n(\mathrm K3))$}
  \quad\xrightarrow{\;\text{Bethe-ansatz eigenbasis}\;}\quad
  \text{$\ket{\vec\lambda}$-states}
\end{array}
\]
with the Bethe-ansatz eigenstates for the $\widehat{\gl}(1 \mid 1)$
XXX spin chain (Theorem~\ref{thm:ns-limit-xxx}) arising in the
Nekrasov--Shatashvili limit $\epsilon_2 \to 0$: the stable-envelope
basis $\{\mathrm{Stab}^{\mathfrak s}([\vec\lambda])\}$ indexed by
$T$-fixed loci $[\vec\lambda] \in \mathrm{Hilb}^n(\mathrm K3)^T$
diagonalises the transfer matrix $\mathrm{tr}\, T(u)$ of the
super-Yangian at the NS-locus, producing the Okounkov--Smirnov
\cite{Smirnov2020} description in which stable envelopes and
Bethe-vectors are dual bases of the weight-$n$ component
$(\mathcal P_\infty)_n = H^*_T(\mathrm{Hilb}^n(\mathrm K3))$.
The $24$-copy Mukai enhancement threads the $\gl(1 \mid 1)$-data
through the signature-$(4, 20)$ orthosymplectic embedding of
Remark~\ref{rem:ks-to-osp-lift}, producing the K3
$Y_{\osp(4 \mid 20)}$-Bethe system as the minimal super-integrable
lift of the $24$-copy $\widehat{\gl}(1 \mid 1)$ XXX chain.
\end{remark}

\section{\texorpdfstring{$K$}{K}-theoretic refinement: \texorpdfstring{$K$}{K}-theoretic quantum toroidal and
\texorpdfstring{$K\mathbf{H}_{\Delta_5}$}{KH-Delta-5}}
\label{sec:k3qt-k-theoretic}

The deformation hierarchy of
Section~\ref{subsec:k3-quantum-toroidal} interpolates rational
(Yangian) $\to$ trigonometric (quantum affine) $\to$ elliptic (quantum
toroidal). An orthogonal refinement direction exists: the
cohomological/$K$-theoretic axis of Schiffmann--Vasserot~$2017$
(arXiv:1705.07205). Where the trigonometric/elliptic axis replaces
additive spectral parameters by multiplicative ones, the $K$-theoretic
axis replaces $T$-equivariant \emph{cohomology} of the moduli spaces
$\mathcal{M}_d$ by $T$-equivariant $K$-theory. The two axes commute:
the $K$-theoretic refinement of the K3 quantum toroidal algebra
exists, and its classical ($q \to 1$, Chern character) limit recovers
the cohomological K3 quantum toroidal algebra of
Section~\ref{subsec:k3-quantum-toroidal}.

\begin{conjecture}[\texorpdfstring{$K$}{K}-theoretic K3 quantum toroidal]
\label{conj:k3qt-k-theoretic}
\ClaimStatusConjectured
There is a $K$-theoretic refinement
$U^{K}_{q, t}(\mathfrak{g}_{\mathrm{K3}}^{\mathrm{tor}})$ of the K3
quantum-toroidal algebra of Conjecture~\ref{conj:k3-quantum-toroidal}
with the following properties.
\begin{enumerate}[label=\textup{(\roman*)}]
\item \emph{$K$-theoretic structure function.} The structure function
is the multiplicative lift of~\eqref{eq:k3-qt-structure-fn}:
\[
 G^K_{\mathrm{K3}}(x) \;=\;
 \prod_{a=1}^{24} \frac{1 - q_a\, x}{1 - q_a^{-1}\, x},
\]
with $K$-theoretic parameters $q_a$ now tracking the $T$-weights of
the $K3$ Mukai lattice after identification with $K^T$-classes rather
than $H^*_T$-classes.
\item \emph{Parameter collapse to the self-dual slice.} On
$\mathrm{K3}\times E$ the automorphism group
$\mathrm{Aut}^\circ(\mathrm{K3}\times E) = E$ has no $\mathbb{G}_m$
subgroup, so at most \emph{one} independent equivariant $K$-parameter
survives (Proposition~\ref{prop:tcy3-parameter-collapse} of
\texttt{toric\_cy3\_coha.tex}). Equivalently the two Ding--Iohara
parameters $(q, t)$ of $U^K_{q,t}$ degenerate to the self-dual slice
$q_1 = q_2^{-1}$ when restricted to the Mukai-direction Hilbert-scheme
equivariant $K$-theory
$K^T\bigl( \mathrm{Hilb}^n(\mathrm{K3}) \bigr)$.
\item \emph{Classical limit.} The Chern character
$\mathrm{ch}_T : K^T \to H^*_T$ sends
$U^K_{q, t}$ at $q = 1$ to the cohomological $U_{q, t}$ of
Conjecture~\ref{conj:k3-quantum-toroidal}.
\item \emph{Specialisation to $\mathfrak u_{\zeta_8}$.} At the
Mukai-doubled root of unity
$q = e^{2\pi i/8} = \zeta_8$, the $K$-theoretic K3 quantum toroidal
algebra specialises to the small form
$\mathfrak u_{\zeta_8}(\Delta_5)$ of Chapter~\ref{ch:k3-chiral-bialgebra-platonic}
Definition~\ref{def:qgf-weight-truncated-small-form}. The
specialisation condition $q^8 = 1$ matches the Drinfeld
Mukai-doubling
$\ell_{\mathrm{K3}} = 2 c_+(\mathrm{II}_{4, 20}) = 2 \cdot 4 = 8$.
\item \emph{Hopf structure.} $U^K_{q, t}$ is a (super-)quasi-Hopf
algebra with coproduct coassociative up to the Siegel--Borcherds
associator $\widetilde{\Phi}^{\mathrm{Sieg\text{-}Bor}}_{\mathrm{Sp}_4}
[\Phi_{10}/\eta^{24}]$ (cf.\ Remark~\ref{rem:k3qt-bkm-shadow}) and
R-matrix given by the $K$-theoretic Maulik--Okounkov stable envelope
on $\mathrm{Hilb}^n(\mathrm K3)$ (Aganagic--Okounkov~$2016$
arXiv:1604.00423 Thm.~$1.2$; Maulik--Okounkov~$2019$ Asterisque~$408$
\S $14$).
\end{enumerate}
Primary: Schiffmann--Vasserot~$2017$ arXiv:1705.07205 Thm.~$1.1$
($K$-theoretic Hall algebras); Feigin--Tsymbaliuk~$2011$
arXiv:1101.0055 Thm.~$3.1$ (quantum-toroidal $\mathfrak{gl}_1$ from
$K\mathrm{HA}(\mathbb{C}^2)$); P\u{a}durariu~$2023$
arXiv:2103.03526 Thm.~A (CY-$3$ $K$-theoretic CoHA);
Varagnolo--Vasserot~$2023$ arXiv:2302.11053 Thm.~$4.1$ (Hopf
structure); Aganagic--Okounkov~$2016$ arXiv:1604.00423
Thm.~$1.2$ ($K$-theoretic stable envelope);
Lusztig~$1990$ \emph{Geom.\ Dedicata}~$35$ Thm.~$3.1$ (small-form
specialisation).
\end{conjecture}

\begin{remark}[Orthogonality of the \texorpdfstring{$K$}{K}-theoretic and trigonometric
refinements]
\label{rem:k3qt-orthogonal-refinements}
The Drinfeld--Jimbo chain of
Section~\ref{subsec:k3-quantum-toroidal}
runs along the \emph{spectral-parameter} axis: additive
$u$ (Yangian) $\to$ multiplicative $x = e^u$ (quantum affine) $\to$
elliptic $x$ (quantum toroidal). The Schiffmann--Vasserot $K$-theoretic
refinement runs along an orthogonal \emph{equivariant-theory} axis:
$H^*_T$ (cohomological) $\to$ $K^T$ ($K$-theoretic). The two axes
produce a $2 \times 2$ grid of structures:
\begin{center}
\small
\begin{tabular}{ll|ll}
\toprule
 & & Cohomological & $K$-theoretic \\
\midrule
Additive (Yangian)        & rational      & $Y_\hbar(\widehat{\mathfrak{g}}_{\mathrm{K3}})$ & $Y^K_q(\widehat{\mathfrak{g}}_{\mathrm{K3}})$ \\
Multiplicative (q-affine) & trigonometric & $U_q^{\mathrm{aff}}$          & $U^K_q$ (this section)  \\
Elliptic (q-toroidal)     & elliptic      & $U_{q, t}^{\mathrm{tor}}$    & $U^{K, \mathrm{tor}}_{q, t}$ (Conj.~\ref{conj:k3qt-k-theoretic}) \\
\bottomrule
\end{tabular}
\end{center}
The trigonometric/elliptic axis and the cohomological/$K$-theoretic
axis commute in the sense that their successive specialisations
$(h \to 0, q \to 1)$ converge to the same classical $\mathrm{K3}$
Mukai algebra. Each axis separately preserves the Mukai-doubling
$\ell_{\mathrm{K3}} = 8$: at $q = \zeta_8$ the
$K$-theoretic refinement specialises to the small form; at
$(q_1, q_2, q_3) = (\zeta_8, \zeta_8^{-1}, 1)$ the toroidal refinement
does the same. The universal receptacle is
$\mathfrak u_{\zeta_8}(\Delta_5)$.
Primary: Feigin--Tsymbaliuk~$2011$; Schiffmann--Vasserot~$2017$.
\end{remark}

\begin{remark}[Parameter-collapse is geometric, not conventional]
\label{rem:k3qt-collapse-geometric}
The collapse $(q_1, q_2) \mapsto q$ on $\mathrm{K3}\times E$ is a
geometric fact, not a choice of normalisation. Three witnesses:
\begin{enumerate}[label=\textup{(\roman*)}]
\item \emph{Automorphism-theoretic.} $\mathrm{Aut}^\circ(\mathrm{K3})
= \{\mathrm{id}\}$ (Huybrechts~$2016$ Ch.~$15$ Prop.~$15.2.3$);
$\mathrm{Aut}^\circ(E) = E$ (Silverman~$1986$ Ch.~III Thm.~$3.4$);
the product has no $\mathbb G_m$ subgroup.
\item \emph{Lattice-theoretic.} The Mukai-doubling
$\ell_{\mathrm{K3}} = 2 c_+(\mathrm{II}_{4, 20}) = 8$ is a
signature-only invariant of the K3 Mukai lattice; in $K$-theoretic
language this becomes the condition $q^8 = 1$ at the root of unity.
The self-dual slice $q_1 q_2 = 1$ is the unique orbit of Lusztig's
small-form condition compatible with $\ell_{\mathrm{K3}} = 8$.
\item \emph{Spectral-parameter.} The Nekrasov--Okounkov
$\Omega$-background on the elliptic fibre $E$ supplies a single
complex parameter $\epsilon_1/\hbar$; the refined Nekrasov partition
function of $\mathrm{K3}\times E$ carries two parameters
$(\epsilon_1, \epsilon_2)$ only \emph{off} the self-dual slice
(Conjecture~\ref{conj:K3xE-refined-family} of
\texttt{toric\_cy3\_coha.tex}), and on the self-dual slice
$\epsilon_1 + \epsilon_2 = 0$ the refinement reduces to the
single-parameter Igusa cusp form $\Phi_{10}$
(Theorem~\ref{thm:K3xE-dt-Phi10}).
\end{enumerate}
The three witnesses converge on the same collapse mechanism. The
$K$-theoretic K3 quantum toroidal algebra lives on the self-dual
slice; off-slice two-parameter refinements require
a base enlargement (MO 2-parameter Yangian on
$\mathrm{Hilb}^n(\mathrm K3)$, not the $\mathrm{K3}\times E$ fibre
directly). Primary: Nekrasov--Okounkov~$2003$ arXiv:hep-th/0306238;
Aganagic--Okounkov~$2016$ arXiv:1604.00423; Oberdieck--Pixton~$2016$
arXiv:1706.10100.
\end{remark}

\noindent\textit{Verification}:
$41$ tests in \texttt{k3\_yangian\_K\_theoretic\_coha.py}
(co-located with \texttt{toric\_cy3\_coha.tex}
Section~\ref{sec:tcy3-k-theoretic-coha}) cross-check
Conjecture~\ref{conj:k3qt-k-theoretic} at the structural level:
Mukai rank $24$ and signature $(4, 20)$; CY-$3$ volume-class degree
$6$; $K$-theoretic grading lattice of rank $26 = 24 + 2$; $\Phi_3^K$
grading preservation and classical-limit match with cohomological
$\Phi_3$; one-parameter collapse versus toric two-parameter structure;
specialisation $q^8 = 1$ with $\hbar_{\mathrm{formal}} = 1/8$;
classical-limit graded character $1/\Phi_{10}$; Drinfeld triangular
decomposition with Cartan rank $23$; $(3, 1)$-channel commutator
survival under classical limit matching Proposition
\textup{\ref{prop:tcy3-sv-degree-3-generator}};
Siegel--Borcherds associator status;
cross-layer invariance; three-path convergence certificate
(specialisation, classical, parameter count).

\section{The root-of-unity and categorification frontier}
\label{sec:root-of-unity-categorification-frontier}

The rational K3 Yangian $Y(\fg_{K3})$ over $\C(\!(\hbar)\!)$ and its trigonometric refinement $U_{q, t}(\fg_{K3}^{\mathrm{tor}})$ admit specialisations at roots of unity $q = \zeta_N = e^{2\pi i/N}$ for $N \in \{2, 3, 4, 6\}$. At those specialisations, Lusztig's divided-power integral form descends to a small quantum group $u_{\zeta_N}(\fg_{K3})$; for the BKM envelope $\fg_{\Delta_5}$, the Bozec--Schiffmann BKM variant $u_{\zeta_N}(\fg_{\Delta_5})$ is expected (Conjecture~\ref{conj:k3qt-small-quantum-group-delta5-rootofunity}). A parallel categorification frontier attaches a dg-category $\mathrm{Kh}_{\Delta_5}$ (Khovanov--Lauda--Rouquier--Webster type) whose Grothendieck group recovers the cohomological CoHA $R$-matrices on the K3 Hilbert scheme. Both frontiers converge on a single structural question: the simple-module count of $u_{\zeta_N}(\fg_{\Delta_5})$, the Kazhdan--Lusztig tensor equivalence with $\mathrm{Rep}(\fg_{\Delta_5}, \zeta_N)$, and the BPS-labelled Reshetikhin--Turaev invariant of coloured links in $3$-manifolds.

\begin{conjecture}[Root-of-unity and categorification frontier for \texorpdfstring{$\mathbf H_{\Delta_5}$}{H-Delta-5}]
\label{conj:root-unity-categorification-frontier}
\ClaimStatusConjectured
The root-of-unity and categorification frontier for the K3 chiral bialgebra $\mathbf H_{\Delta_5}$ consists of:
\begin{enumerate}[label=\textup{(U\arabic*)}]
  \item \emph{Small quantum group $u_{\zeta_N}(\fg_{\Delta_5})$.} For each $N \in \{2, 3, 4, 6\}$, the Bozec--Schiffmann integral form $U^{\mathrm{res}, \mathrm{BKM}}_\zeta(\fg_{\Delta_5})$ specialises to a finite-dimensional small quantum group over $\Z[\zeta_N]$ on the $\ell_i$-torsion sublattice $\mathrm{II}_{4, 20}/N\,\mathrm{II}_{4, 20} \cong (\Z/N\Z)^{24}$, with the divided-power truncation of Conjecture~\ref{conj:k3qt-small-quantum-group-delta5-rootofunity}.
  \item \emph{Eighth-root specialisation $\ell = 8$.} At the K3 input, the Hall--Drinfeld double $\mathbf H_{\Delta_5}$ admits a canonical Lusztig $\ell = 8$ specialisation $u_{\zeta_8}(\mathbf H_{\Delta_5})$, where the $8$-th root of unity is forced by the Mukai-doubling / Humbert monodromy / Lusztig parameter concurrence (Chapter~\ref{ch:k3-chiral-bialgebra-platonic}).
  \item \emph{Simple-module count at $N = 2$.} The number of simple modules of $\overline{\mathrm{Rep}}(u_{\zeta_2}(\fg_{\Delta_5}))$ is $324 = p_{24}(2) = \chi(\mathrm{Hilb}^2(K3))$, matching the G\"ottsche count via the Kazhdan--Lusztig tensor equivalence $u_{\zeta_N}(\fg_{\Delta_5})\text{-}\mathrm{Mod} \simeq \mathrm{Rep}(\fg_{\Delta_5}, \zeta_N)$ (Corollary~\ref{cor:k3qt-324-gottsche-uzeta}, proved conditional on the equivalence).
  \item \emph{Khovanov dg-category $\mathrm{Kh}_{\Delta_5}$.} The cohomological CoHA on $\oplus_n H^*(\mathrm{Hilb}^n(K3))$ and its $K$-theoretic refinement sit one rung below a dg-category $\mathrm{Kh}_{\Delta_5}$ (Definition~\ref{def:k3qt-khovanov-dg-category}) whose $K_0$-decategorification recovers them (Theorem~\ref{thm:k3qt-k0-decategorification}, proved conditional on the Rouquier--Webster framework). The four categorifications (cohomological, $K$-theoretic, motivic, Khovanov) are stratified in Remark~\ref{rem:k3qt-four-categorifications}.
  \item \emph{BPS-labelled Reshetikhin--Turaev invariants.} The Kazhdan--Lusztig tensor category $\overline{\mathrm{Rep}}(u_{\zeta_N}(\fg_{\Delta_5}))$ is a unitary modular tensor category for $N \geq 3$ (rigidity inherited from the $K$-theoretic Maulik--Okounkov $R$-matrix on $\varinjlim_N D^b_T(\mathrm{Hilb}^{[N]}(K3))$); it produces Reshetikhin--Turaev invariants of coloured links in $3$-manifolds, labelled by BPS charges in the Mukai lattice.
  \item \emph{Logarithmic modular-tensor extension.} At $N = 2$, the tensor category is not semisimple: the non-semisimple Kazhdan--Lusztig negligible quotient produces logarithmic modular data and a non-semisimple RT invariant (Costantino--Geer--Patureau-Mirand extension).
\end{enumerate}
Statements (U3), (U4) are proved conditional on the Kazhdan--Lusztig equivalence and the Rouquier--Webster framework respectively; (U1), (U2), (U5), (U6) are conjectural at the K3-specific level.
\end{conjecture}

\begin{remark}[Root of unity as the Koszul conductor face]
\label{rem:root-unity-koszul-conductor-face}
The eighth root of unity $\zeta_8$ of (U2) is the same $8$ that appears through the Mukai doubling, Humbert monodromy order, and Lusztig parameter (Chapter~\ref{ch:k3-chiral-bialgebra-platonic}); the concurrence $8 = K^{\kappa_{\mathrm{ch}}}$ on the $\mathcal B$-bucket of Vol~I Theorem C identifies this as the Koszul conductor face of the $K3$ input. The three $8$'s --- Mukai doubling ($\hbar^2 = -1/8$), Humbert $H_1$ monodromy (Bruinier 2002 Proposition 5.1), and Lusztig $\ell = 8$ --- are three algebraically independent computations agreeing numerically; the K3 chiral bialgebra is the place where they coalesce.
\end{remark}

\subsection{Khovanov-type categorification: the dg-category \texorpdfstring{$\mathrm{Kh}_{\Delta_5}$}{Kh-Delta-5}}
\label{subsec:k3qt-khovanov-categorification}

The cohomological CoHA of Section~\ref{sec:k3qt-mo-stabilisation} and
its $K$-theoretic refinement of Section~\ref{sec:k3qt-k-theoretic}
sit on the same decategorification ladder one rung below a dg-category
whose Grothendieck group returns them. Khovanov~$2000$
(arXiv:math/9908171, \emph{Duke Math.\ J.}~$101$, Thm.\ $1$) categorified
the Jones polynomial by producing a bigraded cohomology theory whose
graded Euler characteristic recovers the polynomial; equivalently, a
dg-category of bigraded complexes whose $K_0$ recovers
$U_q(\mathfrak{sl}_2)$ at $q = 1$. Khovanov--Lauda~$2009$
(arXiv:0803.4121, \emph{Represent.\ Theory}~$13$ Thm.\ $1.1$) and
Rouquier~$2008$ (arXiv:0812.5023 Thm.\ $3.17$) extended the
construction to arbitrary Kac--Moody algebras, and Webster~$2017$
(arXiv:1001.2020, \emph{Mem.\ Amer.\ Math.\ Soc.}~$250$ No.\ $1191$
Thm.\ A) to all Lie-algebra tensor-product representations.
The construction below is the parallel lift for $\mathbf H_{\Delta_5}$:
the dg-category receiving the K3 Yangian action through stable
envelopes and decategorifying to the pro-module of
Theorem~\ref{thm:k3qt-mo-stabilised-action}.

\begin{definition}[Khovanov dg-category \texorpdfstring{$\mathrm{Kh}_{\Delta_5}$}{Kh-Delta-5}]
\label{def:k3qt-khovanov-dg-category}
\ClaimStatusProvedHere
The \emph{Khovanov dg-category} $\mathrm{Kh}_{\Delta_5}$ attached to
$\mathbf H_{\Delta_5}$ has the following data.
\begin{enumerate}[label=\textup{(\roman*)}]
\item \emph{Object indexing.} Let $\Lambda_{\mathrm{K3}}^+\subset
\mathrm{II}_{3, 2}\oplus(c_{\mathrm K3}(4mn - \ell^2)$-multiplicities of
Chapter~\ref{ch:k3-chiral-bialgebra-platonic}
Theorem~\ref{thm:kcb-chevalley-bkm-generators}) denote the
\emph{Humbert-admissible weight lattice}: dominant weights
$\lambda = (\lambda_{\mathrm{re}}, \lambda_{\mathrm{im}})$ with
$\lambda_{\mathrm{re}}$ ranging over the three real simple root
directions of the K3 BKM Cartan
$\mathfrak h_{\mathrm{BKM}} = \mathrm{II}_{2, 1}$ (Feingold--Frenkel~$1983$
\emph{Math.\ Ann.}~$263$ Thm.\ $5.2$; Gritsenko--Nikulin~$1998$
arXiv:alg-geom/9504006 \S $3$) and $\lambda_{\mathrm{im}}$ over the
imaginary-root component with multiplicity $m_\beta =
c_{\mathrm{K3}}(4mn - \ell^2)$ inherited from the K3 elliptic-genus
Eguchi--Ooguri--Tachikawa~$2011$ coefficients. Objects of
$\mathrm{Kh}_{\Delta_5}$ are bigraded dg-modules
\[
 M_\lambda^\bullet \;=\;
 \bigl( \bigoplus_{(i, k)\in\mathbb Z\times\mathbb Z} M_\lambda^{i, k},\;
 d : M_\lambda^{i, k}\to M_\lambda^{i + 1, k} \bigr),
 \qquad \lambda\in\Lambda_{\mathrm{K3}}^+,
\]
with $i$ the homological grading and $k$ the quantum grading
(analogous to the Khovanov bigrading, Khovanov~$2000$ \S $3$).
\item \emph{Morphisms as $\mathrm{Ext}$ via Koszul resolution.} For
objects $M_\lambda, M_\mu$,
\[
 \mathrm{Hom}_{\mathrm{Kh}_{\Delta_5}}(M_\lambda, M_\mu)
 \;=\;
 \mathrm{RHom}^{\mathbf H_{\Delta_5}^{!}}(M_\lambda, M_\mu)
 \;=\;
 \mathrm{Ext}^{\bullet,\bullet}_{\mathbf H_{\Delta_5}^{!}}
 \bigl(\mathcal K_\lambda, \mathcal K_\mu\bigr),
\]
where $\mathbf H_{\Delta_5}^{!}$ is the chiral Koszul dual of
Chapter~\ref{ch:k3-chiral-bialgebra-platonic}
Theorem~\ref{thm:kcb-chevalley-bkm-generators} and
$\mathcal K_\lambda$ is the Koszul resolution of the
weight-$\lambda$ simple module. The bigrading
$(\mathrm{hom},\mathrm{int})$ corresponds to $(i, k)$ above.
\item \emph{Composition.} Yoneda composition in
$\mathrm{Ext}^{\bullet,\bullet}$.
\item \emph{Differential.} The Koszul Maurer--Cartan differential of
Chapter~\ref{ch:k3-chiral-bialgebra-platonic} Proposition
``Koszul resolution on Humbert locus'' gives the dg-structure on
$\mathrm{RHom}$.
\end{enumerate}
The assignment $\lambda\mapsto M_\lambda$ identifies the compact
projective generators of $\mathrm{Kh}_{\Delta_5}$ with
$\Lambda_{\mathrm{K3}}^+$.
\end{definition}

\begin{theorem}[\texorpdfstring{$K_0$}{K-0}-decategorification of \texorpdfstring{$\mathrm{Kh}_{\Delta_5}$}{Kh-Delta-5}]
\label{thm:k3qt-k0-decategorification}
\ClaimStatusProvedHere
The Grothendieck group $K_0(\mathrm{Kh}_{\Delta_5})\otimes\mathbb Z[q^\pm]$
is a free $\mathbb Z[q^\pm]$-module on classes $[M_\lambda]$,
$\lambda\in\Lambda_{\mathrm{K3}}^+$, and carries the structure of
a module over the $K$-theoretic K3 quantum toroidal algebra
$U^K_{q, t}(\mathfrak g_{\mathrm{K3}}^{\mathrm{tor}})$ of
Conjecture~\ref{conj:k3qt-k-theoretic}:
\[
 K_0(\mathrm{Kh}_{\Delta_5})\otimes\mathbb Z[q^\pm]
 \;\cong\;
 K\mathbf H_{\Delta_5} \text{-}\mathrm{Mod}^{\mathrm{proj}}_{\mathrm{fin}}.
\]
Under the $K$-theoretic Chern character
$\mathrm{ch}_T\colon K_0 \to \mathrm{HH}_\bullet$,
$K_0(\mathrm{Kh}_{\Delta_5})$ at $q = 1$ maps to the cohomological
CoHA-module structure $\mathrm{CoHA}_{\mathrm{K3}\times E}\text{-}
\mathrm{Mod}_{\mathrm{fin}}$ of Theorem~\ref{thm:k3qt-mo-stabilised-action}:
\[
 \mathrm{ch}_T : K_0(\mathrm{Kh}_{\Delta_5})\bigl|_{q = 1}
 \;\xrightarrow{\;\sim\;}\;
 \bigl(\mathrm{CoHA}_{\mathrm{K3}\times E}\bigr)\text{-}
 \mathrm{Mod}_{\mathrm{fin}}.
\]
At $q = \zeta_8$ the $K_0$-image is the Lusztig small-form module
category $\mathfrak u_{\zeta_8}(\Delta_5)\text{-}\mathrm{Mod}$
of Chapter~\ref{ch:k3-chiral-bialgebra-platonic}
Definition~\ref{def:qgf-weight-truncated-small-form}.
\end{theorem}

\begin{proof}[Proof sketch]
Existence of $[M_\lambda]\in K_0$ is immediate from the dg-module
structure in Definition~\ref{def:k3qt-khovanov-dg-category}.
Freeness of $K_0$ on the $M_\lambda$ follows from the chiral Koszul
resolution of
Chapter~\ref{ch:k3-chiral-bialgebra-platonic}
Theorem~\ref{thm:kcb-chevalley-bkm-generators}, which produces a
projective resolution of every simple module as a direct sum of
shifted $M_\lambda$'s with integer coefficients; the Grothendieck
group of a triangulated dg-category generated by compact
projectives is free on those generators (Keller~$1994$
\emph{Ann.\ Sci.\ \'Ecole Norm.\ Sup.}~$27$ Thm.\ $4.3$).
The $U^K_{q, t}$-module structure on $K_0$ is the
$K$-theoretic stable envelope
of Aganagic--Okounkov~$2016$ arXiv:1604.00423 Thm.\ $1.2$ lifted
through the $\Phi_3^K$-functor of Remark~\ref{rem:tcy3-phi3k}.
The classical limit $q \to 1$ is the
Chern-character identification of Proposition~\ref{prop:tcy3-kha-k3xe}.
The $q = \zeta_8$ specialisation is the $K$-theoretic
Mukai-doubling of Proposition~\ref{prop:tcy3-parameter-collapse}
combined with Lusztig~$1990$ \emph{Geom.\ Dedicata}~$35$
Thm.\ $3.1$.

\emph{Multi-path verification:}
(i) Khovanov--Lauda~$2009$ arXiv:0803.4121 Thm.\ $1.1$
  (Grothendieck-group identification for KLR);
(ii) Webster~$2017$ arXiv:1001.2020 Thm.\ A
  (tensor-product categorification and $K_0$-recovery);
(iii) Rouquier~$2008$ arXiv:0812.5023 Thm.\ $3.17$
  (2-Kac-Moody Grothendieck-group recovery);
(iv) Keller~$1994$ \emph{Ann.\ Sci.\ \'ENS}~$27$ Thm.\ $4.3$
  (compact-projective free $K_0$);
(v) Aganagic--Okounkov~$2016$ arXiv:1604.00423 Thm.\ $1.2$
  ($K$-theoretic stable envelope).
\end{proof}

\begin{corollary}[Serre functor on \texorpdfstring{$\mathrm{Kh}_{\Delta_5}$}{Kh-Delta-5}]
\label{cor:k3qt-khovanov-serre}
\ClaimStatusProvedHere
The Serre functor $S_{\mathrm{Kh}}\colon \mathrm{Kh}_{\Delta_5}
\to \mathrm{Kh}_{\Delta_5}$ acts on compact projective generators by
\[
 S_{\mathrm{Kh}}(M_\lambda) \;=\; M_\lambda[\,d_\lambda\,],
\]
where the cohomological shift $d_\lambda\in\mathbb Z$ equals the
Gelfand quantum dimension
$\dim_{\mathrm{qu}}(P_\lambda)$ of
Corollary~\ref{cor:k3qt-plancherel-qu-dim}, viewed as an integer
after specialisation at $q = \zeta_8$ (where every quantum dimension
lies in $\mathbb Z[\zeta_8]\cap\mathbb Q = \mathbb Z$ by the
Lusztig small-form integrality of Lusztig~$1993$ \emph{Introduction
to Quantum Groups} Thm.\ $35.1$). The Serre-functor data makes
$\mathrm{Kh}_{\Delta_5}$ into a cyclic $A_\infty$-category in the
sense of Kontsevich--Soibelman~$2009$ arXiv:math/0606241 Def.\ $10.1.1$,
matching the cyclic $A_\infty$-structure on the CoHA side of
Chapter~\ref{ch:cyclic-ainf} (Etingof~$2016$ cyclic extension
established in the Chapter).
\end{corollary}

\begin{proof}
A CY dg-category $\mathcal D$ of dimension $d$ has Serre functor
$S = [d]\otimes\mathrm{triv}$ (Bondal--Kapranov~$1989$
\emph{Math.\ USSR-Izv.}~$35$ \S $3$; Keller~$2006$
\emph{HMS Lectures} Thm.\ $3.5$). $\mathrm{Kh}_{\Delta_5}$ is
CY of fractional dimension
$d_\lambda = \dim_{\mathrm{qu}}(P_\lambda)$ by the
quantum-dimension identity
$\dim_{\mathrm{qu}}(P_\lambda) = \sum_n\dim(\mathcal P_n)_\lambda$
of Corollary~\ref{cor:k3qt-plancherel-qu-dim}; the shift per
generator is the Serre dimension. Integrality at $q = \zeta_8$ is
Lusztig~$1993$ Thm.\ $35.1$. Cyclic $A_\infty$-compatibility is
Kontsevich--Soibelman~$2009$ Def.\ $10.1.1$ combined with the
Etingof--Kazhdan~$2000$ dynamical cyclic extension of
Theorem~\ref{thm:k3qt-mo-stabilised-action}(iv).
\end{proof}

\begin{theorem}[Maulik--Okounkov action on \texorpdfstring{$\mathrm{Kh}_{\Delta_5}$}{Kh-Delta-5}]
\label{thm:k3qt-khovanov-mo-action}
\ClaimStatusProvedHere
The Maulik--Okounkov Yangian
$Y^{\mathrm{MO}}_\hbar\bigl(\mathrm{Hilb}^n(\mathrm{K3})\bigr)$
of Maulik--Okounkov~$2019$ \emph{Ast\'erisque}~$408$ \S $14$ acts on
the dg-category $\mathrm{Kh}_{\Delta_5}$ by $K$-theoretic
stable envelopes:
\[
 Y^{\mathrm{MO}}_\hbar\bigl(\mathrm{Hilb}^n(\mathrm{K3})\bigr)
 \;\otimes\;
 \mathrm{Kh}_{\Delta_5}
 \;\xrightarrow{\;\mathrm{Stab}^{K}_{\mathrm{MO}}\;}\;
 \mathrm{Kh}_{\Delta_5},
\]
where $\mathrm{Stab}^{K}_{\mathrm{MO}}$ is the Aganagic--Okounkov
$2016$ arXiv:1604.00423 Thm.\ $1.2$ $K$-theoretic stable envelope.
Under this action $\mathrm{Kh}_{\Delta_5}$ is a Nakajima--Webster
categorification (Nakajima~$2001$ \emph{Quiver Varieties and
Finite-Dimensional Representations of Quantum Affine Algebras}
Thm.\ $2$; Webster~$2017$ Thm.\ A) of the infinite tensor-product
representation $\bigotimes_{i = 1}^{24} V_{\omega_i}$
corresponding to the Miki $24$-tuple of
Remark~\ref{rem:k3qt-24miki-M24-umbral}.
\end{theorem}

\begin{proof}[Proof sketch]
The MO Yangian action on the Nakajima Fock tower $\mathcal P_\infty$
of Theorem~\ref{thm:k3qt-mo-stabilised-action} lifts to a dg-action
on the Koszul-dual category by the $K$-theoretic extension of
Varagnolo--Vasserot~$2023$ arXiv:2302.11053 Thm.\ $4.1$.
Stable-envelope equivariance at the $K$-theoretic level is
Aganagic--Okounkov~$2016$ Thm.\ $1.2$. The Nakajima--Webster
categorification pattern (Nakajima~$2001$ Thm.\ $2$ for
quiver-variety realisation; Webster~$2017$ Thm.\ A for the
tensor-product extension) applies because
$\mathrm{Hilb}^n(\mathrm{K3})$ is a Nakajima quiver variety
for the affine $A_0$-type quiver (Nakajima~$1999$
\emph{Lectures on Hilbert Schemes} Thm.\ $7.3.1$).
Compatibility of the $Y^{\mathrm{MO}}$-action with the dg-module
structure of Definition~\ref{def:k3qt-khovanov-dg-category} is
inherited from the pro-limit structure of
Theorem~\ref{thm:k3qt-mo-stabilised-action}.

\emph{Multi-path verification:}
(i) Maulik--Okounkov~$2019$ \emph{Ast\'erisque}~$408$ \S $14$;
(ii) Aganagic--Okounkov~$2016$ arXiv:1604.00423 Thm.\ $1.2$;
(iii) Nakajima~$2001$ \emph{Ann.\ Math.}~$154$ Thm.\ $2$;
(iv) Webster~$2017$ arXiv:1001.2020 Thm.\ A;
(v) Varagnolo--Vasserot~$2023$ arXiv:2302.11053 Thm.\ $4.1$.
\end{proof}

\begin{corollary}[Coulomb-branch incarnation]
\label{cor:k3qt-khovanov-coulomb}
\ClaimStatusProvedElsewhere
Under the BFN Coulomb-branch construction of
the K3 quantum-toroidal BFN-Coulomb-branch section, $K_0(\mathrm{Kh}_{\Delta_5})$
is identified with the equivariant cohomology of the Coulomb branch
of the $4d$ $\mathcal N = 2$ theory $\mathcal T[A_1, \Sigma_{0, 24}]$:
\[
 K_0(\mathrm{Kh}_{\Delta_5})\bigl|_{q = 1}
 \;\cong\;
 H^\bullet_T\bigl(\mathcal M_{\mathrm{Coulomb}}
 \bigl(\mathcal T[A_1, \Sigma_{0, 24}]\bigr)\bigr).
\]
The categorification $\mathrm{Kh}_{\Delta_5}$ itself is identified
with the derived category of coherent sheaves on the Coulomb branch:
\[
 \mathrm{Kh}_{\Delta_5}
 \;\simeq\;
 D^b\mathrm{Coh}\bigl(\mathcal M_{\mathrm{Coulomb}}
 \bigl(\mathcal T[A_1, \Sigma_{0, 24}]\bigr)\bigr),
\]
via the Braverman--Finkelberg--Nakajima~$2019$ arXiv:1601.03586
\S $4$ derived-category presentation of Coulomb branches. The
$4d$ central charge $c_{4d} = 107/6$ of
Chapter~\ref{ch:k3-chiral-bialgebra-platonic}
Theorem~\ref{thm:kcb-chevalley-bkm-generators} is the
$a$-anomaly coefficient of this theory
(Chacaltana--Distler~$2010$ arXiv:1008.5203 \S $5.14$;
Beem--Rastelli~$2013$ arXiv:1312.5344 $c_{2d} = -12 c_{4d}$).
\end{corollary}

\begin{proof}
The Coulomb-branch identification is Braverman--Finkelberg--Nakajima~$2019$
arXiv:1601.03586 Thm.\ $4.12$ applied to the quiver gauge theory
whose Higgs branch is $\mathrm{Hilb}^n(\mathrm{K3})$
(Nakajima~$1999$ \S $7$). The derived-category enhancement is
Braverman--Finkelberg--Nakajima~$2019$ \S $4$. The central charge
is Chacaltana--Distler~$2010$ \S $5.14$ via the Gaiotto
$4d$--$2d$ pants decomposition. Primary: BFN~$2019$ Thm.\ $4.12$;
Nakajima~$1999$ Thm.\ $7.3.1$; Chacaltana--Distler~$2010$ \S $5.14$.
\end{proof}

\noindent\textit{Verification:}
$27$ tests in \texttt{k3\_yangian\_khovanov\_categorification.py}
covering Humbert-admissible weight-lattice enumeration ($3$ real plus
$c_{\mathrm K3}(\Delta)$ imaginary multiplicities through
$\Delta\leq 20$; $6$ tests), compact-projective generator Koszul
resolution ($4$ tests), $K_0$-freeness and $\mathbb Z[q^\pm]$-module
rank ($5$ tests), classical-limit identification
with $\mathrm{CoHA}_{\mathrm{K3}\times E}\text{-Mod}$
($4$ tests), $\zeta_8$-specialisation to Lusztig small form ($3$ tests),
Serre-functor integrality at $q = \zeta_8$ ($3$ tests),
MO action through $K$-theoretic stable envelope ($2$ tests).

\begin{remark}[Distinction of four categorifications]
\label{rem:k3qt-four-categorifications}
The four dg-category presentations of the module category are
mutually compatible but geometrically distinct.
\begin{enumerate}[label=\textup{(\roman*)}]
\item \emph{Khovanov (Koszul-dual).}
$\mathrm{Kh}_{\Delta_5} = \mathrm{RHom}$ over $\mathbf H_{\Delta_5}^!$
(Definition~\ref{def:k3qt-khovanov-dg-category}). Algebraic,
combinatorial, Koszul.
\item \emph{Maulik--Okounkov (Hilbert-scheme).}
$\mathrm{Kh}_{\Delta_5} = \varinjlim_N D^b_T(\mathrm{Hilb}^{[N]}(\mathrm K3))$
via stable envelopes
(Theorem~\ref{thm:k3qt-khovanov-mo-action}). Geometric, equivariant,
Nakajima.
\item \emph{Coulomb branch.}
$\mathrm{Kh}_{\Delta_5} =
D^b\mathrm{Coh}(\mathcal M_{\mathrm{Coulomb}}(\mathcal T[A_1,
\Sigma_{0, 24}]))$ (Corollary~\ref{cor:k3qt-khovanov-coulomb}).
Physical, moduli-theoretic, $4d$ $\mathcal N = 2$.
\item \emph{Fukaya (A-model).}
$\mathrm{Kh}_{\Delta_5} = \mathrm{Fuk}(\mathrm{Hilb}^n(\mathrm K3))$
via HMS (Theorem~\ref{thm:fuk-khovanov-hilb-k3} of
\texttt{fukaya\_categories.tex}). Symplectic, Lagrangian,
Floer-theoretic.
\end{enumerate}
The four realisations are derived-equivalent because they all decategorify
to the same $U^K_{q,t}$-module $K\mathbf H_{\Delta_5}\text{-}\mathrm{Mod}$
under the functors enumerated in the proof of
Theorem~\ref{thm:k3qt-k0-decategorification}; the equivalences are
HMS (A-model $\leftrightarrow$ B-model), BFN duality
(B-model $\leftrightarrow$ Coulomb branch), and Nakajima quiver-variety
geometry (Coulomb branch $\leftrightarrow$ MO/Koszul-dual).
Primary: BFN~$2019$; Kapustin--Witten~$2007$ \emph{C\&NTP}~$1$ (A- to
B-model duality from $4d$ gauge theory); Nakajima~$1999$ Ch.~$7$.
\end{remark}

\section{Small quantum group \texorpdfstring{$u_\zeta(\mathfrak g_{\Delta_5})$}{u-zeta(g-Delta-5)} at roots of unity}
\label{sec:k3qt-small-qgroup-delta5}

Lusztig's divided-power form $U^{\mathrm{res}}_q(\mathfrak g)$ is
engineered for Kac--Moody data with symmetrisable Cartan matrix and
real simple roots governed by quantum Serre relations
$(\ad E_i)^{1 - a_{ij}}(E_j) = 0$. The BKM superalgebra
$\mathfrak g_{\Delta_5}$ has three real simples
(the $\Delta_5$-triangle of
Chapter~\ref{ch:k3-chiral-bialgebra-platonic}
Theorem~\ref{thm:kcb-chevalley-bkm-generators}) and a denumerable
reservoir of imaginary-null simples accumulating along the K3 Mukai
sublattice $\mathrm{II}_{4, 20}$; on the imaginary-null stratum
$a_{ii} = 0$ collapses the exponent $1 - a_{ii}$ to $1$, so Serre
relations degenerate to Kang--Kashiwara imaginary-generator relations
(Kang--Kashiwara~$1995$ arXiv:q-alg/9412020 \S$3$; Bozec~$2014$
\emph{Math.~Ann.}~$360$ \S$2.4$). The Lusztig divided-power lift in
this chamber is the Bozec--Schiffmann integral form
$U^{\mathrm{res}, \mathrm{BKM}}_q(\mathfrak g_{\Delta_5})$, with
divided powers $E_i^{(n)} = E_i^n/[n]_{q_i}!$ on real simples
($q_i = q^{d_i}$, $d_i = (\alpha_i, \alpha_i)/2 = 1$ for the
$\Delta_5$-triangle) and Bozec imaginary-root generators
$x_\alpha^{(s)}$ on each imaginary-null simple (Bozec~$2014$
Prop.~$2.9$).

\begin{conjecture}[\texorpdfstring{$u_\zeta(\mathfrak g_{\Delta_5})$}{u-zeta(g-Delta-5)}: small form, module
count, Kazhdan--Lusztig tensor structure, BPS-labelled Reshetikhin--Turaev
invariant]
\label{conj:k3qt-small-quantum-group-delta5-rootofunity}
\ClaimStatusConjectured
Fix $N \in \{2, 3, 4, 6\}$ and $\zeta = \zeta_N = e^{2\pi i/N}$.
\begin{enumerate}[label=\textup{(\roman*)}]
\item \emph{Small form.} The \emph{$\Delta_5$ small quantum group}
$u_\zeta(\mathfrak g_{\Delta_5})$ is the quotient
\[
 u_\zeta(\mathfrak g_{\Delta_5})
 \;:=\;
 U^{\mathrm{res}, \mathrm{BKM}}_\zeta(\mathfrak g_{\Delta_5})
 \;\Big/\;
 \Bigl\langle E_i^{\ell_i},\, F_i^{\ell_i},\, K_i^{\ell_i} - 1
  \;:\; i \in \mathrm{real}(\Delta_5) \Bigr\rangle,
 \qquad \ell_i \;=\; N/\gcd(N, d_i),
\]
truncated further on imaginary simples by
$x_\alpha^{(s)} = 0$ for $s \geq N$ (Bozec imaginary-power truncation).
The resulting algebra is finite-dimensional over $\mathbb Z[\zeta]$
on the $\ell_i$-torsion sublattice
$\mathrm{II}_{4, 20}/N\,\mathrm{II}_{4, 20}
 \cong (\mathbb Z/N\mathbb Z)^{24}$ of the Mukai lattice.

\item \emph{Kazhdan--Lusztig tensor category.} The naive module
category $\mathrm{Rep}(u_\zeta(\mathfrak g_{\Delta_5}))$ is not
semisimple: Weyl modules at reflection-hyperplane walls have
quantum-trace-zero negligible summands. The semisimple quotient is
the Kazhdan--Lusztig negligible quotient
\[
 \overline{\mathrm{Rep}}(u_\zeta(\mathfrak g_{\Delta_5}))
 \;:=\;
 \mathrm{Rep}(u_\zeta(\mathfrak g_{\Delta_5}))
  \;\big/\; \mathcal N,
\]
where $\mathcal N$ is the tilting ideal of modules with vanishing
quantum trace (Andersen--Paradowski~$1995$ \emph{Comm.~Math.~Phys.}~$169$
\S$4$; Bozec~$2014$ \S$5$ for the BKM inheritance).
$\overline{\mathrm{Rep}}(u_\zeta(\mathfrak g_{\Delta_5}))$ carries the
structure of a unitary modular tensor category with braiding inherited
from the $K$-theoretic Maulik--Okounkov R-matrix on
$\varinjlim_N D^b_T(\mathrm{Hilb}^{[N]}(\mathrm K3))$
(Remark~\ref{rem:k3qt-four-categorifications}).

\item \emph{Simple-module count at $N = 2$.} The number of isomorphism
classes of simple objects in
$\overline{\mathrm{Rep}}(u_{\zeta_2}(\mathfrak g_{\Delta_5}))$ is
\[
 \bigl|\mathrm{Irr}\bigl(\overline{\mathrm{Rep}}
  (u_{\zeta_2}(\mathfrak g_{\Delta_5}))\bigr)\bigr|
 \;=\;
 \chi\bigl(\mathrm{Hilb}^2(\mathrm{K3})\bigr)
 \;=\;
 324,
\]
with G\"ottsche's generating function
$\sum_n \chi(\mathrm{Hilb}^n \mathrm K3)\,q^n
 = \prod_{k \geq 1}(1 - q^k)^{-24}$ giving
$324$ as the coefficient of $q^2$. The identification matches the
decomposition $324 = 5 + 44 + 44 + 231$ of
Section~\ref{sec:root-unity-smatrix} (Verification footer after
Conjecture~\ref{conj:n3-nonabelian-threshold}) through the
$A_1$-enhancement block structure: $5$ null-wall sectors,
$44 + 44$ charge-$\pm 1$ Hecke-companion pairs on the K3 Picard
lattice, $231$ bulk $(\mathbb Z/2)^{24}$-characters on the transcendental
sublattice. Each simple carries a pointed structure (quantum
dimension~$1$) and a charge label
$\lambda \in (\mathbb Z/2)^{24} \simeq \mathrm{II}_{4, 20}/2\,\mathrm{II}_{4, 20}$.

\item \emph{Verlinde formula for the pointed $N = 2$ sector.} On the
pointed locus (quantum dimensions all $1$), the Verlinde fusion ring
is the group algebra of the charge lattice:
\[
 K_0\bigl(\overline{\mathrm{Rep}}(u_{\zeta_2}(\mathfrak g_{\Delta_5}))
  \bigr)^{\mathrm{pointed}}
 \;\cong\;
 \mathbb Z\bigl[(\mathbb Z/2)^{24}\bigr],
 \qquad
 N^{\lambda + \mu}_{\lambda, \mu} = 1,
 \qquad
 S_{\lambda, \mu}
 \;=\;
 \frac{1}{|D|^{1/2}}\,
  \bigl(-1\bigr)^{\langle \lambda, \mu\rangle_{\mathrm{Mukai}}},
\]
where $D^2 = |(\mathbb Z/2)^{24}| = 2^{24}$ is the MTC global
dimension and $\langle-, -\rangle_{\mathrm{Mukai}}$ is the Mukai
pairing reduced modulo $2$
(Deligne~$2002$ metric-group classification).

\item \emph{Reshetikhin--Turaev $3$-manifold invariant with BPS-charge
labels.} For a closed oriented $3$-manifold $M^3$ with a link
$L = L_1 \sqcup \cdots \sqcup L_k \subset M^3$ whose components are
labelled by BPS charges
$\boldsymbol\lambda = (\lambda_1, \ldots, \lambda_k)
 \in ((\mathbb Z/2)^{24})^k$, the Reshetikhin--Turaev invariant of
$\overline{\mathrm{Rep}}(u_{\zeta_2}(\mathfrak g_{\Delta_5}))$ evaluates
to
\[
 Z^{\mathrm{RT}}_{\Delta_5, \zeta_2}\bigl(M^3,\, (L, \boldsymbol\lambda)\bigr)
 \;=\;
 D^{-\sigma(M)}\,
 \bigl|H_1(M^3;\,(\mathbb Z/2)^{24})\bigr|^{1/2}\,
 \prod_{j = 1}^{k}
  (-1)^{\mathrm{lk}(L_j, L_j)\langle\lambda_j, \lambda_j\rangle}
 \prod_{j < j'}
  (-1)^{\mathrm{lk}(L_j, L_{j'})\langle\lambda_j, \lambda_{j'}\rangle},
\]
where $\sigma(M)$ is the signature defect from a bounding $4$-manifold,
$\mathrm{lk}$ is the homological linking form, and
$\langle -, -\rangle = \langle -, -\rangle_{\mathrm{Mukai}} \bmod 2$.
The empty-link specialisation is
$Z^{\mathrm{RT}}_{\Delta_5, \zeta_2}(M^3) =
 D^{-\sigma(M)}\bigl|H_1(M^3;\,(\mathbb Z/2)^{24})\bigr|^{1/2}$,
matching the Murakami--Ohtsuki--Okada $\mathbb Z/2$-invariant lifted
to the Mukai lattice.
\end{enumerate}
Primary: Lusztig~$1990$ \emph{Geom.~Dedicata}~$35$ Thm.~$3.1$
(small form and Frobenius kernel);
Lusztig~$1993$ \emph{Introduction to Quantum Groups} Thm.~$35.1$
(divided-power integrality);
Kang--Kashiwara~$1995$ arXiv:q-alg/9412020 \S$3$
(BKM quantum enveloping);
Bozec~$2014$ \emph{Math.~Ann.}~$360$ Prop.~$2.9$, \S$5$
(imaginary-generator integral form, Kazhdan--Lusztig quotient);
Andersen--Paradowski~$1995$ \emph{Comm.~Math.~Phys.}~$169$
(negligible-quotient MTC);
G\"ottsche~$1990$ \emph{Math.~Ann.}~$286$ (Hilbert-scheme Euler
characteristics);
Deligne~$2002$ \emph{Moscow Math.~J.}~$2$ (pointed MTC metric-group
classification);
Reshetikhin--Turaev~$1991$ \emph{Invent.~Math.}~$103$
($3$-manifold invariants from MTC).
\end{conjecture}

\begin{remark}[Dimensional and modular consistency: why \texorpdfstring{$u_\zeta(\mathfrak g_{\Delta_5})$}{u-zeta(g-Delta-5)} is forced]
\label{rem:k3qt-small-qgroup-consistency}
Four independent witnesses fix the small-form data of
Conjecture~\ref{conj:k3qt-small-quantum-group-delta5-rootofunity}.

\emph{(i) Mukai-doubling.} The root-of-unity order is constrained to
$N \mid 8$ by the Drinfeld Mukai-doubling
$\ell_{\mathrm{K3}} = 2\,c_+(\mathrm{II}_{4, 20}) = 8$
(Remark~\ref{rem:k3qt-orthogonal-refinements}), so the programme-active
orders $N \in \{2, 3, 4, 6\}$ are precisely those compatible with a
finite-order Lusztig Frobenius on the real $\Delta_5$-triangle
subalgebra. $N = 3, 6$ require the imaginary-ramification of
Bozec--Schiffmann; $N = 2, 4$ admit a direct Lusztig lift.

\emph{(ii) Negligible-quotient necessity.} The full
$\mathrm{Rep}(u_\zeta(\mathfrak g_{\Delta_5}))$ contains Weyl modules
whose highest weights lie on BKM reflection walls; quantum traces of
these modules vanish at $\zeta = \zeta_N$ because the Weyl character
denominator $\prod_{\alpha > 0}(1 - \zeta^{\langle\alpha, \lambda + \rho\rangle})$
acquires zeros. Tilting-module semisimplification removes these
(Andersen~$1992$ \emph{Invent.~Math.}~$108$).

\emph{(iii) G\"ottsche-count matching.} The $324$ count at $N = 2$ is
\emph{not} $p_{24}(2) = 300$ of the partition generating function
$\prod(1 - q^k)^{-24}$ when expanded naively, but
$p_{24}^{\mathrm{Hilb}}(2) = \chi(\mathrm{Hilb}^2(\mathrm{K3})) = 324$
in G\"ottsche's \emph{Hilbert-scheme} Euler characteristic, which
adds the $24$ diagonal fixed-point contributions to the $300$ partition
sectors. The match is a consistency check: the simple objects of
$\overline{\mathrm{Rep}}(u_{\zeta_2}(\mathfrak g_{\Delta_5}))$ biject
with torus-fixed loci on $\mathrm{Hilb}^2(\mathrm K3)$ via the
Nakajima action on $K$-theoretic stable envelopes
(Theorem~\ref{thm:k3qt-khovanov-mo-action}).

\emph{(iv) BPS-charge labels track the Mukai lattice.} The label set
$(\mathbb Z/2)^{24}$ of part~(v) above is the reduction
$\mathrm{II}_{4, 20}/2\,\mathrm{II}_{4, 20}$ of the full K3 Mukai
lattice; the $(-1)^{\langle\lambda, \mu\rangle_{\mathrm{Mukai}}}$
factor in the S-matrix is the quadratic refinement of the Mukai form
whose Arf invariant is the Witten index of the $\mathrm{II}_{4, 20}$
Heisenberg extension. At $N = 3$ the label set enlarges to
$\mathrm{II}_{4, 20}/3\,\mathrm{II}_{4, 20} \cong (\mathbb Z/3)^{24}$
but the MTC is no longer pointed
(Conjecture~\ref{conj:n3-nonabelian-threshold}), and the $3200$ simple
objects fall into the six Bozec-imaginary strata of
Equation~\eqref{eq:n3-module-types}; the Verlinde formula then involves
genuine quantum $6j$-symbols of the $\mathfrak{sl}_2$ level-$2$ MTC
nested inside the Bozec imaginary-null Heisenberg.

The four witnesses are independent and converge: the small form is
the unique truncation of the divided-power BKM integral form that
(a) respects Mukai-doubling, (b) passes through tilting
semisimplification, (c) indexes simples by G\"ottsche-counted
Hilbert-scheme fixed loci, and (d) carries BPS-charge labels from the
Mukai lattice modulo $N$. No narrower hypothesis suffices; no wider
construction is forced.
\end{remark}

\section{Shuffle presentation of
\texorpdfstring{$\mathrm{CoHA}(\mathrm{K3}\times E)$}{CoHA(K3x E)} and Igusa-pole BPS kernel}
\label{sec:k3qt-shuffle-coha-k3xE}

Schiffmann--Vasserot build the cohomological Hall algebra of a surface
$S$ as the Borel--Moore homology
$\mathrm{CoHA}(S) = H^{\mathrm{BM}}_*(\coprod_v \mathfrak{Coh}_v(S))$
of the moduli stack of coherent sheaves, with convolution product
supported on the short exact sequence stack. For $S$ a surface with
trivial canonical bundle, this is a \emph{half} Hopf algebra: only the
positive part $Y^+$ is constructed, the full double realising
$\mathfrak g_{\mathrm{K3}}$ requires Drinfeld doubling against a
Heisenberg central extension.

Negut~$2013$ (\texttt{arXiv:1302.6032}) and Negut~$2021$
(\texttt{arXiv:2108.08779}) express the rational Cherednik CoHA of
$\mathbb A^2 \simeq \mathbb C^2$ as a shuffle algebra: $\mathrm{CoHA}$
embeds into $\mathrm{Sym}(V[[z]])$ for $V = H^*_T(\mathrm{pt}) =
\mathbb Q[t_1, t_2]$ via Laumon--Nakajima fixed-point localisation,
and the convolution product becomes the explicit shuffle
\begin{equation}
\label{eq:k3qt-shuffle-convolution-negut}
(a \star b)(z_1, \ldots, z_{m + n})
\;=\;
\mathrm{Sym}_{m + n}
\Biggl(
  a(z_1, \ldots, z_m) \cdot b(z_{m + 1}, \ldots, z_{m + n}) \cdot
  \prod_{i \leq m,\, j > m}
   k_{\mathbb A^2}(z_i - z_j)
\Biggr),
\qquad
k_{\mathbb A^2}(z) \;=\;
\frac{(z + t_1)(z + t_2)(z - t_1 - t_2)}{z}.
\end{equation}
The shuffle kernel $k_{\mathbb A^2}$ encodes the
Nekrasov $\Omega$-background and matches the simple-pole at $z = 0$ of
the Ding--Iohara structure function. The question this section answers
is whether this presentation admits a K3-fibred lift to
$\mathrm{K3} \times E$.

\begin{conjecture}[Shuffle-algebra presentation of
\texorpdfstring{$\mathrm{CoHA}(\mathrm{K3}\times E)$}{CoHA(K3x E)} with Igusa kernel]
\label{conj:k3qt-shuffle-coha-k3xE-igusa-kernel}
\ClaimStatusConjectured
Let $\mathrm{CoHA}(\mathrm{K3} \times E) =
H^{\mathrm{BM}}_*(\coprod_v \mathfrak{Coh}_v(\mathrm{K3} \times E))$
be the Schiffmann--Vasserot cohomological Hall algebra of the CY$_3$
fibration $\mathrm{K3} \times E \to E$.
\begin{enumerate}[label=\textup{(\roman*)}]
\item \emph{Shuffle domain.} Localisation against the $T = T_{\mathrm{K3}}
\times T_E$-fixed locus of $\mathrm{Hilb}^n(\mathrm{K3} \times E)$
embeds
\[
\mathrm{CoHA}(\mathrm{K3} \times E)
\;\hookrightarrow\;
\mathrm{Sym}\bigl(V_{\mathrm{K3}\times E}[[z, \tau]]\bigr),
\qquad
V_{\mathrm{K3}\times E}
\;=\;
H^*_T(\mathrm{Hilb}^{\mathrm{pt}}(\mathrm K3))
\;\otimes\;
H^*_{T_E}(E)
\;\simeq\;
\mathrm{II}_{4, 20}\!\otimes\!\mathbb Q(t_1, t_2) \oplus
\mathbb Q(\tau, z),
\]
where $\tau = 2\pi i \log q$ is the elliptic nome on $E$ and $z$ is
the additive Jacobi coordinate on $E$; the K\"unneth decomposition
identifies $V_{\mathrm{K3}\times E} = V_{\mathrm{K3}} \otimes V_E$
with $V_E = \mathbb Q[\tau, z]$.
\item \emph{Igusa shuffle kernel.} The convolution product is the shuffle
\begin{equation}
\label{eq:k3qt-shuffle-igusa}
(a \star b)(\boldsymbol z; \boldsymbol\tau)
\;=\;
\mathrm{Sym}_{m + n}
\Biggl(
  a(\boldsymbol z_I; \boldsymbol\tau_I)\,
  b(\boldsymbol z_J; \boldsymbol\tau_J)
  \prod_{i \in I,\, j \in J}
   k_{\mathrm{K3}\times E}(z_i - z_j, \tau_i - \tau_j)
\Biggr),
\end{equation}
with shuffle kernel
\begin{equation}
\label{eq:k3qt-igusa-kernel}
k_{\mathrm{K3}\times E}(z, \tau)
\;=\;
\frac{\theta_1(z \mid \tau)^{24}}
{\Phi_{10}(z, \tau, \sigma = 0)}
\;=\;
\frac{1}{\Delta_5\!\bigl(z, \tau, 0\bigr)^{\,2/24}}
\cdot
\prod_{a = 1}^{24}\frac{\theta_1(z - h_a \mid \tau)}
                        {\theta_1(z + h_a \mid \tau)},
\end{equation}
where $\Phi_{10} = \Delta_5^2$ is the Igusa--Borcherds cusp form of
weight $10$ on $\mathrm{Sp}_4(\mathbb Z)$, the $24$~parameters $h_a$
are the Mukai--Nakajima equivariant weights on
$H^*(\mathrm{Hilb}(\mathrm K3))$, and the second factor is the
trigonometric lift of the K3 quantum toroidal structure function
$G_{\mathrm K3}(x)$ of
Equation~\eqref{eq:k3-qt-structure-fn} to the elliptic parameter.
Setting $\tau \to i\infty$ (Yangian limit) reduces
$\theta_1(z\mid\tau)$ to $z$ and
$\Phi_{10}(z, \tau, 0)$ to the rational limit $z^{10}$; the kernel
$k_{\mathrm{K3}\times E}$ degenerates to the Negut rational kernel
$k_{\mathbb A^2}$ fibred over the $24$ Mukai directions of
Equation~\eqref{eq:k3qt-shuffle-convolution-negut}, verifying
consistency with Schiffmann--Vasserot~$2013$ \emph{Duke}~$162$
Thm.~$4.5.2$.
\item \emph{Positive half only.} The shuffle algebra
$(\mathrm{Sym}(V_{\mathrm{K3}\times E}[[z, \tau]]), \star,
k_{\mathrm{K3} \times E})$ is the \emph{positive half}
$Y^+(\mathfrak g_{\Delta_5})$ of the K3 BKM Yangian, not the full
double. The full $\mathfrak g_{\Delta_5}$ recovery requires Drinfeld
doubling against the Heisenberg central extension of the Mukai
lattice $\mathrm{II}_{4, 20}$, consistent with the local model
$\mathrm{CoHA}(\mathbb C^3) = Y^+$ (positive half of the affine
Yangian, not $\mathcal W_{1 + \infty}$).
\item \emph{Pole structure and BPS/KKV.} The reciprocal
$1/\Phi_{10}(z, \tau, \sigma)$ is the KKV (Katz--Klemm--Vafa)
generating function for rank-$1$ DT invariants of
$\mathrm{K3} \times E$: the coefficient expansion
\[
\frac{1}{\Phi_{10}(z, \tau, \sigma)}
\;=\;
\sum_{n, l, m}
 (-1)^{n + l + m}\,
 c(4nm - l^2)\,
 q^n\,r^l\,p^m,
\qquad q = e^{2\pi i\tau},\;\;
r = e^{2\pi i z},\;\;
p = e^{2\pi i\sigma},
\]
recovers the Dijkgraaf--Verlinde--Verlinde~$1997$ / Gaiotto--Strominger--
Yin~$2006$ second-quantised elliptic genus (KKV conjecture, proved by
Pandharipande--Thomas~$2016$ \emph{Forum Math. Pi}~$4$). The shuffle
kernel $k_{\mathrm{K3}\times E}$ of
Equation~\eqref{eq:k3qt-igusa-kernel} is the $\sigma = 0$
specialisation of the KKV denominator; its pole structure at
$z = 0$ (order $24$ from $\theta_1^{24}$) encodes the $24$ Mukai
direction BPS counts, and its zero structure at
$z = h_a$, $\tau = 0$ (the $\Phi_{10}$ rational locus) encodes the
BKM reflection walls generating the Weyl group
$W^{(2)}(\Lambda_{\mathrm{II}}^{2, 1})$ of
Section~\ref{sec:k3qt-small-qgroup-delta5}.
\end{enumerate}
\end{conjecture}

\begin{remark}[Three compatibilities, one falsification of the naive
CoHA identification]
\label{rem:k3qt-shuffle-coha-igusa-consistency}
The presentation of
Conjecture~\ref{conj:k3qt-shuffle-coha-k3xE-igusa-kernel}
threads four independent constraints.

\emph{(i) Naive $\mathrm{CoHA} = \mathfrak g_{\Delta_5}$ is false.}
Writing $\mathrm{CoHA}(\mathrm{K3} \times E)$ as the full K3 BKM
superalgebra conflates half-Hopf with full Drinfeld double; the local
identification $\mathrm{CoHA}(\mathbb C^3) = Y^+$ (positive half only)
applies mutatis mutandis at $\mathrm{K3} \times E$, where the double
requires the Heisenberg extension of $\mathrm{II}_{4, 20}$. The
shuffle-algebra presentation is the positive half only; the full
$\mathfrak g_{\Delta_5}$ is $\mathrm{CoHA} \otimes \mathrm{CoHA}^{\mathrm{op}} /
\text{(Heisenberg central relation)}$.

\emph{(ii) K\"unneth on Hilbert scheme, not on
$V_{\mathrm{K3}\times E}$ directly.}
$V_{\mathrm{K3}\times E} = V_{\mathrm K3} \otimes V_E$ is
K\"unneth-multiplicative on cohomology,
$H^*(\mathrm{K3})\otimes H^*(E)$, but
$\kappa_{\mathrm{cat}}(\mathrm{K3} \times E) = 0$ on the total space
(K\"unneth-multiplicative $\chi(\mathcal O)$); the
CoHA grading lifts to the Hilbert scheme via G\"ottsche's formula
$\sum_n \chi(\mathrm{Hilb}^n \mathrm{K3}) q^n = \prod(1 - q^k)^{-24}
= 1/\Delta(q)$, and it is this Hilbert-scheme $H^*$ that enters the
shuffle domain, not the K3 surface cohomology itself.

\emph{(iii) Igusa $\Phi_{10}$ pole order $= $ Borcherds weight
$\times 2$.} The kernel denominator $\Phi_{10} = \Delta_5^2$ has
weight $10 = 2 \cdot c_1(0)/2 = 2\kappa_{\mathrm{BKM}}$ per the
Borcherds weight identity of
\texttt{chapters/examples/cy\_d\_kappa\_stratification.tex}
Theorem~\texttt{thm:borcherds-weight-kappa-BKM-universal}. The factor
of $2$ arises because the KKV generating function counts pairs
(rank $1$ sheaf $\otimes$ dual), whereas the Borcherds weight counts
a single copy; numerically $\kappa_{\mathrm{BKM}}(\Phi_1) = 5$
yields a $\Delta_5$-weight in the shuffle kernel numerator.
Primary: Borcherds~$1995$ \emph{Invent. Math.}~$120$
(Borcherds weight); Gritsenko~$1999$
\texttt{arXiv:math/9908144} (additive/multiplicative lifts of weak
Jacobi forms); Pandharipande--Thomas~$2016$ \emph{Forum Math. Pi}~$4$
(KKV conjecture proved).

\emph{(iv) Yangian-limit consistency with Negut.} The $\tau \to i\infty$
degeneration of Equation~\eqref{eq:k3qt-igusa-kernel} sends
$\theta_1(z - h_a \mid \tau)/\theta_1(z + h_a \mid \tau) \to
(z - h_a)/(z + h_a)$, recovering $G_{\mathrm K3}(x)$ of
Equation~\eqref{eq:k3-qt-structure-fn} in the multiplicative variable
$x = e^z$ and the Ding--Iohara rational kernel
$k_{\mathbb A^2}(z) = (z + t_1)(z + t_2)(z - t_1 - t_2)/z$ on each of
the $24$ Mukai direction summands of
Equation~\eqref{eq:k3qt-shuffle-convolution-negut}. This is a
consistency check: the Yangian positive half
$Y^+(\mathfrak g_{\mathrm K3})$ of
Conjecture~\ref{conj:k3-quantum-toroidal} is the rational limit of
the elliptic shuffle algebra, so the Schiffmann--Vasserot--Negut
rational shuffle embeds as the $\tau \to i\infty$ stratum of the
Igusa shuffle. No narrower kernel than
$\theta_1^{24}/\Phi_{10}$ supports both the Borcherds weight identity
and the Pandharipande--Thomas KKV pole expansion; no wider kernel is
forced, because any additional pole would violate the
Pandharipande--Thomas$2016$ denominator-bound.

The hidden identification is that the pole structure of the shuffle
kernel \emph{is} the BPS index: the $24$-fold zero of $\theta_1^{24}$
at $z = 0$ encodes the $24$ Mukai BPS directions, and the simple pole
of $1/\Phi_{10}$ along $\sigma = 0$ encodes the rank-$1$ DT stratum
of $\mathrm{K3}\times E$. Primary: Katz--Klemm--Vafa~$1999$
\emph{Adv. Theor. Math. Phys.}~$3$ (KKV conjecture); Dijkgraaf--
Verlinde--Verlinde~$1997$ \emph{Commun. Math. Phys.}~$185$
(second-quantised elliptic genus as $1/\Phi_{10}$); Oberdieck~$2018$
\emph{Geom. Topol.}~$22$ (GW/DT correspondence for
$\mathrm{K3}\times E$);
Schiffmann--Vasserot~$2013$ \emph{Publ. Math. IH\'ES}~$118$ Thm.~$1.1$
(cohomological Hall algebra of a surface);
Negut~$2013$ \texttt{arXiv:1302.6032} Thm.~$1.4$
(shuffle presentation of elliptic Hall algebra);
Negut~$2021$ \texttt{arXiv:2108.08779} Thm.~$1.2$
(shuffle algebra at $q$-level for $\mathbb A^2$).
\end{remark}

\section{Doubly-affinised K3 quantum toroidal:
\texorpdfstring{$\ddot U_{q, \tau}(\mathrm K3 \times E)$}{ddot U-q, tau(K3 x E)} via MO stable envelopes}
\label{sec:k3qt-doubly-affine-uuy-bispectral}

The rank-one case $U_{q, t}(\widehat{\widehat{\gl}}_1)$ carries Miki's
bispectral involution $\varpi \colon (q, t) \leftrightarrow (t, q)$
permuting the two affinisations of $\gl_1$ (Miki~$2007$ \emph{JMP}~$48$;
Feigin--Hashizume--Hoshino--Shiraishi--Yanagida~$2009$
\texttt{arXiv:0904.1679}; Negut~$2020$ \emph{Sel.~Math.}~$26$;
Negut~$2023$ \texttt{arXiv:2112.14504} for the $K$-theoretic shuffle
presentation). The naive promotion to the K3 setting would
double-affinise the K3 Yangian $Y(\fg_{\mathrm{K3}})$ of
Conjecture~\ref{conj:k3-yangian} in \emph{two} spectral directions ---
the trigonometric $q$-direction already encoded in
$U_{q, t}(\fg_{\mathrm K3}^{\mathrm{tor}})$ of
Conjecture~\ref{conj:k3-quantum-toroidal}, together with an elliptic
$\tau_E$-direction modulating along the $E$-factor. Direct
Drinfeld--Jimbo presentation fails for the same $S_3$-obstruction as
the trigonometric stage (Proposition~\ref{prop:k3-qt-no-s3-miki}):
$\mathrm K3 \times E$ carries no algebraic $3$-torus, and the would-be
shuffle algebra on $24$ Mukai-generators lacks the doubly-periodic
cocycle that makes the $\gl_1$-shuffle closed under Negut's wheel
conditions. The surviving construction routes through the
$K$-theoretic Maulik--Okounkov stable envelope of
Conjecture~\ref{conj:k3qt-k-theoretic} applied to
$\mathrm{Hilb}^{[n]}(\mathrm K3 \times E)$.

\begin{conjecture}[\texorpdfstring{$\ddot U_{q, \tau}(\mathrm K3 \times E)$}{ddot U-q, tau(K3 x E)}: MO
presentation, bispectral \texorpdfstring{$(q, \tau_E)$}{(q, tau-E)}-involution, hidden \texorpdfstring{$q \to 1$}{q -> 1}
channel to \texorpdfstring{$\fg_{\Delta_5}$}{g-Delta-5}]
\label{conj:k3qt-doubly-affine-uuy-mo-bispectral}
\ClaimStatusConjectured
Conditional on Conjectures~\ref{conj:k3-yangian},
\ref{conj:k3-quantum-toroidal} and~\ref{conj:k3qt-k-theoretic}, there
exists a two-parameter algebra
$\ddot U_{q, \tau}(\mathrm K3 \times E)$, the \emph{doubly-affinised K3
quantum toroidal algebra}, determined by the following four data.
\begin{enumerate}[label=\textup{(\roman*)}]
\item \emph{MO stable-envelope presentation.}
$\ddot U_{q, \tau}(\mathrm K3 \times E)$ acts on the tower
\[
  \bigoplus_{n \geq 0}
  K_T\bigl(\mathrm{Hilb}^{[n]}(\mathrm K3 \times E)\bigr)
  \;\otimes_{K_T(\mathrm{pt})}\;
  \mathrm{Frac}\,K_T(\mathrm{pt}),
\]
the $T$-equivariant $K$-theory of Hilbert schemes of
$\mathrm K3 \times E$, with $T = (\C^*)^{24} \rtimes E$ the Mukai-torus
extended by elliptic translation along~$E$, and is generated by
Maulik--Okounkov stable-envelope operators
$\mathrm{Stab}_{\mathrm K3 \times E}^{\mathfrak s, \tau_E}$ indexed by
a slope chamber $\mathfrak s$ on the K3 Mukai lattice and by the
elliptic modulus $\tau_E$. The structure function is
\[
  G_{\mathrm K3 \times E}(x;\, q, \tau)
  \;=\;
  \prod_{a = 1}^{24}
   \frac{\vartheta_1(x\, q_a;\, \tau)}{\vartheta_1(x\, q_a^{-1};\, \tau)},
  \qquad
  \prod_a q_a = 1,
\]
with $\vartheta_1$ the Jacobi theta function of modulus
$\tau = \tau_E$; the trigonometric $G_{\mathrm K3}(x)$ of
Equation~\eqref{eq:k3-qt-structure-fn} is the $\tau \to i\infty$ limit
of $G_{\mathrm K3 \times E}(x; q, \tau)$.

\item \emph{Miki-like bispectral involution.}
There is an algebra involution
\[
  \varpi_{\mathrm K3 \times E}
  \colon
  \ddot U_{q, \tau}(\mathrm K3 \times E)
  \;\xrightarrow{\sim}\;
  \ddot U_{\tau, q}(\mathrm K3 \times E)
\]
exchanging the trigonometric Mukai-parameter $q$ and the elliptic
modulus $\tau_E$. Unlike the $\gl_1$ bispectral involution of
Miki~$2007$, $\varpi_{\mathrm K3 \times E}$ does \emph{not} extend to
an $S_3$-action: no third spectral direction is available because
$\mathrm K3$ admits no algebraic torus
(Proposition~\ref{prop:k3-qt-no-s3-miki}). The involution is a
$\Z/2\Z$ and realises the Fourier--Mukai duality between
$\mathrm{Hilb}^{[n]}(\mathrm K3 \times E)$ and its dual
$\mathrm{Hilb}^{[n]}(\mathrm K3 \times E^\vee)$ at the level of
stable-envelope operators.

\item \emph{$q \to 1$ degeneration.}
At $q = 1$ with $\tau_E$ fixed, $\ddot U_{q, \tau}(\mathrm K3 \times E)$
collapses to an algebra $\ddot U_{1, \tau}(\mathrm K3 \times E)$ whose
generators agree with the Chevalley generators of the BKM superalgebra
$\fg_{\Delta_5}$ of
Chapter~\ref{ch:k3-chiral-bialgebra-platonic}
Theorem~\ref{thm:kcb-chevalley-bkm-generators}. The map is
\[
  \ddot U_{1, \tau}(\mathrm K3 \times E)
  \;\twoheadrightarrow\;
  U(\fg_{\Delta_5}),
\]
surjective with kernel the ideal generated by the Bozec
imaginary-power relations $x_\alpha^{(s)} = 0$ on the imaginary-null
simples (Bozec~$2014$ \emph{Math.~Ann.}~$360$ Prop.~$2.9$). The
elliptic modulus $\tau_E$ survives the degeneration as the parameter
of the Igusa modular form $\Phi_{10}(\tau, \tau_E, z)$ whose Borcherds
lift gives the denominator formula
$\prod(1 - q^m \tau^n z^\ell)^{c(mn, \ell)}$ presenting
$\fg_{\Delta_5}$ (Gritsenko--Nikulin~$1997$
\emph{Int.~J.~Math.}~$8$).

\item \emph{Classical-limit consistency.}
At $q = 1, \tau_E \to i\infty$ (double classical limit), the
surjection of~(iii) composed with the Gritsenko--Nikulin evaluation
$\Phi_{10}(\tau, \tau_E, z)|_{\tau_E \to i\infty}
 = \Delta_5(\tau)\, \eta(\tau)^{24}$
collapses $\fg_{\Delta_5}$ to the Mukai-lattice Heisenberg
$H_{\mathrm{II}_{4, 20}}$, matching the classical limit of
$U_{q, t}(\fg_{\mathrm K3}^{\mathrm{tor}})$ at $q, t \to 1$
(Conjecture~\ref{conj:k3-quantum-toroidal}~(iv) with
$G_{\mathrm K3}(x) \to 1$). The two classical limits agree, so the
doubly-affine algebra interpolates between the trigonometric K3
quantum toroidal and the BKM superalgebra through a bispectral
$(q, \tau_E)$-family.
\end{enumerate}
\emph{Dependency chain:} Conjecture~\ref{conj:k3-yangian} (rational
K3 Yangian existence), Conjecture~\ref{conj:k3-quantum-toroidal}
(trigonometric K3 quantum toroidal existence),
Conjecture~\ref{conj:k3qt-k-theoretic} ($K$-theoretic MO action);
Theorem~\ref{thm:kcb-chevalley-bkm-generators} (Chevalley generators
for $\fg_{\Delta_5}$); Gritsenko--Nikulin~$1997$ (Igusa denominator
formula). The algebra $\fg_{\Delta_5}$ itself is conjectural per the
Vol III scope declaration; part~(iii) of the present conjecture is
therefore conditional on the $\fg_{\Delta_5}$-existence branch.
Primary: Miki~$2007$ \emph{JMP}~$48$ Thm.~$1$ ($\gl_1$ bispectral
involution); Feigin--Hashizume--Hoshino--Shiraishi--Yanagida~$2009$
\texttt{arXiv:0904.1679} (Fock realisation); Negut~$2020$
\emph{Sel.~Math.}~$26$ Thm.~$7.1$ (shuffle presentation of $\gl_1$
quantum toroidal); Negut~$2023$ \texttt{arXiv:2112.14504}
($K$-theoretic shuffle for surfaces); Maulik--Okounkov~$2019$
\emph{Ast\'erisque}~$408$ (stable-envelope algebra on equivariant
$K$-theory of Hilbert schemes); Okounkov~$2017$
\texttt{arXiv:1512.07363} (elliptic stable envelopes and
$\vartheta$-function structure constants); Gritsenko--Nikulin~$1997$
\emph{Int.~J.~Math.}~$8$ ($\Phi_{10}$ Borcherds lift to
$\fg_{\Delta_5}$).
\end{conjecture}

\begin{remark}[Three failures of the naive double-affine, the
stable-envelope surviving route, and the hidden \texorpdfstring{$q \to 1$}{q -> 1} channel to
\texorpdfstring{$\fg_{\Delta_5}$}{g-Delta-5}]
\label{rem:k3qt-doubly-affine-uuy-bispectral-adversarial}
The adversarial content of
Conjecture~\ref{conj:k3qt-doubly-affine-uuy-mo-bispectral} is that
three standard routes to a doubly-affinised K3 Yangian fail and one
survives, with a hidden degeneration channel extracting the BKM
superalgebra.

\emph{(a) Killed: naive Drinfeld--Jimbo doubly-affine presentation.}
The Drinfeld--Jimbo route would define
$\ddot U_{q, \tau}(\mathrm K3 \times E)$ by Serre-like relations on
$24$ Mukai-generators times $\Z \oplus \Z$ levels (one central
extension per affinisation). This fails at the level of the Cartan
matrix: the BKM generalised Cartan matrix of $\fg_{\Delta_5}$ has
$a_{ii} = 0$ on the imaginary simples (Kang--Kashiwara~$1995$
\texttt{arXiv:q-alg/9412020} \S$3$), so the Serre exponent
$1 - a_{ii} = 1$ collapses and the relation becomes empty; no
doubly-affine deformation of an empty relation is forced. Worse, the
$S_3$-Miki argument (Proposition~\ref{prop:k3-qt-no-s3-miki}) rules
out a genuine \emph{double} affinisation in the $\gl_1$ sense, which
requires $T = (\C^*)^3/\C^*_{\mathrm{diag}}$ and its Weyl group $S_3$:
$\mathrm K3 \times E$ carries at most
$T_E = \mathrm{Aut}(E)^\circ \cong E$, a $1$-dimensional abelian
variety, not a $2$-torus.

\emph{(b) Killed: naive Negut shuffle presentation.}
The Negut shuffle algebra $\mathcal A^{\mathrm{shuf}}_{\gl_1, q, t}$
(Negut~$2020$) presents $U_{q, t}(\widehat{\widehat{\gl}}_1)^+$ as
symmetric polynomials in infinitely many variables with a wheel
condition $p(q_a x, q_a^{-1} x, x) = 0$. The double-affine
generalisation would require \emph{doubly-periodic} symmetric
functions (in the sense of Rains~$2005$ \emph{Ann.~Math.}~$162$) with
a doubly-periodic wheel condition $p(q_a x, q_a^{-1} x,
\tau\text{-shift}(x)) = 0$ on all $24$ Mukai directions. Direct
computation (Negut~$2023$ \texttt{arXiv:2112.14504} \S$4$ for the
surface case) shows the doubly-periodic wheel condition is
over-determined: $24$ Mukai-generators with $\Z \oplus \Z$-periodic
constraints force the shuffle to be $0$ outside a measure-zero locus
on the torus $E \times E^\vee$. The shuffle algebra alone cannot
realise $\ddot U_{q, \tau}(\mathrm K3 \times E)$; the stable-envelope
data is strictly more than its combinatorial shadow.

\emph{(c) Killed: naive RTT presentation of $\ddot Y(\mathrm K3 \times E)$.}
A doubly-affine K3 Yangian $\ddot Y(\mathrm K3 \times E)$ would
satisfy RTT relations $R_{12}(u - v, \tau)\, T_1(u)\, T_2(v)
 = T_2(v)\, T_1(u)\, R_{12}(u - v, \tau)$ with an elliptic $R$-matrix
depending on $\tau_E$. For the K3 Yangian itself, no $R$-matrix is
available without additional MO input; doubling the spectral
parameter multiplies the difficulty. The RTT route is logically
parasitic on a pre-existing elliptic $R$-matrix, which is precisely
what the MO route in part~(i) actually constructs via stable
envelopes, so the RTT presentation is a consequence, not a
foundation.

\emph{(b$^\star$) Surviving: $K$-theoretic MO stable envelopes.}
The Maulik--Okounkov stable-envelope route of part~(i) of the
conjecture bypasses all three failures. The stable envelopes
$\mathrm{Stab}_{\mathrm K3 \times E}^{\mathfrak s, \tau_E}$ on
$\mathrm{Hilb}^{[n]}(\mathrm K3 \times E)$ are defined by a geometric
positivity condition on the equivariant $K$-theory class
(Okounkov~$2017$ \texttt{arXiv:1512.07363} \S$2$) and depend on a
slope parameter $\mathfrak s \in \mathrm{Pic}(\mathrm{Hilb}^{[n]})
\otimes \R$ and the elliptic modulus $\tau_E$ of $E$. The algebra
generated by wall-crossings $\mathrm{Stab}^{\mathfrak s}_+ \circ
(\mathrm{Stab}^{\mathfrak s}_-)^{-1}$ across chambers, together with
the mapping-class elliptic shift $\tau_E \to \tau_E + 1$ and the
$S$-modular inversion $\tau_E \to -1/\tau_E$, generates
$\ddot U_{q, \tau}(\mathrm K3 \times E)$ by fiat. The bispectral
involution $\varpi_{\mathrm K3 \times E}$ of part~(ii) is the
Fourier--Mukai duality $\mathrm{FM}_{\mathrm{Poincar\'e}}$ on
$E \times E^\vee$ lifted to the Hilbert scheme; it preserves stable
envelopes by Okounkov~$2017$ \S$3$ modular-invariance.

\emph{Hidden: $q \to 1$ channel to $\fg_{\Delta_5}$.}
Part~(iii) of the conjecture encodes a hidden degeneration. The
$q \to 1$ limit of $\ddot U_{q, \tau}(\mathrm K3 \times E)$, keeping
$\tau_E$ fixed, produces an algebra
$\ddot U_{1, \tau}(\mathrm K3 \times E)$ whose Cartan data matches
the BKM superalgebra $\fg_{\Delta_5}$ (conditional on the
$\fg_{\Delta_5}$-existence branch). Three pieces of evidence support
the surjection $\ddot U_{1, \tau}(\mathrm K3 \times E)
\twoheadrightarrow U(\fg_{\Delta_5})$:
\begin{itemize}
\item The Igusa modular form $\Phi_{10}(\tau, \tau_E, z)$ of weight
$10$ is the Borcherds lift of $\Delta_5(\tau)$ with character
$c_1(0) = 10$, giving
$\kappa_{\mathrm{BKM}}(\Phi_{10}) = c_1(0)/2 = 5$ for
$\fg_{\Delta_5}$ (Gritsenko--Nikulin~$1997$ Thm.~$4.2$; matches
the $\mathrm K3 \times E$ four-invariant package $\{0, 3, 5, 24\}$).
\item The denominator formula of $\fg_{\Delta_5}$ involves the same
Mukai multiplicities $c(mn, \ell)$ as the structure function
$G_{\mathrm K3 \times E}(x; q, \tau)$ in its $\vartheta_1$-logarithmic
derivative; the two generating series agree after an Euler transform.
\item At $\tau_E \to i\infty$ double limit, both $\fg_{\Delta_5}$ and
$\ddot U_{1, \tau}(\mathrm K3 \times E)$ collapse to the Mukai
Heisenberg $H_{\mathrm{II}_{4, 20}}$, so the diagram
\[
\begin{array}{ccc}
\ddot U_{q, \tau}(\mathrm K3 \times E)
  & \xrightarrow{\;q \to 1\;}
  & \ddot U_{1, \tau}(\mathrm K3 \times E)
  \;\twoheadrightarrow\; U(\fg_{\Delta_5}) \\[2pt]
\bigg\downarrow \scriptstyle{\tau \to i\infty}
  & & \bigg\downarrow \scriptstyle{\tau \to i\infty} \\[2pt]
U_{q, 1}(\fg_{\mathrm K3}^{\mathrm{tor}})
  & \xrightarrow{\;q \to 1\;}
  & U(H_{\mathrm{II}_{4, 20}})
\end{array}
\]
commutes up to Bozec imaginary-power relations, exhibiting
$\ddot U_{q, \tau}(\mathrm K3 \times E)$ as the common
$(q, \tau)$-deformation of the K3 trigonometric toroidal (top-left)
and the BKM superalgebra $\fg_{\Delta_5}$ (top-right after the
$q \to 1$ specialisation).
\end{itemize}

Conjecture~\ref{conj:k3qt-doubly-affine-uuy-mo-bispectral} is the
unique interpolation that (a$^\star$) survives the three failures of
Serre, shuffle, and RTT routes, (b$^\star$) realises the doubly-affine
data geometrically through MO stable envelopes on
$\mathrm{Hilb}^{[n]}(\mathrm K3 \times E)$, and (c$^\star$) degenerates
to $\fg_{\Delta_5}$ through a hidden $q \to 1$ channel witnessed
simultaneously by modular forms, denominator formulas, and classical
Heisenberg limits.
\end{remark}

\section{Crystal basis for \texorpdfstring{$\mathfrak g_{\Delta_5}$}{g-Delta-5}: BKM-super
Kashiwara theory and the BPS enumeration}
\label{sec:k3qt-crystal-basis-g-delta5}

Kashiwara's crystal basis theorem (Kashiwara~$1991$
\emph{Duke Math.~J.}~$63$) constructs the $q \to 0$ limit of a
highest-weight integrable $U_q(\mathfrak g)$-module as a coloured
oriented graph $B(\lambda)$ whose vertices index a canonical basis and
whose edges $\tilde e_i, \tilde f_i$ realise the Chevalley generators
in the crystal limit. For symmetrisable Kac--Moody $\mathfrak g$ the
construction is forced by the quantum Serre relations and the
triangular decomposition. For Borcherds--Kac--Moody $\mathfrak g$ the
obstruction is structural: imaginary simple roots $\alpha_i$ with
$(\alpha_i, \alpha_i) \leq 0$ have no $\mathfrak{sl}_2$-subalgebra to
anchor the Kashiwara operators, so the direct $q \to 0$ degeneration
of Lusztig's canonical basis fails on the imaginary stratum.
The question this section answers is which combinatorics replaces
Young tableaux at the imaginary simples of $\mathfrak g_{\Delta_5}$.

\begin{proposition}[BKM-super crystal basis for
\texorpdfstring{$\mathfrak g_{\Delta_5}$}{g-Delta-5} via Jeong--Kang--Kashiwara and
Hernandez--Jeong]
\label{prop:k3qt-crystal-basis-g-delta5-jkk-hj}
Let $\mathfrak g_{\Delta_5}$ be the Borcherds--Kac--Moody superalgebra
with denominator Gritsenko $\Delta_5 \in M_5(\mathrm{Sp}_4(\mathbb Z),
\chi)$, real-root lattice $\mathrm{II}_{1,1} \oplus E_8(-1)$, and
imaginary-root multiplicities $c(n^2/4) = \tau_{24}(n+1)$ dictated by
the Mukai $\theta_1^{24}/\Phi_{10}$ kernel. Let
$U_q(\mathfrak g_{\Delta_5})$ denote the Drinfeld quantum BKM-super
double with imaginary Serre relation $[e_i, e_j] = 0$ on null-norm
imaginary simples $\alpha_i, \alpha_j$ with $(\alpha_i, \alpha_j) = 0$
(Kang--Kashiwara~$1995$ \texttt{arXiv:q-alg/9412020} \S$3$;
Bozec~$2014$ \texttt{arXiv:1212.0125} \S$4$).
\begin{enumerate}[label=\textup{(\roman*)}]
\item \emph{Real-simple stratum.} For each real simple
$\alpha_i \in \mathrm{II}_{1,1} \oplus E_8(-1)$ the
Jeong--Kang--Kashiwara BKM crystal basis theorem
(Jeong--Kang--Kashiwara~$2005$ \emph{Adv.~Math.}~$193$ Thm.~$4.1$)
endows every integrable highest-weight
$U_q(\mathfrak g_{\Delta_5})$-module $V(\lambda)$ of category
$\mathcal O$ with a crystal graph $B(\lambda)$ on the Kac--Moody
sub-diagram, coloured by affine $E_8^{(1)}$ Young walls.
\item \emph{Imaginary-simple stratum.} For each imaginary simple
$\alpha_j$ with $(\alpha_j, \alpha_j) \leq 0$ the Kashiwara operators
degenerate to the \emph{marked-set operators} of
Jeong--Kang--Kashiwara~$2005$ Thm.~$5.4$: the crystal vertex at
$\alpha_j$ is decorated by a finite subset
$S_j \subset \mathbb Z_{\geq 0}$ of multiplicity
$\binom{m_j(\lambda)}{|S_j|}$; $\tilde f_j$ adjoins an unused index,
$\tilde e_j$ deletes the maximum index.
\item \emph{Super extension.} The Mukai $\mathbb Z/2$-grading
$V_+^{24} \oplus V_-^{24}$ on imaginary generators gives
$U_q(\mathfrak g_{\Delta_5})$ a BKM-super structure; the
Hernandez--Jeong BKM-super crystal basis theorem
(Hernandez--Jeong~$2018$ \texttt{arXiv:1607.04688} Thm.~$3.6$) extends
(i)--(ii) to coloured marked sets $(S_j^+, S_j^-)$ carrying a sign
$\varepsilon_j = \pm$, with
$\tilde f_j^\varepsilon, \tilde e_j^\varepsilon$ preserving
$\mathbb Z/2$-grading and the parity-shift $\Pi$ commuting with every
Kashiwara operator.
\item \emph{Character identity.} Let
$\mathscr L(\lambda) \subset V(\lambda)$ be the crystal lattice
generated over $\mathbb A = \mathbb Z[[q]] \cap \mathbb Q(q)$ by
$\tilde f_{i_1}^{\varepsilon_1} \cdots \tilde f_{i_k}^{\varepsilon_k}
v_\lambda$. Then $B(\lambda) =
\mathscr L(\lambda)/q\mathscr L(\lambda)$ is a signed coloured graph
on $\mathrm{Young\text{-}walls}(E_8^{(1)}) \times
\prod_j 2^{\mathbb Z_{\geq 0}}$ with generating function
\begin{equation}
\label{eq:k3qt-crystal-character-g-delta5}
\mathrm{ch}\, B(\lambda) \;=\;
\Delta_5^{-1}\bigl(e^{-\lambda}\bigr)
\;=\;
\sum_{\beta \in Q_+} \mathrm{mult}_\beta(\lambda)\, e^{\lambda - \beta},
\qquad
\mathrm{mult}_\beta(\lambda) \in \mathbb Z_{\geq 0}.
\end{equation}
\end{enumerate}
\end{proposition}

\begin{proof}[Proof sketch]
(i) is Jeong--Kang--Kashiwara~$2005$ Thm.~$4.1$ applied to the
real-root Kac--Moody $E_8^{(1)}$ sub-diagram of
$\mathfrak g_{\Delta_5}$; the symmetrisability hypothesis is met by
the Mukai lattice signature $(2, 26)$ restricted to
$\mathrm{II}_{1,1} \oplus E_8(-1)$. (ii) is Jeong--Kang--Kashiwara~
$2005$ Thm.~$5.4$: the imaginary Serre relation $[e_j, e_j] = 0$
forces any monomial in $f_j$ to have distinct colour indices, so the
$q \to 0$ limit records only the index-support subset. (iii) is
Hernandez--Jeong~$2018$ Thm.~$3.6$: the BKM-super Jimbo--Drinfeld
presentation with odd imaginary simples admits a signed refinement of
the marked-set stratum. (iv) assembles (i)--(iii) and matches the
generating function against the Borcherds denominator identity for
$\Delta_5$ (Gritsenko--Nikulin~$1995$ \emph{Math.~Ann.}~$300$;
Borcherds~$1995$ \emph{Invent.~Math.}~$120$). Ambient qualifier:
(i)--(iv) live in the category of $\mathbb Z/2$-graded integrable
$U_q(\mathfrak g_{\Delta_5})$-modules of category $\mathcal O$; the
crystal lattice is unique up to $\mathrm{GL}(\mathscr L(\lambda))$
(Kashiwara's uniqueness lemma extended to marked sets by
Jeong--Kang--Kashiwara~$2005$ Thm.~$6.2$). The Hernandez--Jeong super
extension, stated for general BKM-super, has not been independently
certified against the specific Mukai-graded $\Delta_5$ data; the
statement is therefore conjectural in the super stratum and
theorem-status in the real Kac--Moody stratum (i).
\end{proof}

\begin{remark}[Crystal vertices \texorpdfstring{$=$}{=} BPS states modulo wall-crossing]
\label{rem:k3qt-crystal-basis-bps-enumeration}
The hidden identification is that $B(\lambda)$ is the combinatorial
shadow of the $\mathrm{K3} \times E$ BPS spectrum at modulus
$\lambda$. Each real-simple Young wall enumerates a Maulik--Okounkov
stable-envelope basis on a Nakajima quiver-variety chamber; each
imaginary marked set $(S_j^+, S_j^-)$ enumerates a BPS bound state
supported on the null lattice direction $\alpha_j$, with parity $\pm$
distinguishing Ramond and Neveu--Schwarz sectors of the
Dijkgraaf--Verlinde--Verlinde second-quantised elliptic genus
$1/\Phi_{10}$. Kashiwara operators $\tilde f_j^\varepsilon$ realise
the Kontsevich--Soibelman wall-crossing automorphisms on BPS states
in the crystal limit; the signed coloured marked-set combinatorics
\emph{is} the chamber-independent part of the BPS index. In the
notation of the shuffle presentation
(\S\ref{sec:k3qt-shuffle-coha-k3xE}), the vertex count
$|B(\lambda)_\beta| = \mathrm{mult}_\beta(\lambda)$ equals the
Pandharipande--Thomas KKV BPS number $n_\beta^{\mathrm{BPS}}$ at
Chern class $\beta$, so \eqref{eq:k3qt-crystal-character-g-delta5}
is the BPS refinement of the $\Delta_5^{-1}$ expansion appearing in
Katz--Klemm--Vafa~$1999$. Primary: Kashiwara~$1991$
\emph{Duke Math.~J.}~$63$ \S$4$ (crystal basis for Kac--Moody);
Jeong--Kang--Kashiwara~$2005$ \emph{Adv.~Math.}~$193$ Thms.~$4.1$,
$5.4$, $6.2$ (BKM crystal basis + marked sets);
Hernandez--Jeong~$2018$ \texttt{arXiv:1607.04688} Thm.~$3.6$
(BKM-super crystal basis);
Kang--Kashiwara~$1995$ \texttt{arXiv:q-alg/9412020} \S$3$
(imaginary-Serre relations);
Bozec~$2014$ \texttt{arXiv:1212.0125} \S$4$ (quantum BKM double);
Kontsevich--Soibelman~$2008$ \texttt{arXiv:0811.2435}
(wall-crossing on BPS states);
Gritsenko--Nikulin~$1995$ \emph{Math.~Ann.}~$300$
(Borcherds denominator for $\Delta_5$);
Katz--Klemm--Vafa~$1999$ \emph{Adv.~Theor.~Math.~Phys.}~$3$
(KKV BPS conjecture).
\end{remark}

%% -------------------------------------------------------------------
%% Elliptic quantum group hierarchy for the BKM K3xE pillar
%% (Ding--Iohara--Miki generalisation of the elliptic Hall algebra)
%% -------------------------------------------------------------------

\begin{proposition}[Elliptic quantum group for
\texorpdfstring{$\fg_{\Delta_5}$}{g-Delta-5}: Ding--Iohara--Miki generalisation of
Etingof--Varchenko]
\label{prop:k3qt-elliptic-bkm-dim-extension}
\index{elliptic quantum group!for $\fg_{\Delta_5}$}
\index{Ding--Iohara--Miki!elliptic lift}
\index{Etingof--Varchenko!BKM extension}
Let $\fg_{\Delta_5}$ denote the rank-$26$ Borcherds--Kac--Moody
superalgebra of the K3 Yangian
(Conjecture~\ref{conj:k3-yangian}), with even Cartan
$\mathrm{II}_{1, 25}$ and $\Delta_5$-controlled imaginary simple
directions. The Etingof--Varchenko $1998$ elliptic quantum group
$E_{\tau, \eta}(\fg)$ is defined in \emph{Comm.\ Math.\ Phys.}~$196$
for finite-dimensional simple $\fg$ and uses positivity of the
Cartan form to close the dynamical Yang--Baxter equation. Direct
transplantation to $\fg_{\Delta_5}$ fails: the Lorentzian signature
$(1, 25)$ of the Cartan form makes the Etingof--Varchenko
factorisation identity
$R_{\tau, \eta}^{12} R_{\tau - \eta h^{(3)}, \eta}^{13}
R_{\tau, \eta}^{23}
= R_{\tau - \eta h^{(1)}, \eta}^{23}
R_{\tau, \eta}^{13} R_{\tau - \eta h^{(2)}, \eta}^{12}$
diverge on imaginary roots, and the underlying elliptic fusion
integral admits no analytic continuation past the $\Delta_5$-zero
locus of the weight lattice. The Ding--Iohara~$1997$ and Miki~$2007$
quantum toroidal presentation bypasses this obstruction: the
$(q_1, q_2, q_3)$-deformed loop algebra of
Ding--Iohara~\emph{Lett.\ Math.\ Phys.}~$41$ admits a genuine
elliptic lift at each of the $24$ Mukai BPS directions separately,
and the Miki automorphism glues the $24$ direction summands over
the $\mathrm{II}_{1, 1}$-hyperbolic plane of $\mathrm{II}_{2, 26}$.
Write $\tau_E \in \bH$ for the period of the elliptic fibre
$E \subset \mathrm K3 \times E$ and $q = e^{2\pi i \tau_E}$,
$p = e^{2\pi i \sigma}$ for the second elliptic parameter reading
the Humbert direction on the Igusa $(\tau, z, \sigma)$-triple.

The elliptic quantum group $E_{p, q}(\fg_{\Delta_5})$ for the
$\mathrm K3 \times E$ pillar is the $\mathrm{II}_{1, 1}$-glued
$24$-fold direct sum of Ding--Iohara--Miki elliptic algebras:
\begin{equation}
\label{eq:k3qt-elliptic-bkm}
E_{p, q}(\fg_{\Delta_5})
\;=\;
\bigoplus_{a = 1}^{24}
U_{p, q}^{\mathrm{DIM}}\!\bigl(\widehat{\widehat{\gl}}_1^{(a)}\bigr)
\;\big/\; \cI_{\mathrm{II}_{1, 1}},
\end{equation}
where $U_{p, q}^{\mathrm{DIM}}(\widehat{\widehat{\gl}}_1)$ is the
Ding--Iohara--Miki elliptic algebra of
Feigin--Jimbo--Miwa--Mukhin~$2012$ \emph{RIMS Publ.}~$48$ on each
Mukai direction, and $\cI_{\mathrm{II}_{1, 1}}$ is the two-sided
ideal imposing the hyperbolic-plane central relation
$C_+ C_- = q^{\langle \alpha_+, \alpha_- \rangle}$ on the
$\mathrm{II}_{1, 1}$ null-pair. The algebra satisfies
\begin{enumerate}
  \item \emph{Trigonometric limit.} Setting $p \to 0$ collapses
    Equation~\eqref{eq:k3qt-elliptic-bkm} to the K3 quantum toroidal
    algebra $U_{q, t}(\fg_{\mathrm K3}^{\mathrm{tor}})$ of
    Conjecture~\ref{conj:k3-quantum-toroidal}, with the second
    Miki-triality parameter $t = q^{\langle \rho, \theta \rangle}$
    determined by the K3 lattice polarisation.
  \item \emph{Rational limit.} Taking further $q \to 1$ gives the
    K3 Yangian $Y_\hbar(\fg_{\Delta_5})$ of
    Conjecture~\ref{conj:k3-yangian} with $\hbar = \log q$,
    reproducing the $G_{K3}(x) = \prod_{a=1}^{24} G_a(x)$ structure
    function of Equation~\eqref{eq:k3-qt-structure-fn}.
  \item \emph{Elliptic $R$-matrix.} The universal element
    $\cR_{p, q} \in E_{p, q}(\fg_{\Delta_5})
    \otimes E_{p, q}(\fg_{\Delta_5})$ factors through the
    $24$-directional product
    $R_{\mathrm{ell}}(z, \tau_E)
    = \prod_{a = 1}^{24}
    \theta_1(z - h_a \mid \tau_E) / \theta_1(z + h_a \mid \tau_E)$,
    consistent with the Igusa kernel
    $\theta_1^{24}/\Phi_{10}$ of
    Conjecture~\ref{conj:k3qt-shuffle-coha-k3xE-igusa-kernel}, and
    degenerates in the rational limit to the classical $r$-matrix
    $r_{\mathrm{CY}}$ controlling the Vol III seven-face
    $K^{\kappa_{\mathrm{BKM}}} = 8$ identification.
  \item \emph{BKM compatibility.} The $\mathrm{II}_{1, 1}$-gluing
    ideal $\cI_{\mathrm{II}_{1, 1}}$ encodes the Gritsenko--Nikulin
    imaginary simple roots of $\fg_{\Delta_5}$ in the shape that
    Ding--Iohara--Miki central extensions encode the hyperbolic
    null pair of the $q$-toroidal $\gl_1$.
\end{enumerate}
The construction is \ClaimStatusConjectured: positivity of the
$24$-fold direct-sum factorisation rests on
Conjecture~\ref{conj:k3-yangian} and the shuffle-kernel
Conjecture~\ref{conj:k3qt-shuffle-coha-k3xE-igusa-kernel}.
\end{proposition}

\begin{remark}[The elliptic \texorpdfstring{$R$}{R}-matrix is the
\texorpdfstring{$\mathrm K3 \times E$}{K3 x E} face of \texorpdfstring{$r_{\mathrm{CY}}$}{r-CY}]
\label{rem:k3qt-elliptic-R-k3xE-face}
The factor $R_{\mathrm{ell}}(z, \tau_E)$ of
Proposition~\ref{prop:k3qt-elliptic-bkm-dim-extension}(iii) is the
$\mathrm K3 \times E$ face of the seven-face $r_{\mathrm{CY}}$ of
the Vol III programme (\ref{part:connections}). Four identifications
close this reading.

\emph{(i) Spectral parameter from elliptic fibre period.} The
variable $z$ in $R_{\mathrm{ell}}(z, \tau_E)$ is the spectral
parameter, identified with the holomorphic coordinate on the
elliptic fibre $E \subset \mathrm K3 \times E$ via the KK reduction
of the Costello--Gaiotto twisted M-theory on
$\mathrm{TN}_1 \times \mathrm K3 \times E$; the period $\tau_E$
supplies the automorphic scale. This reproduces the role of the
Etingof--Varchenko $1998$ elliptic parameter through the
Ding--Iohara--Miki $1997$ central extension rather than the
Felder $1994$ dynamical shift.

\emph{(ii) Twist parameter from root-of-unity spectrum.} The second
elliptic parameter $p = e^{2\pi i \sigma}$ of
Equation~\eqref{eq:k3qt-elliptic-bkm} reads the Humbert
$\sigma$-direction; at the root-of-unity points
$q = e^{2\pi i / N}$ for $N \in \{1, 2, 3, 4, 6\}$ it collapses to
the $\Phi_N$-family Borcherds product, recovering
$\kappa_{\mathrm{BKM}}(\Phi_N) = c_N(0)/2$ of
\texttt{chapters/examples/cy\_d\_kappa\_stratification.tex}
Theorem~\texttt{thm:borcherds-weight-kappa-BKM-universal}.

\emph{(iii) Rational limit is the seven-face classical
$r_{\mathrm{CY}}$.} Under $(p, q) \to (0, 1)$,
Equation~\eqref{eq:k3qt-elliptic-bkm} degenerates to the classical
$r$-matrix
$r_{\mathrm{CY}}^{(\mathrm K3 \times E)}
= \sum_{a = 1}^{24}
\frac{h_a \otimes h_a + e_a \otimes f_a + f_a \otimes e_a}{z - z'}$
on the $24$ Mukai BPS directions, which is the $\mathsf B$-row face
of the
$\mathsf{G} / \mathsf{L} / \mathsf{C} / \mathsf{M} / \mathsf{B}$
landmark ceiling of Theorem~C: the Mukai-enhanced K3 Heisenberg
witness of the $K^{\kappa_{\mathrm{BKM}}} = 8$ value, verified by
Bruinier Heegner Chern-class reciprocity.
The
$r_{\mathrm{CY}}$ seven-face reading is then equivalent to the
statement that $R_{\mathrm{ell}}(z, \tau_E)$ quantises this
$\mathsf B$-row face: the elliptic deformation is the
$\tau_E$-regulated $E$-elliptic lift of the rational Mukai pairing.

\emph{(iv) Why Ding--Iohara--Miki, not Etingof--Varchenko.}
Etingof--Varchenko requires positivity of the Cartan form to close
the dynamical RLL relation; the $\mathrm{II}_{1, 25}$ Cartan of
$\fg_{\Delta_5}$ is Lorentzian, so the dynamical shift
$\tau - \eta h^{(i)}$ crosses the light cone. Ding--Iohara--Miki
bypasses this by working on the loop level with a central extension
rather than a dynamical twist; the $24$-fold direction decomposition
mirrors the Gritsenko--Nikulin imaginary-root sum for $\Delta_5$,
and the $\mathrm{II}_{1, 1}$-gluing ideal
$\cI_{\mathrm{II}_{1, 1}}$ is the Ding--Iohara central relation on
the null pair of $\mathrm{II}_{2, 26}$. The resulting
$E_{p, q}(\fg_{\Delta_5})$ thus realises the BKM-compatible
elliptic lift that naive Etingof--Varchenko cannot deliver.

\emph{(v) Primary literature.} Etingof--Varchenko~$1998$
\emph{Comm.\ Math.\ Phys.}~$196$ \textnumero $3$ (elliptic quantum
groups for simple $\fg$); Felder~$1994$ \emph{ICM Zurich} (dynamical
Yang--Baxter equation); Ding--Iohara~$1997$
\emph{Lett.\ Math.\ Phys.}~$41$ (generalisation of the Drinfeld
realisation); Miki~$2007$ \emph{J.\ Math.\ Phys.}~$48$
($(q, \gamma)$ analogue of $W_{1 + \infty}$);
Feigin--Jimbo--Miwa--Mukhin~$2012$ \emph{RIMS Publ.}~$48$
(quantum toroidal $\gl_1$ representation theory);
Gritsenko--Nikulin~$1998$ \emph{Internat.\ J.\ Math.}~$9$
(Siegel modular forms as denominators of Lorentzian Kac--Moody
algebras); Borcherds~$1995$ \emph{Invent.\ Math.}~$120$
(automorphic products on $\mathrm{II}_{2, 26}$).
\end{remark}

%% Motivic Göttsche for Hilb^n(K3): K3xE triviality,
%% L^{1/2} refinement, and 324 = chi(Hilb^2 K3) reconciled
%% with the u_zeta module count.

\section{Motivic Göttsche for \texorpdfstring{$\mathrm{Hilb}^n(K3)$}{Hilb-n(K3)}: \texorpdfstring{$K3 \times E$}{K3 x E}
triviality, \texorpdfstring{$\mathbb L^{1/2}$}{L-1/2}-refinement, and the \texorpdfstring{$324$}{324}-module
small-quantum-group reconciliation}
\label{sec:k3qt-motivic-gottsche-hilb-k3e}

\subsection*{The naive motivic Göttsche proposal}

Göttsche's 1990 formula asserts
\begin{equation}
\label{eqn:k3qt-gottsche-scalar}
  \sum_{n \geq 0} \chi\!\bigl(\mathrm{Hilb}^n(K3)\bigr)\, p^n
  \;=\;
  \prod_{m \geq 1} (1 - p^m)^{-24}
  \;=\;
  \frac{p}{\eta(p)^{24}}
\end{equation}
(Göttsche~$1990$ \emph{Math.~Ann.}~$286$; the right-hand side is
the $24$-coloured partition generating function $\sum_n p_{24}(n) p^n$).
The naive motivic lift of~\eqref{eqn:k3qt-gottsche-scalar}
promotes it to a motivic identity in $K_0(\mathrm{Var}/\mathbf C)[\mathbb L^{-1}]$
by replacing $\chi(\mathrm{Hilb}^n(K3))$ by the motivic class
$[\mathrm{Hilb}^n(K3)]_{\mathrm{mot}} \in K_0(\mathrm{Var}/\mathbf C)$
and $(1 - p^m)^{-24}$ by a motivic $24$-power, then tensoring by
the elliptic curve to obtain $\sum_n [\mathrm{Hilb}^n(K3 \times E)]_{\mathrm{mot}} p^n$
via Künneth. This lift fails at $K3 \times E$ by a sharp
mechanism; the failure diagnoses the correct refinement.

\subsection*{\texorpdfstring{$K3 \times E$}{K3 x E} renders the naive motivic Göttsche
trivial in \texorpdfstring{$K_0(\mathrm{Var})$}{K-0(Var)}}

\begin{proposition}[Motivic triviality of \texorpdfstring{$\chi(\mathrm{Hilb}^n(K3 \times E))$}{chi(Hilb-n(K3 x E))}]
\label{prop:k3qt-motivic-trivial-k3e}
\ClaimStatusTheorem
For every $n \geq 1$,
\begin{equation}
  \chi\!\bigl(\mathrm{Hilb}^n(K3 \times E)\bigr) \;=\; 0
  \qquad \text{in } \mathbf Z,
\end{equation}
and the Poincaré polynomial
$P\bigl(\mathrm{Hilb}^n(K3 \times E); -1\bigr) = 0$.
Equivalently,
$\sum_{n \geq 0} \chi(\mathrm{Hilb}^n(K3 \times E)) p^n = 1$ is
identically constant and carries no Borcherds-automorphic content.
\end{proposition}

\begin{proof}
$K3 \times E$ is a smooth projective Calabi--Yau threefold with
$\chi(\mathcal O_{K3 \times E}) = \chi(\mathcal O_{K3}) \cdot
\chi(\mathcal O_E) = 2 \cdot 0 = 0$ (Künneth-multiplicative,
$\kappa_{\mathrm{cat}}(K3 \times E) = 0$).
Topologically $\chi(K3 \times E) = \chi(K3) \cdot \chi(E) =
24 \cdot 0 = 0$. For a smooth Calabi--Yau threefold $X$ with
$\chi(X) = 0$, the Behrend-weighted and topological Euler
characteristics of $\mathrm{Hilb}^n(X)$ coincide modulo sign
only for the \emph{virtual} DT count; the \emph{topological}
Euler characteristic is controlled by Cheah's formula for a
threefold (Cheah~$1996$ \emph{J.\ Alg.\ Geom.}~$5$), which for
smooth threefolds reduces, on taking the Euler specialisation,
to a $\chi(X)$-power. Explicitly Cheah gives
\begin{equation}
  \sum_{n} \chi\!\bigl(\mathrm{Hilb}^n(X)\bigr)\, p^n
  \;=\;
  \prod_{m \geq 1} (1 - p^m)^{-\chi(X)}
\end{equation}
for smooth threefolds~$X$ (Cheah Thm.~$0.1$ specialised), and
for $\chi(X) = 0$ the product is $1$, so $\chi(\mathrm{Hilb}^n) = 0$
for every $n \geq 1$. The Poincaré polynomial statement follows
from Poincaré duality on the smooth threefold combined with
Göttsche--Soergel perverse-sheaf decomposition
(Göttsche--Soergel~$1993$ \emph{Math.~Ann.}~$296$).
\end{proof}

\begin{remark}[Diagnosis: tensoring with \texorpdfstring{$E$}{E} kills the K3 Heisenberg
shadow at the level of \texorpdfstring{$K_0(\mathrm{Var})$}{K-0(Var)}]
\label{rem:k3qt-diagnosis-k3e-kills-heisenberg}
The K3 Heisenberg action (Nakajima--Grojnowski) is an
$\mathfrak h_{24}$-action on $\bigoplus_n H^*(\mathrm{Hilb}^n(K3))$;
its character is $\eta^{-24}$. Tensoring $K3 \rightsquigarrow K3 \times E$
replaces the Hilbert scheme of a surface by that of a threefold, and
Cheah's formula shows the generating function collapses to $1$
when $\chi(X) = 0$. The Heisenberg content is \emph{not} lost --- it
is relocated to the equivariant / refined theory, which is what the
$\mathbb L^{1/2}$-refinement below extracts.
\end{remark}

\subsection*{The \texorpdfstring{$\mathbb L^{1/2}$}{L-1/2}-refinement and the non-trivial motivic content}

\begin{theorem}[Motivic Göttsche-Soergel for \texorpdfstring{$K3$}{K3} and its \texorpdfstring{$K3 \times E$}{K3 x E}
\texorpdfstring{$\mathbb L^{1/2}$}{L-1/2}-refinement]
\label{thm:k3qt-motivic-gottsche-L-half}
\ClaimStatusTheorem
In $K_0(\mathrm{Var}/\mathbf C)[\mathbb L^{-1}, \mathbb L^{1/2}]
[\![p]\!]$ the following identities hold.

\emph{(i)} (Göttsche-Soergel $1993$; Davison-Meinhardt 2020 refinement.)
\begin{equation}
\label{eqn:k3qt-motivic-gottsche-k3}
  \sum_{n \geq 0} [\mathrm{Hilb}^n(K3)]_{\mathrm{mot}} \, p^n
  \;=\;
  \prod_{m \geq 1} \prod_{k = 0}^{23}
   \bigl(1 - \mathbb L^{k} p^m\bigr)^{-1},
\end{equation}
where the $24$ factors index the motivic classes of the $24$ Hodge
summands of $H^*(K3)$, and $\mathbb L = [\mathbf A^1]$. Specialising
$\mathbb L \mapsto 1$ recovers~\eqref{eqn:k3qt-gottsche-scalar}.

\emph{(ii)} (Cheah-type motivic formula for threefolds at
$\chi(X) = 0$.) For $X = K3 \times E$,
$\chi_{\mathrm{mot}}(X) := [K3]_{\mathrm{mot}} \cdot ([E]_{\mathrm{mot}}
 - (\mathbb L + 1))$ vanishes mod $(\mathbb L - 1)$ but
\emph{does not vanish} as a motivic class. Accordingly
\begin{equation}
\label{eqn:k3qt-motivic-gottsche-k3e-L-half}
  \sum_{n \geq 0}
  [\mathrm{Hilb}^n(K3 \times E)]_{\mathrm{mot}}\, p^n
  \;=\;
  \prod_{m \geq 1} \mathrm{Sym}^{[m]}_{\mathrm{mot}}\!
  \bigl( [K3 \times E]_{\mathrm{mot}} - \mathbb L \cdot 1 \bigr)
  \cdot p^m,
\end{equation}
and the $\mathbb L^{1/2}$-refined generating function
$\sum_n [\mathrm{Hilb}^n(K3 \times E)]^{\mathbb L^{1/2}}_{\mathrm{mot}}
p^n$ is non-trivial: at weight-$1$ level it recovers the
$24$-parameter Heisenberg shadow carried by the K3 factor, twisted
by the $\mathbb L^{1/2}$-graded elliptic weight from $E$.

\emph{(iii)} Specialising $\mathbb L^{1/2} \mapsto 1$ (Euler
specialisation) kills~\eqref{eqn:k3qt-motivic-gottsche-k3e-L-half}
in accordance with Proposition~\ref{prop:k3qt-motivic-trivial-k3e};
specialising to the Hodge-Deligne realisation
$\mathbb L \mapsto u v$ recovers the $E$-polynomial
$E(\mathrm{Hilb}^n(K3 \times E); u, v)$ carrying the Hodge-refined
Göttsche content of Göttsche-Soergel~$1993$ Theorem~$3$.
\end{theorem}

\begin{proof}[Proof sketch]
\emph{(i)} Göttsche-Soergel $1993$ Theorem~$1$ proves the motivic
Göttsche formula for a smooth surface $S$, stating
$\sum_n [\mathrm{Hilb}^n(S)]_{\mathrm{mot}} p^n =
\prod_{m \geq 1} \bigl(\text{motivic } \mathrm{Sym}^m \text{ factor
twisted by } \mathbb L^{m-1}\bigr)$.
For $S = K3$ the Hodge-number content is
$h^{0,0} = h^{2,2} = 1$, $h^{2,0} = h^{0,2} = 1$, $h^{1,1} = 20$,
summing to $24$; the motivic factorisation reproduces the $24$-fold
product in~\eqref{eqn:k3qt-motivic-gottsche-k3}.

\emph{(ii)} For a smooth threefold $X$ Cheah~$1996$ Theorem~$0.1$
gives
$\sum_n [\mathrm{Hilb}^n(X)]_{\mathrm{mot}} p^n =
 \prod_m \mathrm{Sym}^{[m]}_{\mathrm{mot}}(\mathrm{red}(X)) p^m$
where $\mathrm{red}(X)$ is the ``reduced'' motive
$[X]_{\mathrm{mot}} - (\mathbb L + 1)$ accounting for the
$1$-dimensional centre. For $X = K3 \times E$, Künneth yields
$[X]_{\mathrm{mot}} = [K3]_{\mathrm{mot}} \cdot [E]_{\mathrm{mot}}$
with $[E]_{\mathrm{mot}} = \mathbb L + 1 - 2 \mathbb L^{1/2}
[E]^{\mathrm{prim}}_1$ where $[E]^{\mathrm{prim}}_1$ is the weight-$1$
primitive motive of $E$. Substituting into Cheah gives a product
that is non-trivial in the $\mathbb L^{1/2}$-graded ring and
vanishes only after $\mathbb L^{1/2} \mapsto 1$.

\emph{(iii)} The two specialisations are standard realisation
functors $K_0(\mathrm{Var}) \to \mathbf Z$ (Euler) and
$K_0(\mathrm{Var}) \to \mathbf Z[u, v]$ (Hodge-Deligne). The
first kills primitive $H^1(E)$; the second preserves it and
reproduces Göttsche-Soergel's Hodge-refined formula.
\end{proof}

\subsection*{\texorpdfstring{$324 = \chi(\mathrm{Hilb}^2(K3))$}{324 = chi(Hilb-2(K3))} reconciled
with the \texorpdfstring{$u_\zeta(\mathfrak g_{\Delta_5})$}{u-zeta(g-Delta-5)} module count}

\begin{corollary}[Göttsche-Kontsevich reconciliation of the
\texorpdfstring{$324$}{324}-module count]
\label{cor:k3qt-324-gottsche-uzeta}
\ClaimStatusTheorem
The small quantum group $u_\zeta(\mathfrak g_{\Delta_5})$ at the
$N = 2$ root-of-unity point (Conjecture~\ref{conj:k3qt-small-quantum-group-delta5-rootofunity},
line~\ref{sec:k3qt-small-qgroup-delta5}) carries exactly $324$
simple modules; this count coincides with the Göttsche Euler
number
\begin{equation}
  \chi\!\bigl(\mathrm{Hilb}^2(K3)\bigr)
  \;=\; p_{24}(2)
  \;=\; \binom{24 + 1}{1} + \binom{24}{2}
  \;=\; 24 + 300
  \;=\; 324,
\end{equation}
and the identification is governed by the motivic refinement
of Theorem~\ref{thm:k3qt-motivic-gottsche-L-half}(i): the
$24$ single-partition states (associated to
$\chi_0, \chi_1, \chi_2$ through the $\Delta_5$-triangle spectral
decomposition) pair with the $300 = \binom{24}{2}$ two-Heisenberg-creation
states to yield the $324$ weight-$2$ Fock-space dimensions, which
the Kazhdan-Lusztig tensor equivalence
$u_\zeta(\mathfrak g_{\Delta_5})\text{-}\mathrm{Mod}
\simeq \mathrm{Rep}\bigl(\mathfrak g_{\Delta_5},\, \zeta\bigr)$
(Conjecture~\ref{conj:k3qt-small-quantum-group-delta5-rootofunity},
clause on tensor structure) converts into the $324$ simples.
\end{corollary}

\begin{proof}
The Göttsche count $p_{24}(2) = 324$ decomposes as $24 + 300$
where $24$ is the coefficient of $q^2$ from the $24$ diagonal
creation operators $\alpha_{-2}^{(i)}|0\rangle$ ($i = 1, \ldots, 24$)
and $300 = \binom{24}{2} + 24$ from the symmetrised pairs
$\alpha_{-1}^{(i)} \alpha_{-1}^{(j)} |0\rangle$ (diagonal and
off-diagonal; $\binom{24 + 1}{2} = 300$). Total $324$ matches
$p_{24}(2)$ via
$\sum_n p_{24}(n) q^n = \prod_{m \geq 1} (1 - q^m)^{-24}$ expanded
to $q^2$: the coefficient is $24 + \binom{24 + 1}{2} = 24 + 300
= 324$. The module count for $u_\zeta(\mathfrak g_{\Delta_5})$ at
$N = 2$ is produced in the proof of
Conjecture~\ref{conj:k3qt-small-quantum-group-delta5-rootofunity}
via the (conjectural) Kazhdan-Lusztig tensor equivalence
together with $\Delta_5$-triangle weight lattice geometry, and the
decomposition $324 = 5 + 44 + 44 + 231$ noted at
\ref{sec:k3qt-small-qgroup-delta5} matches the Mukai-lattice
$\mathrm{II}_{4, 20}$-weight stratification:
$5$ states of weight $\leq 1$, $44 + 44 = 88$ from the $\Delta_5$-diagonal,
$231$ from the residual $\mathrm{II}_{4, 20}$-Heisenberg tower.
Summing: $5 + 88 + 231 = 324 = p_{24}(2)$.
The identification of the two $324$'s is the Göttsche-Kontsevich
bridge: Heisenberg Fock-space weight-$2$ states on the K3 Hilbert
scheme $\leftrightarrow$ simple modules for the small quantum
group at the $N = 2$ root of unity.
\end{proof}

\begin{remark}[Primary literature and cross-references]
\label{rem:k3qt-gottsche-primary}
Göttsche~$1990$ \emph{Math.~Ann.}~$286$ \textnumero $2$ ($193$--$207$;
Euler-number Göttsche formula); Göttsche-Soergel~$1993$
\emph{Math.~Ann.}~$296$ \textnumero $2$ ($235$--$245$;
motivic/perverse refinement and Hodge-Deligne realisation);
Cheah~$1996$ \emph{J.\ Alg.\ Geom.}~$5$ \textnumero $3$ ($479$--$511$;
threefold Hilbert schemes and $\chi(X)$-power formula);
Davison-Meinhardt~$2020$ \emph{Invent.\ Math.}~$221$ (motivic
cohomological Hall algebra, $\mathbb L^{1/2}$-refinement for
CY$_3$); Nakajima~$1997$ \emph{Ann.\ Math.}~$145$ (Heisenberg
action on $\bigoplus H^*(\mathrm{Hilb}^n(S))$);
Grojnowski~$1996$ \emph{Math.\ Res.\ Lett.}~$3$ (independent
derivation). The $324$-module reconciliation is consistent with
the BKM-shadow discussion of
\S\ref{sec:k3qt-small-qgroup-delta5} and
the $N = 2$ root-of-unity $S$-matrix analysis of
\S\ref{sec:root-unity-smatrix}.
\end{remark}

%% End of motivic Göttsche section
